% L'article de journal — Allemand 1ère A — Cours signé OnBuch+
% Programme officiel MINESEC. Compiler avec : tectonic ALA18-l-article-de-journal.tex
\newcommand{\DOCMATIERE}{Allemand}
\newcommand{\DOCNIVEAU}{1ère A}
\newcommand{\DOCMODULE}{Chapitre 4 : Média et communication}
\newcommand{\DOCLECON}{ALA18}
\newcommand{\DOCTITRE}{L'article de journal}
% OnBuch+ — préambule « cours signés » ENRICHI (compilé avec tectonic / XeLaTeX)
% Système visuel « Pop » d'OnBuch+ : fond crème (jamais blanc), bordures encre
% épaisses, ombres « dures » décalées, titres Archivo Black en capitales.
% Version MATHÉMATIQUES : graphes (pgfplots/tikz), boîtes pédagogiques
% (définition, propriété, méthode, attention, le savais-tu, exercice résolu,
% exercice, TP, bilan), page de garde et signature OnBuch+ ancrée en bas.
%
% Macros attendues dans main.tex AVANT \input{preamble} :
%   \DOCMATIERE  \DOCNIVEAU  \DOCMODULE  \DOCLECON (numéro)  \DOCTITRE
\documentclass[11pt]{article}

\usepackage{fontspec}
\usepackage[ngerman,french]{babel} % français = langue principale ; l'allemand via \all{..} / textebox
\usepackage[a4paper,margin=2cm,top=3.4cm,bottom=2.8cm,headheight=1.4cm,footskip=1.3cm]{geometry}
\usepackage{amsmath,amssymb,amsfonts}
\usepackage{mathtools} % \xmapsto, \coloneqq... souvent utilisés par les modèles
\usepackage{cancel} % simplifications barrées (bilans de forces)
% \abs{z} (module, valeur absolue) et \norm{u} (norme), utilisés par les modèles
\providecommand{\abs}{}\let\abs\relax\DeclarePairedDelimiter{\abs}{\lvert}{\rvert}
\providecommand{\norm}{}\let\norm\relax\DeclarePairedDelimiter{\norm}{\lVert}{\rVert}
\usepackage[table]{xcolor} % [table] : fournit aussi \rowcolors (alternance de lignes)
\usepackage{pagecolor}
\usepackage{tikz}
\usetikzlibrary{arrows.meta,positioning,calc,shapes.geometric,shapes.misc,shapes.arrows,decorations.pathmorphing,decorations.pathreplacing,decorations.markings,patterns,shadows,mindmap,trees,fit,backgrounds,matrix,intersections,angles,quotes}
\usepackage{pgfplots}
\usepackage[version=4]{mhchem} % équations nucléaires \ce{^{235}_{92}U}
\usepackage[european,siunitx]{circuitikz} % schémas électriques
\pgfplotsset{compat=1.18}
\usepgfplotslibrary{fillbetween,groupplots} % aires entre courbes (calcul intégral)
\usepackage{enumitem}
\usepackage{fancyhdr}
% [most] charge toutes les bibliothèques tcolorbox (listings, minted, raster,
% documentation...) et gonfle inutilement la mémoire TeX sur un document long
% (~300 boîtes) : on ne charge que ce qui est réellement utilisé.
\usepackage[skins,breakable]{tcolorbox}
\usepackage{graphicx}
\usepackage{colortbl}
\usepackage{booktabs}
\usepackage{tabularx}
\usepackage{diagbox} % case d'en-tête diagonale des tableaux à double entrée
% Colonnes extensibles centrée / gauche / droite (C, L, R), souvent utilisées par les modèles
\newcolumntype{C}{>{\centering\arraybackslash}X}
\newcolumntype{L}{>{\raggedright\arraybackslash}X}
\newcolumntype{R}{>{\raggedleft\arraybackslash}X}
\usepackage{array}
\usepackage{multicol}
\usepackage{multirow}
\usepackage{siunitx}
% parse-numbers=false : accepte aussi \SI{2,5 \times 10^{-3}}{...}
\sisetup{locale=FR,output-decimal-marker={,},group-minimum-digits=4,parse-numbers=false}
\DeclareSIUnit\ohm{\ensuremath{\symup{\Omega}}} % U+2126 absent de la police
\DeclareSIUnit\division{div} % réglages d'oscilloscope (ms/div, V/div)
\DeclareSIUnit\tr{tr} % tours (vitesse de rotation, ex. \SI{1500}{\tr\per\minute})
\DeclareSIUnit\molar{M} % concentration molaire (ex. \SI{3}{\molar}, \SI{1}{\micro\molar})
\DeclareSIUnit\an{an} % année (le modèle confond parfois avec \year, un compteur TeX) — ex. \SI{5}{\kilo\meter\per\an}
\DeclareSIUnit\FCFA{FCFA} % franc CFA (ex. \SI{500}{\FCFA\per\kilo\gram})
\DeclareSIUnit\degreeBrix{\textdegree Brix} % teneur en sucre d'un sirop/jus (réfractométrie)
\DeclareSIUnit\month{mois} % le modèle utilise parfois \month, un compteur TeX, comme unité de durée
\usepackage{qrcode}
% \degree : commande fréquemment utilisée par les modèles pour un angle ou une
% température en dehors de \SI{}{} (non fournie par les paquets chargés).
\providecommand{\degree}{^{\circ}}
% \qed (fin de démonstration), utilisé par les modèles sans amsthm
\providecommand{\qed}{\hfill$\square$}
% \barre : double barre des valeurs interdites dans un tableau de variations (tkz-tab)
\providecommand{\barre}{\ensuremath{\|}}
% environnement proof (amsthm non chargé) utilisé par les modèles
\ifdefined\proof\else\newenvironment{proof}[1][Démonstration]{\par\noindent\textit{#1.}\ }{\qed\par}\fi
\usepackage{lastpage}
\usepackage{hyperref}
\hypersetup{hidelinks}

% ============================================================
% Palette « Pop » OnBuch+ — mêmes hex que la classe P (plus_ui.dart)
% ============================================================
\definecolor{poporange}{HTML}{F59321}
\definecolor{popdark}{HTML}{E07A0C}
\definecolor{popcream}{HTML}{FFF6E8}
\definecolor{popink}{HTML}{1C1714}
\definecolor{poppurple}{HTML}{7A5AE0}
\definecolor{popgreen}{HTML}{2E9E6A}
\definecolor{poppink}{HTML}{E0457B}
\definecolor{popblue}{HTML}{2F7DE1}
\definecolor{popgold}{HTML}{C98A12}
\definecolor{popmuted}{HTML}{8A7F76}
\definecolor{popline}{HTML}{EADFD2}
\definecolor{popgray}{HTML}{8A7F76}
\colorlet{popgrey}{popgray}
% Alias tolérants (le modèle nomme parfois les couleurs différemment)
\colorlet{popwhite}{white}
\colorlet{popblack}{popink}
\colorlet{popred}{poppink}
\colorlet{popyellow}{popgold}
% Teintes claires (fonds de tableaux/encadrés) — le modèle nomme parfois
% les couleurs avec un suffixe L (ex. popblueL) sans les avoir définies.
\colorlet{poporangeL}{poporange!20}
\colorlet{popdarkL}{popdark!20}
\colorlet{poppurpleL}{poppurple!20}
\colorlet{popgreenL}{popgreen!20}
\colorlet{poppinkL}{poppink!20}
\colorlet{popblueL}{popblue!20}
\colorlet{popgoldL}{popgold!20}
\colorlet{popredL}{popred!20}
\colorlet{popgrayL}{popgray!20}
% Teintes claires pour fonds de tableaux / zones
\colorlet{popblueL}{popblue!10!white}
\colorlet{poporangeL}{poporange!14!white}
\colorlet{popgreenL}{popgreen!12!white}
\colorlet{poppurpleL}{poppurple!12!white}
\colorlet{poppinkL}{poppink!12!white}
\colorlet{popgoldL}{popgold!14!white}

\pagecolor{popcream}

% ============================================================
% Typographie
% ============================================================
\defaultfontfeatures{Ligatures=TeX}
\setmainfont{PlusJakartaSans-Medium.ttf}[Path=../../../fonts/,BoldFont=PlusJakartaSans-Bold.ttf]
% Maths en sans-serif (Fira Math) avec chiffres et lettres droites de Plus
% Jakarta, pour que les formules se fondent dans le texte.
\usepackage[math-style=ISO,bold-style=ISO]{unicode-math}
\setmathfont{FiraMath-Regular.otf}[Path=../../../fonts/]
\setmathfont{PlusJakartaSans-Medium.ttf}[Path=../../../fonts/,range={up/num,up/{Latin,latin},\mathup}]
\usepackage{icomma} % virgule décimale sans espace en mode math
% Caractères mathématiques Unicode absents des polices de texte (Plus Jakarta,
% Archivo Black) : rendus par leur équivalent mathématique, en texte comme en
% formule — sinon ils s'affichent en carré vide (ex. « Arithmétique dans ℤ »).
\usepackage{newunicodechar}
\newunicodechar{ℝ}{\ensuremath{\mathbb{R}}}
\newunicodechar{ℤ}{\ensuremath{\mathbb{Z}}}
\newunicodechar{ℂ}{\ensuremath{\mathbb{C}}}
\newunicodechar{ℕ}{\ensuremath{\mathbb{N}}}
\newunicodechar{ℚ}{\ensuremath{\mathbb{Q}}}
\newunicodechar{𝔻}{\ensuremath{\mathbb{D}}}
\newunicodechar{α}{\ensuremath{\alpha}}
\newunicodechar{β}{\ensuremath{\beta}}
\newunicodechar{γ}{\ensuremath{\gamma}}
\newunicodechar{δ}{\ensuremath{\delta}}
\newunicodechar{ε}{\ensuremath{\varepsilon}}
\newunicodechar{θ}{\ensuremath{\theta}}
\newunicodechar{λ}{\ensuremath{\lambda}}
\newunicodechar{μ}{\ensuremath{\mu}}
\newunicodechar{σ}{\ensuremath{\sigma}}
\newunicodechar{τ}{\ensuremath{\tau}}
\newunicodechar{φ}{\ensuremath{\varphi}}
\newunicodechar{ψ}{\ensuremath{\psi}}
\newunicodechar{ω}{\ensuremath{\omega}}
\newunicodechar{Σ}{\ensuremath{\Sigma}}
% majuscules produites par \MakeUppercase dans les titres (α-aminés -> Α)
\newunicodechar{Α}{\ensuremath{\alpha}}
\newunicodechar{Β}{\ensuremath{\beta}}
\newunicodechar{Γ}{\ensuremath{\Gamma}}
\newunicodechar{Δ}{\ensuremath{\Delta}}
\newunicodechar{Ε}{\ensuremath{\varepsilon}}
\newunicodechar{Θ}{\ensuremath{\Theta}}
\newunicodechar{Λ}{\ensuremath{\Lambda}}
\newunicodechar{Μ}{\ensuremath{\mu}}
\newunicodechar{Τ}{\ensuremath{\tau}}
\newunicodechar{Φ}{\ensuremath{\Phi}}
\newunicodechar{Ψ}{\ensuremath{\Psi}}
\newunicodechar{Ω}{\ensuremath{\Omega}}
\newunicodechar{↦}{\ensuremath{\mapsto}}
\newunicodechar{∈}{\ensuremath{\in}}
\newunicodechar{∉}{\ensuremath{\notin}}
\newunicodechar{∧}{\ensuremath{\wedge}}
\newunicodechar{⇔}{\ensuremath{\Leftrightarrow}}
\newunicodechar{⟺}{\ensuremath{\Longleftrightarrow}}
\newunicodechar{⇒}{\ensuremath{\Rightarrow}}
\newunicodechar{⟹}{\ensuremath{\Longrightarrow}}
\newunicodechar{∘}{\ensuremath{\circ}}
\newunicodechar{∪}{\ensuremath{\cup}}
\newunicodechar{∩}{\ensuremath{\cap}}
\newunicodechar{≡}{\ensuremath{\equiv}}
\newunicodechar{⊂}{\ensuremath{\subset}}
\newunicodechar{⊥}{\ensuremath{\perp}}
\newunicodechar{∃}{\ensuremath{\exists}}
\newunicodechar{∀}{\ensuremath{\forall}}
\newunicodechar{∛}{\ensuremath{\sqrt[3]{\,}}}
\newunicodechar{∜}{\ensuremath{\sqrt[4]{\,}}}
\newunicodechar{⃗}{\ensuremath{{}^{\to}}}
\newunicodechar{ⁿ}{\ifmmode^{n}\else\textsuperscript{\ensuremath{n}}\fi}
\newunicodechar{ˣ}{\ifmmode^{x}\else\textsuperscript{\ensuremath{x}}\fi}
\newunicodechar{ᵇ}{\ifmmode^{b}\else\textsuperscript{\ensuremath{b}}\fi}
\newunicodechar{ʳ}{\ifmmode^{r}\else\textsuperscript{\ensuremath{r}}\fi}
\newunicodechar{ᵘ}{\ifmmode^{u}\else\textsuperscript{\ensuremath{u}}\fi}
\newunicodechar{ʸ}{\ifmmode^{y}\else\textsuperscript{\ensuremath{y}}\fi}
\newunicodechar{ᵃ}{\ifmmode^{a}\else\textsuperscript{\ensuremath{a}}\fi}
\newunicodechar{ᵉ}{\ifmmode^{e}\else\textsuperscript{\ensuremath{e}}\fi}
\newunicodechar{ᵏ}{\ifmmode^{k}\else\textsuperscript{\ensuremath{k}}\fi}
\newunicodechar{ᵖ}{\ifmmode^{p}\else\textsuperscript{\ensuremath{p}}\fi}
\newunicodechar{ᴺ}{\ifmmode^{N}\else\textsuperscript{\ensuremath{N}}\fi}
\newunicodechar{ᶜ}{\ifmmode^{c}\else\textsuperscript{\ensuremath{c}}\fi}
\newunicodechar{ʰ}{\ifmmode^{h}\else\textsuperscript{\ensuremath{h}}\fi}
\newunicodechar{ᵠ}{\ifmmode^{q}\else\textsuperscript{\ensuremath{q}}\fi}
\newunicodechar{ₙ}{\ifmmode_{n}\else\textsubscript{\ensuremath{n}}\fi}
\newunicodechar{ₐ}{\ifmmode_{a}\else\textsubscript{\ensuremath{a}}\fi}
\newunicodechar{ᵢ}{\ifmmode_{i}\else\textsubscript{\ensuremath{i}}\fi}
\newunicodechar{ᵣ}{\ifmmode_{r}\else\textsubscript{\ensuremath{r}}\fi}
\newunicodechar{ₖ}{\ifmmode_{k}\else\textsubscript{\ensuremath{k}}\fi}
\newunicodechar{ᵦ}{\ifmmode_{\beta}\else\textsubscript{\ensuremath{\beta}}\fi}
\newunicodechar{ₓ}{\ifmmode_{x}\else\textsubscript{\ensuremath{x}}\fi}
\newfontfamily\popdisplay{ArchivoBlack-Regular.ttf}[Path=../../../fonts/]
\newfontfamily\poplogo{SpaceGrotesk-Bold.ttf}[Path=../../../fonts/]
\newfontfamily\popsemi{PlusJakartaSans-SemiBold.ttf}[Path=../../../fonts/]
\newfontfamily\popxbold{PlusJakartaSans-ExtraBold.ttf}[Path=../../../fonts/]
\newfontfamily\popxboldls{PlusJakartaSans-ExtraBold.ttf}[Path=../../../fonts/,LetterSpace=3.0]

\setlist[enumerate]{leftmargin=1.7em,itemsep=5pt,topsep=4pt}
\setlist[itemize]{leftmargin=1.7em,itemsep=4pt,topsep=4pt,label={\color{poporange}\textbullet}}
\setlength{\parindent}{0pt}
\setlength{\parskip}{5pt}
\renewcommand{\arraystretch}{1.25}

% pgfplots : style Pop par défaut
\pgfplotsset{
  every axis/.append style={axis line style={popink,line width=1pt},tick style={popink},
    grid style={popline,line width=0.5pt},label style={font=\small},tick label style={font=\footnotesize}},
  popcurve/.style={poporange,line width=1.8pt,no markers,smooth},
}
% Même style disponible dans un \draw TikZ simple (hors pgfplots)
\tikzset{popcurve/.style={poporange,line width=1.8pt}}
\tikzset{popgrid/.style={popline,very thin}}
\colorlet{popcurve}{poporange} % le nom du style est parfois utilisé comme couleur

% ============================================================
% Pill / squiggle
% ============================================================
\newcommand{\poppill}[3]{%
  \tikz[baseline=-0.55ex]\node[draw=popink,line width=1.2pt,rounded corners=2.6pt,%
    fill=#2,inner xsep=6.5pt,inner ysep=3pt]{%
    \popxboldls\fontsize{7}{7}\selectfont\color{#3}\MakeUppercase{#1}};%
}
\newcommand{\popsquiggle}[1]{%
  \tikz[baseline=-0.4ex]{
    \draw[#1,line width=2.6pt,line cap=round]
      plot [smooth] coordinates {(0,0.09) (0.34,0.01) (0.68,0.21) (1.02,0.01) (1.36,0.10)};
  }%
}
% Mot-clé surligné (marqueur orange sous le texte)
\newcommand{\cle}[1]{{\popxbold\color{popdark}#1}}

% ============================================================
% En-tête / pied
% ============================================================
\pagestyle{fancy}
\fancyhf{}
\renewcommand{\headrulewidth}{0pt}
\renewcommand{\footrulewidth}{0pt}
\lhead{\raisebox{-0.15ex}{\poplogo\fontsize{13}{13}\selectfont\color{popink}OnBuch\color{poporange}+}}
\rhead{\poppill{\DOCMATIERE\ \textbullet\ \DOCNIVEAU\ \textbullet\ Leçon \DOCLECON}{white}{popink}}
\lfoot{\popxbold\fontsize{7.6}{7.6}\selectfont\color{popmuted}\MakeUppercase{Cours signé OnBuch+}\ \textbullet\ {\popsemi\DOCTITRE}}
\rfoot{\tikz[baseline=-0.6ex]\node[rounded corners=8.5pt,fill=popink,minimum height=17pt,minimum width=17pt,inner xsep=5pt,inner sep=0]{\popxbold\fontsize{8}{8}\selectfont\color{popcream}\thepage\,/\,\pageref*{LastPage}};}

% ============================================================
% Titres
% ============================================================
\newcommand{\chaptitre}[2]{%
  \begin{tcolorbox}[enhanced,colback=poporange,colframe=popink,boxrule=2.6pt,arc=11pt,
      drop shadow=popink,left=15pt,right=15pt,top=12pt,bottom=11pt,before skip=3pt,after skip=18pt]
    \poppill{#2}{popink}{popcream}\par\vspace{9pt}
    {\raggedright\hyphenpenalty=10000\exhyphenpenalty=10000\popdisplay\fontsize{22}{24}\selectfont\color{popcream}\MakeUppercase{#1}\par}
  \end{tcolorbox}
}
% Section numérotée : pastille orange avec le numéro + titre Archivo
\newcounter{popsec}
\newcommand{\coursec}[1]{%
  \stepcounter{popsec}\setcounter{popsub}{0}%
  \par\vspace{16pt}\noindent
  \begin{minipage}{\linewidth}
  \tikz[baseline=-0.7ex]\node[draw=popink,line width=1.6pt,fill=poporange,rounded corners=4pt,minimum size=22pt,inner sep=0]{\popdisplay\fontsize{12}{12}\selectfont\color{popcream}\thepopsec};%
  \hspace{8pt}{\popdisplay\fontsize{15}{17}\selectfont\color{popink}\MakeUppercase{#1}}\par
  \vspace{4pt}\noindent{\color{poporange}\rule{36pt}{3.6pt}}
  \end{minipage}\par\nopagebreak\vspace{8pt}
}
\newcounter{popsub}
\newcommand{\courssub}[1]{%
  \stepcounter{popsub}%
  \par\vspace{9pt}\noindent{\popxbold\fontsize{11.5}{13}\selectfont\color{popink}{\color{poporange}\thepopsec.\thepopsub}\hspace{6pt}#1}\par\nopagebreak\vspace{4pt}
}

% ============================================================
% Boîtes pédagogiques — même recette : fond blanc, bordure épaisse colorée,
% ombre dure encre, badge plein. Titre optionnel [#1].
% ============================================================
\tcbset{popbase/.style={breakable,enhanced,colback=white,boxrule=2.3pt,arc=9pt,
  shadow={3pt}{-3pt}{0pt}{popink},left=14pt,right=14pt,top=11pt,bottom=11pt,before skip=11pt,after skip=15pt,
  fontupper=\popsemi\fontsize{10.6}{14.5}\selectfont\color{popink},
  fontlower=\popsemi\fontsize{10.6}{14.5}\selectfont\color{popink}}}
% Coche dessinée (le glyphe ✓ n'existe pas dans les polices du document)
\newcommand{\popcheck}[1]{\tikz[baseline=-0.5ex]\draw[#1,line width=1.6pt,line cap=round,line join=round](-0.13,0.0)--(-0.04,-0.09)--(0.13,0.1);}
\providecommand{\coche}{\popcheck{popgreen}}
\providecommand{\resultat}[1]{\par\smallskip\noindent\fcolorbox{popgreen}{popgreenL}{\parbox{\dimexpr\linewidth-2\fboxsep-2\fboxrule\relax}{\textbf{Résultat.} #1}}\par\smallskip}
\newcommand{\popboxhead}[3]{\poppill{#1}{#2}{white}\ifx\relax#3\relax\else\hspace{6pt}{\popxbold\fontsize{10.6}{12}\selectfont\color{popink}#3}\fi\par\vspace{7pt}}

\newtcolorbox{aretenir}[1][]{popbase,colframe=popgreen,before upper={\popboxhead{À retenir}{popgreen}{#1}}}
% Texte en allemand (document de lecture, dialogue, extrait) : cadre
% d'étude. Un titre optionnel (source, auteur) s'affiche dans l'en-tête.
\newtcolorbox{textebox}[1][]{popbase,colframe=popblue,before upper={\popboxhead{Text}{popblue}{#1}\begin{otherlanguage}{ngerman}},after upper={\end{otherlanguage}}}
% Marques ✓ / ✗ (forme correcte / fautive) souvent utilisées par les modèles
\providecommand{\cross}{\textcolor{poppink}{\ensuremath{\times}}}
\providecommand{\xmark}{\cross}\providecommand{\crossmark}{\cross}\providecommand{\cmark}{\ensuremath{\checkmark}}
% Ligne de réponse / justification à compléter (inventée par les modèles)
\providecommand{\justification}[1]{\par\noindent\textit{Justification~:} \dotfill\par}
% \textipa (package tipa non chargé) : la police ne couvre pas l'API -> simple passage
\providecommand{\textipa}[1]{#1}
% Soulignement (ulem non chargé) : \uline, \ul
\providecommand{\uline}[1]{\underline{#1}}\providecommand{\ul}[1]{\underline{#1}}
% \alert (beamer) inventée par les modèles : texte mis en évidence
\providecommand{\alert}[1]{\textbf{\textcolor{poppink}{#1}}}
% variantes ASCII d'ordinaux écrites en macro
\providecommand{\iere}{\textsuperscript{re}}\providecommand{\ieme}{\textsuperscript{e}}\providecommand{\ere}{\textsuperscript{re}}\providecommand{\eme}{\textsuperscript{e}}\providecommand{\er}{\textsuperscript{er}}\providecommand{\ie}{\textsuperscript{e}}
% Ordinaux français écrits en macro par les modèles : 1\ière, 3\ième
\newcommand{\ière}{\textsuperscript{re}}\newcommand{\ième}{\textsuperscript{e}}
% \exemple (inventée par les modèles) : étiquette « Exemple : »
\providecommand{\exemple}{\textbf{Exemple~:}\ }
% \square utilisable aussi hors mode math (cases à cocher)
\AtBeginDocument{\let\squareMath\square\renewcommand{\square}{\ensuremath{\squareMath}}}
% Mot ou phrase en allemand à l'intérieur d'un paragraphe français
\newcommand{\all}[1]{\ifmmode\text{\textit{#1}}\else\begingroup\selectlanguage{ngerman}\itshape #1\endgroup\fi}
\let\esp\all % alias toléré
\newtcolorbox{exemplebox}[1][]{popbase,colframe=poporange,before upper={\popboxhead{Exemple}{poporange}{#1}}}
\newtcolorbox{definition}[1][]{popbase,colframe=popblue,before upper={\popboxhead{Définition}{popblue}{#1}}}
\newtcolorbox{propriete}[1][]{popbase,colframe=poppurple,before upper={\popboxhead{Propriété}{poppurple}{#1}}}
\newtcolorbox{methode}[1][]{popbase,colframe=popgold,before upper={\popboxhead{Méthode}{popgold}{#1}}}
\newtcolorbox{attention}[1][]{popbase,colframe=poppink,before upper={\popboxhead{Attention}{poppink}{#1}}}
\newtcolorbox{experience}[1][]{popbase,colframe=popblue,colback=popblueL,before upper={\popboxhead{Expérience}{popblue}{#1}}}
% « Le savais-tu ? » : carte violette pleine (contexte, culture, Cameroun)
\newtcolorbox{savaistu}[1][]{popbase,colback=poppurple,colframe=popink,
  fontupper=\popsemi\fontsize{10.4}{14.2}\selectfont\color{white},
  before upper={\poppill{Le savais-tu\,?}{white}{poppurple}\ifx\relax#1\relax\else\hspace{6pt}{\popxbold\color{white}#1}\fi\par\vspace{7pt}}}
% Exercice résolu : énoncé en haut, \tcblower puis la solution sur fond crème
\newcounter{popexo}\newcounter{popexr}
\newtcolorbox{exoresolu}[1][]{popbase,colframe=poporange,colbacklower=popcream,
  segmentation style={popink,line width=1.2pt,dashed},
  before upper={\stepcounter{popexr}\popboxhead{Exercice résolu \thepopexr}{poporange}{#1}},
  before lower={\poppill{Solution}{popink}{popcream}\par\vspace{6pt}}}
% Exercice (non résolu en place) — la correction est dans la partie Corrigés
\newtcolorbox{exercice}[1][]{popbase,colframe=popink,
  before upper={\stepcounter{popexo}\popboxhead{Exercice \thepopexo}{popink}{#1}}}
% Corrigé
\newtcolorbox{corrige}[1][]{popbase,colframe=popgreen,colback=popgreenL,
  before upper={\popboxhead{Corrigé}{popgreen}{#1}}}
% Objectifs / prérequis / bilan
\newtcolorbox{objectifs}{popbase,colframe=popink,colback=poporangeL,
  before upper={\popboxhead{Objectifs de la leçon}{popink}{}}}
\newtcolorbox{prerequis}{popbase,colframe=popink,colback=popblueL,
  before upper={\popboxhead{Prérequis}{popblue}{}}}
\newtcolorbox{bilan}{popbase,colframe=popink,colback=white,boxrule=2.8pt,
  before upper={\popboxhead{Fiche bilan}{poporange}{L'essentiel à connaître pour l'examen}}}
% Figure encadrée avec légende
\newtcolorbox{popfigure}[1][]{enhanced,breakable=false,colback=white,colframe=popline,boxrule=1.4pt,arc=8pt,
  left=10pt,right=10pt,top=10pt,bottom=8pt,before skip=10pt,after skip=12pt,halign=center,
  lower separated=false,halign lower=center,
  fontlower=\popsemi\fontsize{9}{11.5}\selectfont\color{popmuted},
  #1}
\newcommand{\legende}[1]{\par\vspace{6pt}{\popsemi\fontsize{9}{11.5}\selectfont\color{popmuted}#1}}

% ============================================================
% Signature OnBuch+
% ============================================================
\newcommand{\poplogoinline}[1]{{\poplogo\fontsize{#1}{#1}\selectfont\color{popink}OnBuch\color{poporange}+}}
\newcommand{\signatureonbuch}{%
  \begin{tcolorbox}[enhanced,colback=popink,colframe=popink,boxrule=2.2pt,arc=10pt,
      drop shadow=poporange,left=14pt,right=14pt,top=10pt,bottom=10pt,before skip=0pt,after skip=0pt]
    \tikz[baseline=-0.7ex]\node[circle,fill=popgreen,draw=popcream,line width=1.3pt,minimum size=22pt,inner sep=0]{\popcheck{white}};%
    \hspace{9pt}\begin{minipage}[c]{0.62\linewidth}
      {\popxbold\fontsize{12}{13}\selectfont\color{popcream}Signé\ \textbullet\ {\poplogo OnBuch\color{poporange}+}}\par\vspace{2pt}
      {\raggedright\popsemi\fontsize{8}{10.5}\selectfont\color{popline}\MakeUppercase{Équipe pédagogique OnBuch+}\\\MakeUppercase{Conforme au programme officiel MINESEC}\par}
    \end{minipage}\hfill
    {\poplogo\fontsize{22}{22}\selectfont\color{popcream}OnBuch\color{poporange}+}
  \end{tcolorbox}%
}

% ============================================================
% Page de garde
% #1 titre · #2 matière · #3 niveau · #4 module · #5 n° leçon · #6 durée ·
% #7 description / objectifs (texte ou itemize)
% ============================================================
\newcommand{\pagegarde}[7]{%
  \thispagestyle{empty}
  \newgeometry{a4paper,margin=1.7cm,top=1.6cm,bottom=1.4cm}%
  \begin{tikzpicture}[remember picture,overlay]
    \fill[poporange!18] ([xshift=-3.2cm,yshift=-3.6cm]current page.north east) circle (4.6cm);
    \fill[poppurple!14] ([xshift=2.4cm,yshift=7.6cm]current page.south west) circle (3.8cm);
  \end{tikzpicture}%
  {\poplogo\fontsize{30}{30}\selectfont\color{popink}OnBuch\color{poporange}+}\hfill
  \raisebox{0.6ex}{\poppill{Cours signé \textbullet\ Premium}{popink}{popcream}}\par
  \vspace{1pt}\popsquiggle{poppurple}\par
  \vspace{8pt}
  \begin{tcolorbox}[enhanced,colback=poporange,colframe=popink,boxrule=2.8pt,arc=13pt,
      drop shadow=popink,left=18pt,right=18pt,top=12pt,bottom=13pt,after skip=10pt]
    \poppill{#2 \textbullet\ #3}{popink}{popcream}\hspace{5pt}\poppill{#4}{white}{popink}\par\vspace{8pt}
    {\popdisplay\fontsize{14}{14}\selectfont\color{popink}LEÇON #5}\par\vspace{4pt}
    {\raggedright\hyphenpenalty=10000\exhyphenpenalty=10000\popdisplay\fontsize{25}{27}\selectfont\color{popcream}\MakeUppercase{#1}\par}
    \vspace{8pt}
    {\popxbold\fontsize{9.5}{9.5}\selectfont\color{popink}\MakeUppercase{Durée conseillée : #6}}
  \end{tcolorbox}
  \begin{tcolorbox}[enhanced,colback=white,colframe=popink,boxrule=2.2pt,arc=10pt,
      drop shadow=popink,left=16pt,right=16pt,top=10pt,bottom=10pt,after skip=8pt]
    \poppill{Dans cette leçon}{poppurple}{white}\par\vspace{5pt}
    \popsemi\fontsize{9.3}{12.6}\selectfont\color{popink}#7
  \end{tcolorbox}
  \noindent\begin{minipage}[t]{0.487\linewidth}
    \begin{tcolorbox}[enhanced,colback=popgreen,colframe=popink,boxrule=2.2pt,arc=10pt,
        drop shadow=popink,left=11pt,right=11pt,top=10pt,bottom=10pt,height=3.1cm,valign=center]
      \tcbox[colback=white,colframe=popink,boxrule=1.3pt,arc=5pt,boxsep=0pt,nobeforeafter,
        left=3pt,right=3pt,top=3pt,bottom=3pt,valign=center]{\qrcode[height=1.9cm]{https://play.google.com/store/apps/details?id=com.luvvix.onbuch&hl=fr}}%
      \hspace{8pt}\begin{minipage}[c]{0.5\linewidth}
        {\popdisplay\fontsize{11}{12.5}\selectfont\color{white}\MakeUppercase{Télécharge OnBuch}}\par\vspace{4pt}
        {\raggedright\popsemi\fontsize{8.4}{11}\selectfont\color{white}Gratuit \textbullet\ Google Play\\Tuteur IA, annales, concours, résultats\par}
      \end{minipage}
    \end{tcolorbox}
  \end{minipage}\hfill
  \begin{minipage}[t]{0.487\linewidth}
    \begin{tcolorbox}[enhanced,colback=poppurple,colframe=popink,boxrule=2.2pt,arc=10pt,
        drop shadow=popink,left=11pt,right=11pt,top=10pt,bottom=10pt,height=3.1cm,valign=center]
      \tikz[baseline=-0.7ex]\node[circle,fill=white,draw=popink,line width=1.3pt,minimum size=1.6cm,inner sep=0]{\popdisplay\fontsize{24}{24}\selectfont\color{poppurple}+};%
      \hspace{8pt}\begin{minipage}[c]{0.6\linewidth}
        {\popdisplay\fontsize{11}{12.5}\selectfont\color{white}\MakeUppercase{Passe à OnBuch+}}\par\vspace{4pt}
        {\raggedright\popsemi\fontsize{8.4}{11}\selectfont\color{white}Tous les cours signés, Tuteur IA illimité, fiches PDF — onglet OnBuch+ de l'app\par}
      \end{minipage}
    \end{tcolorbox}
  \end{minipage}
  \vspace{8pt}
  \popcontenu
  \vfill
  \signatureonbuch
  \restoregeometry
  \newpage
}

% Rangée « Ce cours contient » (page de garde)
\newcommand{\poptile}[3]{% #1 couleur #2 gros texte #3 légende
  \begin{tcolorbox}[enhanced,colback=#1,colframe=popink,boxrule=2pt,arc=9pt,shadow={2.5pt}{-2.5pt}{0pt}{popink},
      left=6pt,right=6pt,top=9pt,bottom=9pt,halign=center,valign=center,height=1.75cm,width=\linewidth,nobeforeafter]
    {\popdisplay\fontsize{12.5}{13}\selectfont\color{white}\MakeUppercase{#2}}\par\vspace{3pt}
    {\centering\popsemi\fontsize{7.8}{9.5}\selectfont\color{white}#3\par}
  \end{tcolorbox}}
\newcommand{\popcontenu}{%
  \par\noindent{\popxboldls\fontsize{7.5}{7.5}\selectfont\color{popmuted}CE COURS SIGNÉ CONTIENT}\par\vspace{7pt}
  \noindent\begin{minipage}[t]{0.235\linewidth}\poptile{popblue}{Cours}{complet, illustré, pas à pas}\end{minipage}\hfill
  \begin{minipage}[t]{0.235\linewidth}\poptile{poporange}{Méthodes}{et exercices résolus}\end{minipage}\hfill
  \begin{minipage}[t]{0.235\linewidth}\poptile{poppink}{Exercices}{type examen + corrigés}\end{minipage}\hfill
  \begin{minipage}[t]{0.235\linewidth}\poptile{popgreen}{Fiche bilan}{l'essentiel à retenir}\end{minipage}\par}

% Sommaire
\newtcolorbox{sommaire}{enhanced,colback=white,colframe=popink,boxrule=2.2pt,arc=10pt,
  drop shadow=popink,left=18pt,right=18pt,top=13pt,bottom=13pt,before skip=6pt,after skip=20pt,
  before upper={\poppill{Sommaire}{poppurple}{white}\par\vspace{9pt}\popsemi\fontsize{10.6}{14.5}\selectfont\color{popink}}}

% Signature de fin de cours — poussée tout en bas de la dernière page
\newcommand{\findecours}{%
  \par\vfill\nopagebreak
  \begin{center}\popsquiggle{poporange}\end{center}\vspace{4pt}
  \signatureonbuch
}
\AtBeginDocument{\def\llbracket{\mathopen{[\mkern-3.5mu[}}\def\rrbracket{\mathclose{]\mkern-3.5mu]}}} % intervalles d'entiers (glyphe ⟦ absent de la police)
% Glyphes absents de Fira Math : redéfinis à partir de glyphes disponibles
\AtBeginDocument{%
  \def\setminus{\mathbin{\backslash}}\let\smallsetminus\setminus
  \def\vdots{\mathord{\vbox{\baselineskip3.2pt\lineskiplimit0pt\kern4pt\hbox{.}\hbox{.}\hbox{.}}}}%
  \def\ddots{\mathinner{\mkern1mu\raise6.5pt\hbox{.}\mkern2mu\raise3.6pt\hbox{.}\mkern2mu\raise0.7pt\hbox{.}\mkern1mu}}%
  \def\triangle{\mathord{\scalebox{0.9}{$\symup{\Delta}$}}}%
  \def\nabla{\mathord{\raisebox{\depth}{\scalebox{1}[-1]{$\symup{\Delta}$}}}}%
  \def\checkmark{\popcheck{popdark}}%
  \def\star{\mathbin{\ast}}\def\bigstar{\mathord{\ast}}%
  \def\diamond{\mathbin{\scalebox{0.6}{$\Diamond$}}}%
  \def\bigcup{\mathop{\mathchoice{\raisebox{-0.3em}{\scalebox{1.7}{$\cup$}}}{\scalebox{1.25}{$\cup$}}{\cup}{\cup}}\displaylimits}%
  \def\bigcap{\mathop{\mathchoice{\raisebox{-0.3em}{\scalebox{1.7}{$\cap$}}}{\scalebox{1.25}{$\cap$}}{\cap}{\cap}}\displaylimits}%
  \def\lesseqqgtr{\mathrel{\lessgtr}}\def\bigoplus{\mathop{\oplus}\displaylimits}%
}
\providecommand{\ssi}{si et seulement si} % abréviation fréquente
\AtBeginDocument{\providecommand{\wideparen}{\overparen}} % arc de cercle

%% FIGFIX-BEGIN v5 — correctifs globaux des figures (docs/onbuch-generation/tools/figfix)
\usepackage{adjustbox}\usepackage{etoolbox}\tcbuselibrary{raster}
% toute figure TikZ/pgfplots plus large que la ligne est réduite à la largeur disponible
\BeforeBeginEnvironment{tikzpicture}{\begin{adjustbox}{max width=\linewidth}}
\AfterEndEnvironment{tikzpicture}{\end{adjustbox}}
% légendes pgfplots lisibles par-dessus les courbes
\pgfplotsset{every axis/.append style={legend style={fill=white,fill opacity=0.92,text opacity=1,draw=black!15,inner sep=3pt}}}
% étiquettes d'arêtes (syntaxe edge["..."]) avec fond blanc
\tikzset{every edge quotes/.append style={fill=white,inner sep=1.2pt}}
%% FIGFIX-END
\begin{document}

\pagegarde{L'article de journal}{Allemand}{1ère A}{Chapitre 4 : Média et communication}{ALA18}{2 h}{Cette leçon propose la lecture et l'exploitation d'un article de journal allemand : découverte du vocabulaire de la presse, compréhension globale et détaillée, révision du parfait et du prétérit ainsi que des subordonnées, puis production guidée d'un court article.
\vspace{4pt}
\begin{multicols}{2}\small
\begin{itemize}
\item À la fin de cette leçon, je sais identifier la structure d'un article de journal allemand (titre, chapeau, corps, signature).
\item À la fin de cette leçon, je sais comprendre globalement puis en détail un article de presse simple.
\item À la fin de cette leçon, je sais employer le vocabulaire thématique de la presse (der Artikel, die Überschrift, berichten, das Ereignis...) avec article et pluriel.
\item À la fin de cette leçon, je sais distinguer et utiliser correctement le parfait et le prétérit pour raconter des faits passés.
\item À la fin de cette leçon, je sais construire des subordonnées avec weil, dass, als, wenn en plaçant le verbe à la fin.
\item À la fin de cette leçon, je sais répondre par écrit et à l'oral à des questions de compréhension.
\end{itemize}\end{multicols}}

\begin{sommaire}
\begin{multicols}{2}
\begin{enumerate}
\item Le texte support : un article de journal
\item Compréhension du texte
\item Lexique thématique : la presse et l'article
\item Grammaire 1 : le parfait et le prétérit
\item Grammaire 2 : les subordonnées
\item Méthode du commentaire et commentaire modèle
\item Production guidée : écrire un court article
\item Activité d'intégration
\item Méthodes et exercices résolus
\item Exercices
\item Corrigés des exercices
\item Fiche bilan
\end{enumerate}
\end{multicols}
\end{sommaire}

\begin{objectifs}
\begin{itemize}
\item À la fin de cette leçon, je sais identifier la structure d'un article de journal allemand (titre, chapeau, corps, signature).
\item À la fin de cette leçon, je sais comprendre globalement puis en détail un article de presse simple.
\item À la fin de cette leçon, je sais employer le vocabulaire thématique de la presse (der Artikel, die Überschrift, berichten, das Ereignis...) avec article et pluriel.
\item À la fin de cette leçon, je sais distinguer et utiliser correctement le parfait et le prétérit pour raconter des faits passés.
\item À la fin de cette leçon, je sais construire des subordonnées avec weil, dass, als, wenn en plaçant le verbe à la fin.
\item À la fin de cette leçon, je sais répondre par écrit et à l'oral à des questions de compréhension.
\item À la fin de cette leçon, je sais résumer un article en quelques phrases.
\item À la fin de cette leçon, je sais rédiger un court article de journal guidé sur un événement local.
\end{itemize}
\end{objectifs}

\begin{prerequis}
\begin{itemize}
\item Le présent et les auxiliaires haben/sein (Seconde)
\item Les participes passés réguliers et irréguliers de base (Seconde)
\item La place du verbe en position 2 dans la phrase principale (Seconde)
\item Les temps du passé en français (passé composé / imparfait)
\item Le vocabulaire général des médias vu en début d'année (Chapitre 4)
\end{itemize}
\end{prerequis}

\newpage

\coursec{Le texte support : un article de journal}

Chaque matin à Yaoundé, des milliers de Camerounais ouvrent \emph{Cameroon Tribune} ou écoutent la CRTV pour s'informer. En Allemagne, en Autriche et en Suisse, c'est le même rituel : on lit la \emph{Süddeutsche Zeitung}, le \emph{Standard} ou le \emph{Tages-Anzeiger} au petit-déjeuner, dans le tramway ou sur son téléphone. Mais comment un article de journal est-il construit ? Qui écrit, qui lit, que rapporte-t-on ? Et surtout : comment raconter des faits passés de manière claire et objective ?

Dans cette leçon, tu vas lire un vrai article de presse rédigé en allemand sur un événement sportif au Cameroun, découvrir la structure du genre journalistique et apprendre à raconter des faits passés comme un journaliste. La problématique qui guidera notre travail est la suivante : \textbf{comment l'article de journal informe-t-il le lecteur de manière rapide, complète et objective ?}

\courssub{Qu'est-ce qu'un article de journal ?}

\begin{definition}[Der Zeitungsartikel]
Un \cle{article de journal} (\all{der Zeitungsartikel}, pluriel \all{die Zeitungsartikel}) est un texte journalistique qui \textbf{informe} le lecteur sur un \textbf{événement récent} de manière \textbf{objective} et \textbf{structurée}. Il est écrit par un \cle{journaliste} (\all{der Journalist}) et lu par des \cle{lecteurs} (\all{die Leser}). Contrairement au récit littéraire, l'article ne cherche pas à divertir : il répond à des questions précises sur les faits.
\end{definition}

Tout article de presse, qu'il soit publié à Douala, à Berlin ou à Vienne, répond à cinq questions fondamentales. Les journalistes allemands les appellent les \all{W-Fragen} (les « questions en W », car elles commencent toutes par la lettre W) :

\begin{center}
\begin{tabularx}{\linewidth}{|X|X|X|}\hline
\rowcolor{popblueL} \textbf{Question allemande} & \textbf{Français} & \textbf{Information cherchée} \\ \hline
\all{Wer?} & Qui ? & les personnes concernées \\ \hline
\all{Was?} & Quoi ? & l'événement, les faits \\ \hline
\all{Wo?} & Où ? & le lieu \\ \hline
\all{Wann?} & Quand ? & la date, le moment \\ \hline
\all{Warum?} & Pourquoi ? & la cause, l'explication \\ \hline
\end{tabularx}
\end{center}

\begin{popfigure}
\begin{tikzpicture}[mindmap, grow cyclic, every node/.style={concept}, root concept/.append style={concept color=popblue, font=\bfseries, text=white}, level 1 concept/.append style={sibling angle=72, level distance=3.2cm, concept color=poporange, text=popdark}]
\node[root concept]{Der Artikel};
\node[align=center]{\all{Wer?}\\ qui ?};
\node[align=center]{\all{Was?}\\ quoi ?};
\node[align=center]{\all{Wo?}\\ où ?};
\node[align=center]{\all{Wann?}\\ quand ?};
\node[align=center]{\all{Warum?}\\ pourquoi ?};
\end{tikzpicture}
\legende{Carte mentale : les cinq questions du journaliste (die W-Fragen)}
\end{popfigure}

\courssub{La structure d'un article de journal}

Observe une page de journal : les textes ne sont pas disposés n'importe comment. Un article allemand suit une architecture très codifiée, que le lecteur reconnaît immédiatement.

\begin{propriete}[Les quatre parties de l'article]
\begin{enumerate}
\item \cle{die Überschrift} (le titre) : courte, frappante, elle annonce le sujet et attire l'œil. Elle est souvent au présent ou sans verbe conjugué.
\item \cle{der Vorspann} (le chapeau) : un ou deux phrases en caractères gras sous le titre ; il résume l'essentiel (souvent le lieu et la date : \all{Yaoundé. Am Samstag ...}).
\item \cle{der Textkörper} (le corps du texte) : le développement, qui détaille les faits, cite les témoins et donne les explications.
\item \cle{die Quelle / die Unterschrift} (la source / la signature) : le nom du journaliste ou de l'agence, parfois le lieu de rédaction.
\end{enumerate}
\end{propriete}

\begin{popfigure}
\begin{tikzpicture}[node distance=0.35cm]
\node[draw=popblue, fill=popblueL, rounded corners, minimum width=9cm, minimum height=1cm, font=\bfseries] (titre) {\all{die Überschrift} — le titre};
\node[draw=poporange, fill=poporangeL, rounded corners, minimum width=9cm, minimum height=1cm, below=of titre] (chapeau) {\all{der Vorspann} — le chapeau (lieu, date, résumé)};
\node[draw=popgreen, fill=popgreenL, rounded corners, minimum width=9cm, minimum height=1.6cm, below=of chapeau, align=center] (corps) {\all{der Textkörper} — le corps du texte\\ (faits, citations, explications)};
\node[draw=popgold, fill=popgoldL, rounded corners, minimum width=9cm, minimum height=1cm, below=of corps] (source) {\all{die Quelle / die Unterschrift} — la source, la signature};
\draw[-{Stealth}, thick, popmuted] (titre.west) ++(-0.4,0) -- ++(0,-4.6) node[midway, left, align=center, font=\small]{du plus\\ important\\ au plus\\ détaillé};
\end{tikzpicture}
\legende{Schéma : la structure d'un article de journal allemand}
\end{popfigure}

\begin{attention}[Le titre n'est pas une phrase complète]
En allemand comme en français, le titre d'un article est souvent \textbf{elliptique} : on supprime l'article défini ou le verbe. \all{Jungen aus Yaoundé gewinnen Fußballturnier} signifie en réalité \all{Die Jungen aus Yaoundé haben das Fußballturnier gewonnen}. Ne cherche donc pas toujours un verbe conjugué dans le titre !
\end{attention}

\courssub{Lecture globale : première approche du texte}

\begin{methode}[Lire un article en trois étapes]
\begin{enumerate}
\item \textbf{Lire le titre et le chapeau} pour deviner le sujet : de quoi parle l'article ? Où et quand se passe l'action ?
\item \textbf{Lire le corps du texte} en repérant les cinq W : souligne les noms de personnes, de lieux, les dates et les chiffres.
\item \textbf{Vérifier les mots inconnus} d'abord au glossaire, et seulement ensuite au dictionnaire : beaucoup de mots se devinent par le contexte.
\end{enumerate}
\end{methode}

Applique cette méthode maintenant. Lis le texte suivant une première fois \textbf{en silence}, sans t'arrêter sur les mots inconnus, et cherche simplement : de quoi parle cet article ?

\begin{textebox}[Jungen aus Yaoundé gewinnen Fußballturnier]
\textbf{Yaoundé.} Am Samstag hat das große Schulfußballturnier in Yaoundé stattgefunden. Zehn Schulen aus der ganzen Stadt haben an dem Turnier teilgenommen. Die Mannschaft des Lycée de Nkol-Eton hat das Finale 3:1 gewonnen. Viele Zuschauer kamen, weil das Wetter schön war. Eltern, Lehrer und Schüler haben laut applaudiert. Der Trainer berichtete, dass die Jungen jeden Tag trainiert hatten. „Wir sind sehr stolz“, sagte der Kapitän nach dem Spiel. „Dieser Sieg ist das Ergebnis von harter Arbeit.“ Auch der Direktor der Schule gratulierte der Mannschaft. Das Ereignis war ein großer Erfolg für die ganze Schule. Nächstes Jahr will das Lycée den Titel verteidigen.
\end{textebox}

\textbf{Glossaire} (les mots difficiles du texte) :

\begin{center}
\begin{tabularx}{\linewidth}{|X|X|}\hline
\rowcolor{popblueL} \textbf{Deutsch} & \textbf{Français} \\ \hline
\all{das Schulfußballturnier (-e)} & le tournoi de football scolaire \\ \hline
\all{stattfinden (fand statt, hat stattgefunden)} & avoir lieu, se dérouler \\ \hline
\all{teilnehmen (nahm teil, hat teilgenommen)} & participer \\ \hline
\all{die Mannschaft (-en)} & l'équipe \\ \hline
\all{das Finale (-)} & la finale \\ \hline
\all{gewinnen (gewann, hat gewonnen)} & gagner \\ \hline
\all{der Zuschauer (-)} & le spectateur \\ \hline
\all{applaudieren} & applaudir \\ \hline
\all{der Trainer (-)} & l'entraîneur \\ \hline
\all{berichten} & rapporter, raconter, informer \\ \hline
\all{trainieren} & s'entraîner \\ \hline
\all{stolz} & fier \\ \hline
\all{der Kapitän (-e)} & le capitaine \\ \hline
\all{der Sieg (-e)} & la victoire \\ \hline
\all{das Ergebnis (-se)} & le résultat \\ \hline
\all{gratulieren (+ Dativ)} & féliciter \\ \hline
\all{das Ereignis (-se)} & l'événement \\ \hline
\all{der Erfolg (-e)} & le succès \\ \hline
\all{verteidigen} & défendre \\ \hline
\end{tabularx}
\end{center}

\begin{attention}[Ne traduis pas mot à mot !]
Un article se comprend d'abord \textbf{globalement}. Les mots apparentés (les « cognats ») t'aident énormément : \all{das Turnier} = le tournoi, \all{der Erfolg} ≈ le succès, \all{das Finale} = la finale, \all{der Kapitän} = le capitaine, \all{gratulieren} = féliciter. Commence par ces mots familiers, puis devine le reste grâce au contexte. Le dictionnaire n'est que le dernier recours.
\end{attention}

\courssub{Lecture détaillée : les cinq W dans le texte}

Relis maintenant l'article une deuxième fois, plus lentement, et souligne les \textbf{temps du passé} (parfait et prétérit — nous les étudierons en détail dans la partie grammaire). Vérifie ensuite que tu as trouvé les réponses aux cinq questions du journaliste :

\begin{center}
\begin{tabularx}{\linewidth}{|X|X|}\hline
\rowcolor{popblueL} \textbf{W-Frage} & \textbf{Réponse dans le texte} \\ \hline
\all{Wer?} & l'équipe du Lycée de Nkol-Eton, les spectateurs, l'entraîneur, le capitaine \\ \hline
\all{Was?} & un tournoi de football scolaire, une victoire 3:1 en finale \\ \hline
\all{Wo?} & à Yaoundé \\ \hline
\all{Wann?} & samedi dernier \\ \hline
\all{Warum?} & victoire grâce à un entraînement quotidien ; beaucoup de spectateurs grâce au beau temps \\ \hline
\end{tabularx}
\end{center}

Observe aussi la \textbf{structure} : la première phrase (\all{Am Samstag hat das große Schulfußballturnier ... stattgefunden}) joue le rôle de chapeau : elle donne déjà le lieu, la date et l'événement. Le corps du texte ajoute ensuite les détails, les citations (\all{„Wir sind sehr stolz“}) et une ouverture vers l'avenir (\all{Nächstes Jahr will das Lycée den Titel verteidigen}).

\begin{exoresolu}[Repérer les informations essentielles]
\textbf{Énoncé.} Réponds en français, en une phrase complète, aux questions suivantes : 1) Quel événement l'article rapporte-t-il ? 2) Quelle équipe a gagné et sur quel score ? 3) Pourquoi les spectateurs étaient-ils nombreux ?
\tcblower
\textbf{Solution détaillée.}
1) L'article rapporte le grand tournoi de football scolaire qui a eu lieu samedi à Yaoundé (\all{Am Samstag hat das große Schulfußballturnier in Yaoundé stattgefunden}).
2) C'est l'équipe du Lycée de Nkol-Eton qui a gagné la finale sur le score de 3 à 1 (\all{Die Mannschaft des Lycée de Nkol-Eton hat das Finale 3:1 gewonnen}).
3) Les spectateurs étaient nombreux parce qu'il faisait beau (\all{Viele Zuschauer kamen, weil das Wetter schön war}).
\end{exoresolu}

\begin{exercice}[À toi de jouer]
1) Trouve dans le texte trois verbes au parfait (\all{hat ... stattgefunden}, ...) et deux verbes au prétérit.
2) Relève la citation du capitaine : comment reconnaît-on qu'il s'agit de paroles rapportées au discours direct ?
3) Quelle phrase du texte annonce l'avenir ? Quel verbe exprime le projet ?
\end{exercice}

\begin{corrige}[Exercice « À toi de jouer »]
1) Parfait : \all{hat ... stattgefunden}, \all{hat ... teilgenommen}, \all{hat ... gewonnen} (aussi \all{haben ... applaudiert}). Prétérit : \all{kamen}, \all{berichtete}, \all{sagte}, \all{war}, \all{gratulierte}.
2) La citation \all{„Wir sind sehr stolz“} est entre guillemets allemands „...“ et introduite par le verbe déclaratif \all{sagte} : c'est le discours direct.
3) \all{Nächstes Jahr will das Lycée den Titel verteidigen} : le verbe modal \all{wollen} (\all{will}) exprime le projet, l'intention.
\end{corrige}

\begin{savaistu}[La presse germanophone]
La presse allemande est l'une des plus lues d'Europe. Des journaux comme la \emph{Frankfurter Allgemeine Zeitung}, la \emph{Süddeutsche Zeitung} ou le magazine d'information \emph{Der Spiegel} sont connus dans le monde entier pour la qualité de leurs enquêtes. En Autriche, on lit \emph{Der Standard} et \emph{Die Presse} ; en Suisse, le \emph{Tages-Anzeiger} et la \emph{Neue Zürcher Zeitung}. Et le lien avec le Cameroun ? De jeunes journalistes camerounais germanophones travaillent aujourd'hui pour des rédactions bilingues, et l'allemand reste la langue étrangère qui ouvre le plus de portes dans les métiers de la communication en Afrique centrale.
\end{savaistu}

\begin{aretenir}[L'essentiel de la section]
\begin{itemize}
\item Un article de journal (\all{der Zeitungsartikel}) informe objectivement sur un événement récent.
\item Il répond aux cinq questions : \all{Wer? Was? Wo? Wann? Warum?}
\item Sa structure : \all{die Überschrift} (titre) → \all{der Vorspann} (chapeau) → \all{der Textkörper} (corps) → \all{die Quelle} (source).
\item Lecture efficace : d'abord globale (titre + chapeau), puis détaillée (les cinq W), en s'aidant des mots apparentés avant le dictionnaire.
\item Le vocabulaire clé : \all{der Artikel (-)}, \all{die Überschrift (-en)}, \all{der Leser (-)}, \all{berichten}, \all{das Ereignis (-se)}, \all{das Turnier (-e)}, \all{gewinnen}, \all{stattfinden}, \all{der Sieg (-e)}, \all{die Mannschaft (-en)}.
\end{itemize}
\end{aretenir}

\coursec{Compréhension du texte}

Dans la section précédente, nous avons lu l'article de journal \all{„Jugendturnier in Yaoundé: Ein großer Erfolg“} et découvert sa structure : le titre (\all{die Überschrift}), le lieu et la date, le corps de l'article avec les faits et les citations. Mais lire un article, ce n'est pas seulement le déchiffrer : c'est le \cle{comprendre}, c'est-à-dire être capable de dire de quoi il parle, de répondre à des questions précises, de vérifier des affirmations et de le résumer. C'est exactement ce que l'on vous demandera au baccalauréat. Cette section vous donne les outils et les techniques pour y arriver.

\begin{definition}[Vocabulaire clé : la presse et l'article]
Les mots suivants figurent au programme officiel pour ce thème. Apprenez-les avec leur article, leur pluriel et leur traduction :
\begin{center}
\begin{tabularx}{\linewidth}{|l|X|l|}
\hline
\rowcolor{popblueL} \textbf{Allemand} & \textbf{Français} & \textbf{Information} \\
\hline
\all{der Artikel, die Artikel} & l'article (de journal) & nom masculin \\
\hline
\all{der Titel, die Titel} & le titre (général) & nom masculin \\
\hline
\all{die Überschrift, die Überschriften} & le titre (de l'article), la manchette & nom féminin \\
\hline
\all{der Leser, die Leser} & le lecteur & nom masculin \\
\hline
\all{die Leserin, die Leserinnen} & la lectrice & nom féminin (\emph{-in}) \\
\hline
\all{berichten (berichtete, hat berichtet)} & rapporter, faire un reportage & verbe faible + \all{haben} \\
\hline
\all{das Ereignis, die Ereignisse} & l'événement, le fait divers & nom neutre (\emph{-e}) \\
\hline
\all{die Meldung, die Meldungen} & la dépêche, la brève & nom féminin \\
\hline
\all{der Bericht, die Berichte} & le reportage, le compte-rendu & nom masculin \\
\hline
\all{erscheinen (erschien, ist erschienen)} & paraître, être publié & verbe fort + \all{sein} \\
\hline
\end{tabularx}
\end{center}
\end{definition}

\courssub{La compréhension globale : de quoi parle l'article ?}

Avant de répondre à des questions détaillées, il faut d'abord saisir le \cle{sens global} du texte. Pour cela, on se pose trois questions simples :

\begin{enumerate}
\item \textbf{De quoi parle l'article ?} (le sujet, \all{das Thema})
\item \textbf{Où se passe l'action ?} (le lieu, \all{der Ort})
\item \textbf{Quand se passe l'action ?} (le moment, \all{die Zeit})
\end{enumerate}

Reprenons notre article. Une première lecture rapide (sans dictionnaire, sans s'arrêter sur les mots inconnus) suffit pour répondre :

\begin{exemplebox}[Compréhension globale — réponses modèles]
\begin{itemize}
\item \all{Der Artikel berichtet über ein Fußballturnier der Schulen.} — L'article parle d'un tournoi de football des écoles.
\item \all{Das Turnier fand in Yaoundé statt, auf dem Sportplatz des Gymnasiums.} — Le tournoi a eu lieu à Yaoundé, sur le terrain de sport du lycée.
\item \all{Es fand am Samstag statt.} — Il a eu lieu samedi.
\end{itemize}
\end{exemplebox}

\begin{methode}[Lire un article en trois étapes]
\begin{enumerate}
\item \textbf{Première lecture} : lire le titre et le premier paragraphe seulement. Le titre annonce presque toujours le sujet.
\item \textbf{Deuxième lecture} : lire tout le texte rapidement pour repérer le lieu, le moment et les personnes.
\item \textbf{Troisième lecture} : lire lentement, paragraphe par paragraphe, en soulignant les mots-clés (noms propres, dates, verbes d'action).
\end{enumerate}
\end{methode}

\courssub{La compréhension détaillée : répondre aux W-Fragen}

Les questions de compréhension en allemand commencent presque toujours par un mot interrogatif en \textbf{W-} : on les appelle les \all{W-Fragen}. Les voici, avec leur équivalent français :

\begin{center}
\begin{tabularx}{\linewidth}{|l|X|X|}
\hline
\rowcolor{popblueL} \textbf{Mot interrogatif} & \textbf{Français} & \textbf{Ce qu'il demande} \\
\hline
\all{Wer?} & Qui ? & une personne (sujet) \\
\hline
\all{Was?} & Quoi ? Que ? & une chose, une action (objet) \\
\hline
\all{Wann?} & Quand ? & un moment, une date \\
\hline
\all{Wo?} & Où ? & un lieu (position) \\
\hline
\all{Wohin?} & Où ? (direction) & une direction (mouvement vers) \\
\hline
\all{Woher?} & D'où ? (origine) & une provenance (mouvement depuis) \\
\hline
\all{Warum?} & Pourquoi ? & une raison, une cause \\
\hline
\all{Wie?} & Comment ? & une manière, un état \\
\hline
\end{tabularx}
\end{center}

\begin{attention}[La place du verbe dans la question]
En allemand, dans une \all{W-Frage}, le verbe conjugué occupe la \textbf{deuxième position}, juste après le mot interrogatif : \all{Wann \textbf{hat} das Turnier stattgefunden?} Attention à l'ordre des mots dans la réponse : le verbe reste aussi en deuxième position — \all{Das Turnier \textbf{hat} am Samstag stattgefunden.} Ne dites jamais \all{*Das Turnier am Samstag hat stattgefunden}.
\end{attention}

\begin{methode}[Répondre à une W-Frage : la technique de la reprise]
Pour répondre en phrase complète, on \textbf{reprend la structure de la question} et on remplace le mot interrogatif par l'information demandée :
\begin{enumerate}
\item \all{Wann hat das Turnier stattgefunden?} (Quand le tournoi a-t-il eu lieu ?)
\item On reprend : \all{Das Turnier hat \ldots\ stattgefunden.}
\item On insère la réponse à la place du mot interrogatif : \all{Das Turnier hat \textbf{am Samstag} stattgefunden.}
\end{enumerate}
Cette technique garantit une phrase correcte : sujet, verbe en deuxième position, participe (ou infinitif) à la fin.
\end{methode}

\begin{popfigure}
\begin{center}
\begin{tikzpicture}[
  frage/.style={draw=popblue, fill=popblueL, rounded corners, align=center, text width=4.6cm, font=\small},
  antwort/.style={draw=popgreen, fill=popgreenL, rounded corners, align=center, text width=4.6cm, font=\small},
  pfeil/.style={-{Stealth[length=3mm]}, thick, poporange}
]
\node[frage, align=center] (f1) at (0,3){\textbf{Frage}\\ Wann hat das Turnier stattgefunden?};
\node[antwort, align=center] (a1) at (7,3){\textbf{Antwort}\\ Das Turnier hat am Samstag stattgefunden.};
\node[frage, align=center] (f2) at (0,1.2){\textbf{Frage}\\ Wer hat gewonnen?};
\node[antwort, align=center] (a2) at (7,1.2){\textbf{Antwort}\\ Die Mannschaft des Gymnasiums hat gewonnen.};
\node[frage, align=center] (f3) at (0,-0.6){\textbf{Frage}\\ Warum kamen viele Zuschauer?};
\node[antwort, align=center] (a3) at (7,-0.6){\textbf{Antwort}\\ Viele Zuschauer kamen, weil das Finale sehr spannend war.};
\draw[pfeil] (f1) -- (a1);
\draw[pfeil] (f2) -- (a2);
\draw[pfeil] (f3) -- (a3);
\end{tikzpicture}
\end{center}
\legende{De la question à la réponse complète : on reprend les mots de la question.}
\end{popfigure}

\begin{exoresolu}[Questions de compréhension sur l'article]
Répondez en phrases complètes : 
\begin{enumerate}
\item \all{Wann hat das Turnier stattgefunden?} 
\item \all{Wer hat gewonnen?} 
\item \all{Warum kamen viele Zuschauer?} 
\item \all{Was sagte der Kapitän?}
\end{enumerate}
\tcblower
\textbf{Solutions :}
\begin{enumerate}
\item \all{Das Turnier hat am Samstag auf dem Sportplatz des Gymnasiums in Yaoundé stattgefunden.} — Le tournoi a eu lieu samedi sur le terrain de sport du lycée de Yaoundé.
\item \all{Die Mannschaft des Gymnasiums Yaoundé hat gewonnen.} — L'équipe du lycée de Yaoundé a gagné. (On reprend \all{wer} par le sujet de la réponse.)
\item \all{Viele Zuschauer kamen, weil das Finale sehr spannend war und weil die Schüler ihre Mannschaft unterstützen wollten.} — Beaucoup de spectateurs sont venus parce que la finale était très passionnante et parce que les élèves voulaient soutenir leur équipe. (Après \all{warum}, on répond avec \all{weil} + verbe à la fin.)
\item \all{Der Kapitän sagte, dass die Mannschaft hart trainiert \textbf{habe} und dass er stolz auf seine Spieler \textbf{sei}.} — Le capitaine a dit que l'équipe s'était entraînée dur et qu'il était fier de ses joueurs. (Après \all{was}, on rapporte les paroles avec \all{dass} + \textbf{Konjunktiv I} : forme standard du discours indirect écrit.)
\end{enumerate}
\end{exoresolu}

\begin{savaistu}[Le Konjunktiv I dans le discours indirect]
En allemand écrit (journalisme, rapports), le discours indirect utilise le \textbf{Konjunktiv I} (subjonctif I) pour marquer la distance par rapport aux propos tenus : \all{Er sagte: „Ich \textbf{bin} stolz.“} $\rightarrow$ \all{Er sagte, dass er stolz \textbf{sei}.} (au lieu de \all{ist}). Pour \all{haben} : \all{habe, hast, habe, haben, habt, haben}. Pour \all{sein} : \all{sei, seist, sei, seien, seid, seien}. Au niveau A2, l'indicatif (\all{hat/ist}) est toléré à l'oral, mais l'écrit exige le Konjunktiv I.
\end{savaistu}

\begin{attention}[Répondre à \all{Warum?} : le verbe à la fin]
Quand on répond à une question en \all{warum}, la réponse commence par \all{weil} (parce que), et le verbe conjugué va \textbf{à la fin} de la phrase : \all{Viele Zuschauer kamen, weil das Finale spannend \textbf{war}.} Ne calquez pas le français « parce que la finale était passionnante » mot à mot : \all{*weil das Finale war spannend} est une faute grave (ordre des mots principal).
\end{attention}

\courssub{Vrai ou faux : richtig oder falsch ?}

L'exercice \all{richtig/falsch} est un classique des épreuves de compréhension. Il consiste à lire une affirmation et à dire si elle est vraie (\all{richtig}) ou fausse (\all{falsch}) \textbf{d'après le texte} — et non d'après vos connaissances personnelles.

\begin{attention}[Justifier, toujours justifier]
Au baccalauréat, écrire seulement \all{richtig} ou \all{falsch} ne suffit pas : il faut \textbf{toujours justifier par une phrase du texte} (une citation entre guillemets allemands „ … “). Une réponse sans justification perd la moitié des points, même si elle est correcte.
\end{attention}

\begin{exemplebox}[Modèle de réponse richtig/falsch justifiée]
\textbf{Affirmation :} \all{Das Turnier fand am Sonntag statt.}\\
\textbf{Réponse :} \all{Falsch. Im Text steht: „Das Turnier fand am \textbf{Samstag} statt.“}\\
— Faux. Le texte dit : « Le tournoi a eu lieu samedi. »\\[0.3em]
\textbf{Affirmation :} \all{Viele Zuschauer kamen zum Finale.}\\
\textbf{Réponse :} \all{Richtig. Der Artikel sagt: „Hunderte Zuschauer kamen auf den Sportplatz.“}\\
— Vrai. L'article dit : « Des centaines de spectateurs sont venus sur le terrain. »
\end{exemplebox}

\begin{propriete}[Expressions utiles pour justifier]
Pour citer le texte, utilisez ces formules toutes faites :
\begin{itemize}
\item \all{Im Text steht: „\ldots“} — Dans le texte, il est écrit : « … » (verbe \all{stehen} impersonnel, toujours 3\textsuperscript{e} pers. sg.)
\item \all{Der Artikel sagt / berichtet: „\ldots“} — L'article dit / rapporte : « … »
\item \all{Der Autor schreibt: „\ldots“} — L'auteur écrit : « … »
\item \all{Laut Text \ldots} — Selon le texte … (suivi de l'indicatif, verbe en 2\textsuperscript{e} position)
\end{itemize}
\end{propriete}

\begin{exercice}[Richtig oder falsch ?]
Dites si les affirmations suivantes sont \all{richtig} ou \all{falsch} et justifiez par une citation du texte.
\begin{enumerate}
\item \all{Die Mannschaft des Gymnasiums hat verloren.}
\item \all{Der Kapitän war stolz auf seine Mannschaft.}
\item \all{Keine Zuschauer kamen zum Spiel.}
\end{enumerate}
\end{exercice}

\begin{corrige}[Exercice richtig/falsch]
\begin{enumerate}
\item \all{Falsch. Im Text steht: „Die Mannschaft des Gymnasiums hat das Finale gewonnen.“} — L'équipe n'a pas perdu, elle a gagné la finale.
\item \all{Richtig. Der Kapitän sagte: „Ich bin stolz auf meine Spieler.“} — Le capitaine a déclaré être fier de ses joueurs (citation directe = indicatif).
\item \all{Falsch. Der Artikel berichtet: „Viele Zuschauer kamen, um das Finale zu sehen.“} — De nombreux spectateurs sont venus voir la finale (construction \all{um ... zu} + infinitif).
\end{enumerate}
\end{corrige}

\courssub{Le résumé guidé : garder l'essentiel}

Résumer (\all{zusammenfassen}), c'est dire l'essentiel d'un texte en peu de mots, \textbf{avec ses propres phrases}. C'est une compétence clé pour le baccalauréat et pour la vie : savoir raconter un article à un ami en trois phrases.

\begin{methode}[Résumer un article en trois étapes]
\begin{enumerate}
\item \textbf{Garder uniquement les 5 W} : \all{Wer?} (qui), \all{Was?} (quoi), \all{Wann?} (quand), \all{Wo?} (où), \all{Warum?} (pourquoi).
\item \textbf{Supprimer les détails} : les citations directes, les chiffres secondaires, les descriptions, les adjectifs superflus.
\item \textbf{Reformuler} avec ses propres mots, en phrases simples au Perfekt (langage oral/rapport) ou au Präteritum (récit journalistique).
\end{enumerate}
\end{methode}

\begin{exoresolu}[Résumé guidé : complétez les phrases]
Complétez les trois phrases suivantes pour résumer l'article avec vos propres mots :
\begin{enumerate}
\item \all{Der Artikel berichtet über \ldots}
\item \all{Das Turnier fand \ldots\ statt, und gewonnen hat \ldots}
\item \all{Der Kapitän sagte, dass \ldots}
\end{enumerate}
\tcblower
\textbf{Solution possible :}
\begin{enumerate}
\item \all{Der Artikel berichtet über ein Fußballturnier der Schulen in Yaoundé.} — L'article parle d'un tournoi de football scolaire à Yaoundé.
\item \all{Das Turnier fand am Samstag auf dem Sportplatz des Gymnasiums statt, und gewonnen hat die Mannschaft des Gymnasiums.} — Le tournoi a eu lieu samedi sur le terrain du lycée, et c'est l'équipe du lycée qui a gagné (topicalisation de \all{gewonnen} pour l'insister).
\item \all{Der Kapitän sagte, dass seine Mannschaft hart trainiert \textbf{habe} und dass er sehr stolz \textbf{sei}.} — Le capitaine a dit que son équipe s'était beaucoup entraînée et qu'il était très fier (Konjunktiv I dans le discours indirect).
\end{enumerate}
Remarquez que le résumé ne contient ni citation directe, ni score exact, ni description du public : seulement les 5 W.
\end{exoresolu}

\courssub{Compréhension orale : noter les 5 W}

En classe, l'enseignant lit à voix haute un second court article. Votre tâche : écouter \textbf{sans lire}, et noter les 5 W sur votre cahier. C'est l'exercice de \cle{compréhension orale} (\all{das Hörverstehen}).

\begin{methode}[Réussir la compréhension orale]
\begin{enumerate}
\item Avant l'écoute, écrivez sur votre cahier les cinq mots : \all{Wer? Was? Wann? Wo? Warum?}
\item À la première écoute, notez seulement des mots-clés (noms, lieux, dates), pas des phrases.
\item À la deuxième écoute, complétez vos notes et vérifiez les 5 W.
\item Formulez ensuite vos réponses en phrases complètes, comme pour les questions écrites.
\end{enumerate}
\end{methode}

\begin{textebox}[Texte pour la compréhension orale : „Neue Bücher für die Schulbibliothek“]
Die Schüler des Gymnasiums in Douala haben am Freitag eine große Überraschung bekommen: Eine deutsche Organisation hat der Schule zweihundert neue Bücher geschenkt. Die Bücher sind Romane, Sachbücher und Wörterbücher auf Deutsch und auf Englisch.

Die Direktorin, Frau Mbarga, hat die Bücher am Morgen in der Aula empfangen. „Wir sind sehr glücklich“, sagte sie. „Unsere Bibliothek war klein, aber jetzt können alle Schüler deutsche Bücher lesen.“ Auch die Deutschlehrer waren sehr zufrieden. Herr Nkomo, der älteste Lehrer der Schule, erklärte: „Lesen ist der beste Weg, eine Sprache zu lernen.“

Die Organisation hat die Bücher gekauft, weil sie die Zusammenarbeit zwischen Deutschland und Kamerun stärken will. Schon im letzten Jahr hat sie drei Schulen in Yaoundé und Bafoussam Computer geschenkt.

Am Nachmittag haben die Schüler der Klasse 1A die Bücher in die Bibliothek getragen und sie nach Themen sortiert. Ein Schüler sagte: „Ich habe schon ein Buch über den Fußball in Deutschland gefunden. Das lese ich zuerst!“

\emph{Glossaire :} \all{die Überraschung, -en} = la surprise ; \all{schenken (schenkte, hat geschenkt)} = offrir, faire cadeau de (qqch \textbf{à qqn} : Akk + Dat) ; \all{das Sachbuch, -bücher} = le documentaire, le livre documentaire ; \all{das Wörterbuch, -bücher} = le dictionnaire ; \all{empfangen (empfing, hat empfangen)} = recevoir, accueillir ; \all{die Aula, -s} = la salle des fêtes, l'amphithéâtre ; \all{der Weg, -e} = le chemin, le moyen ; \all{die Zusammenarbeit} = la coopération ; \all{stärken} = renforcer ; \all{sortieren} = trier, classer ; \all{die Klasse 1A} = la classe de 1re A.
\end{textebox}

\begin{exercice}[Les 5 W du texte oral]
Après l'écoute, répondez en phrases complètes :
\begin{enumerate}
\item \all{Wer hat Bücher bekommen?}
\item \all{Was hat die Organisation geschenkt?}
\item \all{Wann hat die Direktorin die Bücher empfangen?}
\item \all{Wo hat sie die Bücher empfangen?}
\item \all{Warum hat die Organisation die Bücher gekauft?}
\end{enumerate}
\end{exercice}

\begin{corrige}[Exercice : les 5 W]
\begin{enumerate}
\item \all{Die Schüler des Gymnasiums in Douala haben Bücher bekommen.} — Les élèves du lycée de Douala ont reçu des livres.
\item \all{Die Organisation hat zweihundert neue Bücher geschenkt.} — L'organisation a offert deux cents nouveaux livres.
\item \all{Die Direktorin hat die Bücher am Freitagmorgen empfangen.} — La directrice a reçu les livres vendredi matin.
\item \all{Sie hat die Bücher in der Aula empfangen.} — Elle les a reçus dans la salle des fêtes.
\item \all{Die Organisation hat die Bücher gekauft, weil sie die Zusammenarbeit zwischen Deutschland und Kamerun stärken will.} — L'organisation a acheté les livres parce qu'elle veut renforcer la coopération entre l'Allemagne et le Cameroun.
\end{enumerate}
\end{corrige}

\begin{aretenir}[L'essentiel de la compréhension de texte]
\begin{itemize}
\item Compréhension globale : identifier le sujet (\all{Thema}), le lieu (\all{Ort}) et le moment (\all{Zeit}) dès la première lecture.
\item Répondre à une \all{W-Frage} en reprenant la structure de la question : verbe en deuxième position.
\item \all{Warum?} $\rightarrow$ réponse avec \all{weil} + verbe \textbf{à la fin} ; \all{Was sagte X?} $\rightarrow$ réponse avec \all{dass} + \textbf{Konjunktiv I} (écrit) ou indicatif (oral).
\item \all{Richtig} ou \all{falsch} exige toujours une justification citée : \all{Im Text steht: „\ldots“}.
\item Un bon résumé ne garde que les 5 W et supprime citations, chiffres et détails.
\item À l'oral (\all{Hörverstehen}) : noter des mots-clés sous les 5 W, puis formuler des phrases complètes.
\end{itemize}
\end{aretenir}

\coursec{Lexique thématique : la presse et l'article}

Après avoir lu et compris l'article de journal de la section précédente, nous allons maintenant nous arrêter sur les mots qui composent ce texte. Un bon lecteur de presse allemande doit connaître le vocabulaire du journal : les parties du journal, les personnes qui le fabriquent et qui le lisent, les événements dont on parle, et les verbes de la communication écrite. Ce lexique est indispensable pour comprendre un article, mais aussi pour en rédiger un vous-même à la fin de la leçon.

\courssub{Observation : le vocabulaire dans son contexte}

Avant d'apprendre des listes de mots, observons-les d'abord dans un texte authentique. Lisez ce court article de presse scolaire et repérez les mots en gras.

\begin{textebox}[Ein Artikel in der Schülerzeitung]
\textbf{Neue Bibliothek am Gymnasium eröffnet}

Yaoundé. Die \textbf{Schülerzeitung} „Unsere Stimme" \textbf{berichtet} heute \textbf{über} ein wichtiges \textbf{Ereignis} an unserer Schule: Am Montag hat die Direktorin die neue Bibliothek eröffnet. Viele Schüler und Lehrer waren da.

Der \textbf{Journalist} Paul Mbarga, Schüler der Klasse Première A, hat einen langen \textbf{Artikel} über dieses Ereignis \textbf{geschrieben}. „Ich habe viele \textbf{Informationen} gesammelt und den Direktor interviewt", erklärt er. Sein Artikel steht auf der ersten \textbf{Seite} der Zeitung, direkt unter der großen \textbf{Überschrift}.

Die \textbf{Leser} sind sehr zufrieden. „Endlich gibt es gute \textbf{Nachrichten} aus unserer Schule", sagt Marie, eine Schülerin. „Ich \textbf{lese} die Zeitung jeden Monat."

Die Schülerzeitung \textbf{erscheint} am ersten Tag jedes Monats. Der \textbf{Redakteur} Herr Nkoulou \textbf{veröffentlicht} auch kurze \textbf{Berichte} über Sport und Kultur. Die nächste \textbf{Meldung}: Das Fußballteam der Schule hat am Samstag gewonnen!
\end{textebox}

\textbf{Glossaire :} die Schülerzeitung (-en) = le journal scolaire ; eröffnen = inaugurer ; sammeln = rassembler ; interviewen = interviewer ; zufrieden = satisfait ; gewinnen (gewann, hat gewonnen) = gagner.

\textbf{Questions de repérage :}
\begin{enumerate}
\item Qui a écrit l'article ? \emph{(Der Journalist Paul Mbarga.)}
\item Où se trouve l'article dans le journal ? \emph{(Auf der ersten Seite, unter der Überschrift.)}
\item Quand le journal paraît-il ? \emph{(Am ersten Tag jedes Monats.)}
\end{enumerate}

\courssub{Champ lexical 1 : le journal et ses parties}

\begin{definition}[Les parties du journal]
\begin{itemize}
\item \cle{die Zeitung} (-en) : le journal. \all{Ich lese jeden Morgen die Zeitung.} « Je lis le journal chaque matin. »
\item \cle{der Artikel} (-) : l'article. \all{Der Artikel ist sehr interessant.} « L'article est très intéressant. » Pluriel : \all{die Artikel}.
\item \cle{der Titel} (-) : le titre (d'un livre, film, chanson). \all{Der Titel des Buches ist lang.} « Le titre du livre est long. »
\item \cle{die Überschrift} (-en) : le titre, le titre en gros caractères (d'un article). \all{Die Überschrift steht über dem Text.} « Le titre se trouve au-dessus du texte. »
\item \cle{die Schlagzeile} (-n) : le gros titre, la manchette (à la une). \all{Die Schlagzeilen sind groß gedruckt.} « Les gros titres sont en gros caractères. »
\item \cle{die Seite} (-n) : la page. \all{Der Sport steht auf Seite drei.} « Le sport est à la page trois. »
\item \cle{die Zeile} (-n) : la ligne. \all{Lies bitte Zeile fünf!} « Lis la ligne cinq, s'il te plaît ! »
\end{itemize}
\end{definition}

\begin{attention}[der Titel ou die Überschrift ?]
Les deux mots signifient « le titre », mais on les emploie différemment : \all{der Titel} désigne le titre d'un livre, d'un film ou d'une chanson ; \all{die Überschrift} désigne le titre d'un article de journal, placé au-dessus du texte. Pour un article de presse, on dit donc plutôt \all{die Überschrift}. On parle aussi de \all{die Schlagzeile} (-n) pour les gros titres de la une : \all{die Schlagzeilen lesen} = « lire les gros titres ».
\end{attention}

\courssub{Champ lexical 2 : les personnes}

\begin{definition}[Les personnes de la presse]
\begin{itemize}
\item \cle{der Leser} (-) / \cle{die Leserin} (-nen) : le lecteur / la lectrice.
\item \cle{der Journalist} (-en) / \cle{die Journalistin} (-nen) : le journaliste / la journaliste.
\item \cle{der Redakteur} (-e) : le rédacteur, le rédacteur en chef.
\item \cle{der Autor} (-en) : l'auteur.
\item \cle{der Korrespondent} (-en) : le correspondant (à l'étranger).
\end{itemize}
\end{definition}

\begin{propriete}[Le féminin en -in]
En allemand, le féminin des noms de personnes se forme presque toujours en ajoutant \textbf{-in} au masculin : \all{der Leser} $\rightarrow$ \all{die Leserin}, \all{der Journalist} $\rightarrow$ \all{die Journalistin}. Au pluriel, ce féminin prend \textbf{-nen} : \all{die Leserinnen}, \all{die Journalistinnen}. Attention : \all{der Journalist} est un nom faible, il prend \textbf{-en} à tous les cas sauf au nominatif singulier : \all{Ich frage den Journalisten.} « Je demande au journaliste. »
\end{propriete}

\begin{exemplebox}[Les personnes en contexte]
\all{Die Journalistin interviewt den Bürgermeister von Douala.} « La journaliste interviewe le maire de Douala. »

\all{Der Redakteur kontrolliert alle Artikel.} « Le rédacteur vérifie tous les articles. »

\all{Viele Leser schreiben Briefe an die Zeitung.} « Beaucoup de lecteurs écrivent des lettres au journal. »
\end{exemplebox}

\courssub{Champ lexical 3 : les événements et l'information}

\begin{definition}[das Ereignis]
\cle{das Ereignis} (-se) : l'événement. Attention au genre : ce nom est \textbf{neutre} (das), et son pluriel se forme en \textbf{-se} : \all{die Ereignisse}. \all{Das Ereignis des Jahres war die Fußballmeisterschaft.} « L'événement de l'année a été le championnat de football. »
\end{definition}

\begin{definition}[Les mots de l'information]
\begin{itemize}
\item \cle{die Nachricht} (-en) : la nouvelle, l'information (à la radio, à la télévision : « les informations » = \all{die Nachrichten}).
\item \cle{die Information} (-en) : l'information, le renseignement.
\item \cle{der Bericht} (-e) : le compte rendu, le reportage.
\item \cle{die Meldung} (-en) : la dépêche, la brève, l'annonce courte.
\end{itemize}
\end{definition}

\begin{attention}[Faux amis : die Nachricht et aktuell]
\all{die Nachricht} ne signifie pas seulement « le message » (SMS, courriel) : c'est d'abord « la nouvelle, l'information journalistique ». \all{Ich habe eine gute Nachricht für dich.} « J'ai une bonne nouvelle pour toi. »

Autre piège : \all{aktuell} signifie « actuel » au sens de « d'actualité, qui se passe maintenant » (\all{aktuelle Nachrichten} = « les informations du moment »). Pour dire « actuel » au sens de « réel, véritable », l'allemand emploie \all{wirklich} ou \all{tatsächlich}. Ne traduisez jamais « le problème actuel » par \all{das aktuelle Problem} si vous voulez dire « le problème réel » !
\end{attention}

\courssub{Champ lexical 4 : les verbes de la presse}

\begin{definition}[Les verbes essentiels]
\begin{itemize}
\item \cle{berichten} (über + Akk.) : rendre compte de, rapporter.
\item \cle{informieren} (über + Akk.) : informer sur, s'informer sur (\all{sich informieren}).
\item \cle{lesen} (las, hat gelesen) : lire. Verbe fort : \all{er liest}.
\item \cle{schreiben} (schrieb, hat geschrieben) : écrire.
\item \cle{erscheinen} (erschien, ist erschienen) : paraître, être publié.
\item \cle{veröffentlichen} : publier (verbe faible, auxiliaire haben).
\end{itemize}
\end{definition}

\begin{propriete}[berichten + über + accusatif]
Le verbe \all{berichten} se construit avec la préposition \textbf{über} suivie de l'\textbf{accusatif} : \all{über ein Ereignis berichten} = « rendre compte d'un événement ». Ne dites jamais \all{berichten von} ou \all{berichten das Ereignis} directement. Même construction pour \all{informieren über + Akk.} : \all{Die Zeitung informiert über das Ereignis.}
\end{propriete}

\begin{exemplebox}[Exemples traduits]
\all{Der Journalist berichtet über das Ereignis.} « Le journaliste rend compte de l'événement. »

\all{Die Zeitung erscheint jeden Tag.} « Le journal paraît chaque jour. »

\all{Sie hat einen Artikel über den Kakaomarkt geschrieben.} « Elle a écrit un article sur le marché du cacao. »

\all{Wir informieren uns über die aktuellen Nachrichten.} « Nous nous informons sur les nouvelles du moment. »
\end{exemplebox}

\begin{attention}[erscheinen : auxiliaire sein !]
\all{erscheinen} exprime un changement d'état (le journal passe de « non publié » à « publié ») : il forme donc son parfait avec l'auxiliaire \textbf{sein} et non \textbf{haben} : \all{Die Zeitung ist gestern erschienen.} « Le journal est paru hier. » Erreur typique des francophones : \all{hat erschienen} est faux.
\end{attention}

\courssub{Les familles de mots}

En allemand, les mots d'une même famille se construisent de façon très régulière. Observer ces familles permet d'apprendre trois mots pour le prix d'un.

\begin{center}
\begin{tabularx}{\linewidth}{|l|X|X|}
\hline
\rowcolor{popblueL} Verbe & Nom d'action / personne & Traduction \\
\hline
lesen & der Leser, die Leserin, das Lesen & lire ; le lecteur, la lectrice, la lecture \\
\hline
berichten & der Bericht & rendre compte ; le compte rendu \\
\hline
informieren & die Information & informer ; l'information \\
\hline
schreiben & das Schreiben, der Autor & écrire ; l'écriture, l'auteur \\
\hline
erscheinen & das Erscheinen & paraître ; la parution \\
\hline
veröffentlichen & die Veröffentlichung & publier ; la publication \\
\hline
\end{tabularx}
\end{center}

\courssub{Expressions utiles à mémoriser}

\begin{aretenir}[Expressions de la presse]
\begin{itemize}
\item \all{einen Artikel lesen} : lire un article
\item \all{einen Artikel schreiben} : écrire un article
\item \all{über ein Ereignis berichten} : rendre compte d'un événement
\item \all{in der Zeitung stehen} : figurer dans le journal. \all{Es steht in der Zeitung.} « C'est écrit dans le journal. »
\item \all{die Schlagzeilen lesen} : lire les gros titres
\item \all{auf der ersten Seite stehen} : être en première page
\item \all{sich über etwas informieren} : s'informer sur quelque chose
\end{itemize}
\end{aretenir}

\courssub{Carte mentale : le vocabulaire de la presse}

\begin{popfigure}
\begin{center}\tcbox[colback=popblue!70,colframe=popblue!70,coltext=white,arc=9pt,boxsep=4pt,left=6pt,right=6pt]{\bfseries\small DIE PRESSE}\end{center}
\vspace{-2pt}
\begin{tcbraster}[raster columns=2,raster equal height=rows,raster column skip=6pt,raster row skip=6pt,enhanced,blankest]
\begin{tcolorbox}[enhanced,colback=poporange!60,colframe=poporange!60,coltext=white,boxrule=0.9pt,arc=3pt,left=4pt,right=4pt,top=3pt,bottom=3pt,boxsep=1pt,halign=left]\bfseries\scriptsize das Medium die Zeitung der Artikel die Überschrift die Schlagzeile die Seite\end{tcolorbox}
\begin{tcolorbox}[enhanced,colback=popgreen!60,colframe=popgreen!60,coltext=white,boxrule=0.9pt,arc=3pt,left=4pt,right=4pt,top=3pt,bottom=3pt,boxsep=1pt,halign=left]\bfseries\scriptsize die Personen der Leser der Journalist der Redakteur der Korrespondent\end{tcolorbox}
\begin{tcolorbox}[enhanced,colback=poppurple!50,colframe=poppurple!50,coltext=white,boxrule=0.9pt,arc=3pt,left=4pt,right=4pt,top=3pt,bottom=3pt,boxsep=1pt,halign=left]\bfseries\scriptsize das Ereignis die Nachricht die Information der Bericht die Meldung\end{tcolorbox}
\begin{tcolorbox}[enhanced,colback=poppink!60,colframe=poppink!60,coltext=white,boxrule=0.9pt,arc=3pt,left=4pt,right=4pt,top=3pt,bottom=3pt,boxsep=1pt,halign=left]\bfseries\scriptsize die Verben berichten informieren lesen schreiben erscheinen veröffentlichen\end{tcolorbox}
\end{tcbraster}

\legende{Carte mentale du champ lexical de la presse : quatre branches à connaître.}
\end{popfigure}

\courssub{Tableau récapitulatif}

\begin{center}
\begin{tabularx}{\linewidth}{|l|l|X|}
\hline
\rowcolor{popblueL} Deutsch & Français & Remarque \\
\hline
die Zeitung (-en) & le journal & féminin \\
\hline
der Artikel (-) & l'article & pluriel : die Artikel \\
\hline
die Überschrift (-en) & le titre (d'article) & pas « der Titel » pour un article \\
\hline
die Schlagzeile (-n) & le gros titre, la manchette & souvent au pluriel \\
\hline
der Leser / die Leserin & le lecteur / la lectrice & féminin en -in \\
\hline
der Journalist (-en) & le journaliste & nom faible : den Journalisten \\
\hline
das Ereignis (-se) & l'événement & neutre, pluriel en -se \\
\hline
die Nachricht (-en) & la nouvelle, l'info & faux ami : pas « message » \\
\hline
berichten über + Akk. & rendre compte de & préposition obligatoire \\
\hline
erscheinen (ist erschienen) & paraître & auxiliaire sein \\
\hline
\end{tabularx}
\end{center}

\courssub{Exercice résolu et pièges}

\begin{exoresolu}[Complétez avec le vocabulaire de la presse]
Énoncé. Complétez les phrases avec : \all{Überschrift}, \all{Ereignis}, \all{berichtet}, \all{erschienen}, \all{Leser}.

1. Die \_\_\_\_\_ des Artikels ist sehr lang.
2. Der Journalist \_\_\_\_\_ über das Fußballspiel in Yaoundé.
3. Das \_\_\_\_\_ war sehr wichtig für die Stadt.
4. Die neue Zeitung ist gestern \_\_\_\_\_.
5. Viele \_\_\_\_\_ finden den Artikel interessant.
\tcblower
Solution détaillée.

1. \all{Die Überschrift des Artikels ist sehr lang.} « Le titre de l'article est très long. » (féminin, nominatif)

2. \all{Der Journalist berichtet über das Fußballspiel.} « Le journaliste rend compte du match de football. » (verbe régulier, 3\textsuperscript{e} personne : -t)

3. \all{Das Ereignis war sehr wichtig.} « L'événement était très important. » (neutre : das)

4. \all{Die neue Zeitung ist gestern erschienen.} « Le nouveau journal est paru hier. » (parfait avec sein !)

5. \all{Viele Leser finden den Artikel interessant.} « Beaucoup de lecteurs trouvent l'article intéressant. » (pluriel identique au singulier)
\end{exoresolu}

\begin{attention}[Les trois erreurs classiques des francophones]
\begin{enumerate}
\item Oublier la majuscule : on écrit \all{der Artikel}, \all{die Zeitung}, jamais de minuscule à un nom.
\item Dire \all{berichten das Ereignis} sans préposition : il faut toujours \all{über ein Ereignis berichten}.
\item Conjuguer \all{erscheinen} avec \all{haben} : le parfait correct est \all{ist erschienen}.
\end{enumerate}
\end{attention}

\begin{savaistu}[La majuscule, fierté de la langue allemande]
En allemand, \textbf{tous} les noms communs prennent une majuscule, partout dans la phrase : \all{der Artikel}, \all{die Überschrift}, \all{das Ereignis}. C'est une particularité unique parmi les grandes langues européennes : le français et l'anglais ne le font pas. Cette majuscule est une aide précieuse pour le lecteur : elle permet de repérer immédiatement les noms dans une phrase. Au Moyen Âge, on écrivait aussi avec des majuscules en français ; l'allemand est la seule langue qui a conservé cette habitude jusqu'à aujourd'hui.
\end{savaistu}

\begin{exercice}[À vous de jouer]
1. Traduisez en allemand : a) Le journaliste écrit un article sur l'événement. b) Le journal paraît chaque semaine. c) Les lecteurs lisent les gros titres.

2. Formez le féminin et le pluriel : \all{der Leser}, \all{der Journalist}.

3. Conjuguez au Perfekt : \all{lesen} (ich), \all{schreiben} (sie), \all{erscheinen} (die Zeitung).
\end{exercice}

\begin{corrige}[Exercice « À vous de jouer »]
1. a) \all{Der Journalist schreibt einen Artikel über das Ereignis.} b) \all{Die Zeitung erscheint jede Woche.} c) \all{Die Leser lesen die Schlagzeilen.}

2. \all{der Leser} $\rightarrow$ \all{die Leserin}, pluriel \all{die Leserinnen} ; \all{der Journalist} $\rightarrow$ \all{die Journalistin}, pluriel \all{die Journalistinnen}.

3. \all{ich habe gelesen} ; \all{sie hat geschrieben} ; \all{die Zeitung ist erschienen}.
\end{corrige}

Ce vocabulaire est maintenant votre boîte à outils : gardez-le sous les yeux, car nous en aurons besoin dans la partie grammaticale pour construire des phrases au parfait et au prétérit, puis pour rédiger votre propre article de journal.

\coursec{Grammaire 1 : le parfait et le prétérit}

Après avoir enrichi notre vocabulaire de la presse (\cle{der Artikel}, \cle{die Überschrift}, \cle{der Leser}, \cle{berichten}...), nous allons maintenant observer comment le journaliste raconte les événements passés. En allemand, il existe deux temps du passé qui correspondent au passé composé, à l'imparfait et au passé simple français : le \cle{parfait} (das Perfekt) et le \cle{prétérit} (das Präteritum). Cette section est essentielle : sans elle, impossible de comprendre — et encore moins de rédiger — un article de journal.

\courssub{Observation : les verbes du passé dans le texte}

Relisons notre article de journal et relevons les verbes qui expriment le passé. Voici les formes que nous avons rencontrées :

\begin{center}
\begin{tabularx}{\linewidth}{|X|X|X|}\hline
\rowcolor{popblueL} Forme dans le texte & Infinitif & Français \\ \hline
hat stattgefunden & stattfinden & a eu lieu \\ \hline
hat gewonnen & gewinnen & a gagné \\ \hline
kamen & kommen & sont venus / vinrent \\ \hline
berichtete & berichten & a rapporté / rapporta \\ \hline
sagte & sagen & a dit / dit \\ \hline
war & sein & était / fut \\ \hline
\end{tabularx}
\end{center}

Classons ces formes. Nous constatons deux groupes bien distincts :

\begin{itemize}
\item \textbf{Groupe 1 (formes en deux mots)} : \all{hat stattgefunden}, \all{hat gewonnen} — un auxiliaire conjugué (\all{hat}) + un participe à la fin de la phrase. C'est le \textbf{parfait} (das Perfekt).
\item \textbf{Groupe 2 (formes en un seul mot)} : \all{kamen}, \all{berichtete}, \all{sagte}, \all{war} — le verbe seul, conjugué. C'est le \textbf{prétérit} (das Präteritum).
\end{itemize}

\begin{aretenir}[Les deux temps du passé en allemand]
L'allemand possède deux temps principaux pour raconter le passé : le \cle{parfait} (temps composé : auxiliaire + participe passé) et le \cle{prétérit} (temps simple). Le choix entre les deux ne dépend pas du sens, mais du \textbf{registre de langue} : parfait pour la langue parlée, prétérit pour la langue écrite et le récit.
\end{aretenir}

\courssub{Le parfait (das Perfekt) : formation}

\begin{propriete}[Formation du parfait]
Le parfait se forme avec :
\begin{enumerate}
\item l'auxiliaire \cle{haben} ou \cle{sein}, conjugué au présent, en \textbf{position 2} de la phrase ;
\item le \cle{participe passé} (Partizip II) du verbe principal, placé \textbf{à la fin} de la phrase.
\end{enumerate}
\all{Die Mannschaft hat das Finale gewonnen.} « L'équipe a gagné la finale. »
\end{propriete}

Attention à la construction : l'auxiliaire et le participe forment une sorte de « cadre » (der Satzklammer) qui enserre le reste de la phrase. C'est une règle absolue de l'allemand.

\begin{exemplebox}[Le cadre du parfait]
\all{Die Mannschaft hat gestern Abend in Yaoundé das Finale gewonnen.} « L'équipe a gagné la finale hier soir à Yaoundé. »

\all{Das Turnier hat am Samstag stattgefunden.} « Le tournoi a eu lieu samedi. »

\all{Die Schüler haben viele Artikel gelesen.} « Les élèves ont lu beaucoup d'articles. »
\end{exemplebox}

\textbf{Les participes passés}

Il existe trois grands types de participes passés :

\begin{center}
\begin{tabularx}{\linewidth}{|X|X|X|X|}\hline
\rowcolor{popblueL} Type & Formation & Exemple & Français \\ \hline
Verbes faibles (réguliers) & ge + radical + t & machen $\rightarrow$ ge\textbf{mach}\textbf{t} & fait \\ \hline
Verbes forts (irréguliers) & ge + (radical modifié) + en & gehen $\rightarrow$ ge\textbf{gang}\textbf{en} & allé \\ \hline
Verbes en -ieren & radical + t (sans ge-) & trainieren $\rightarrow$ \textbf{trainier}\textbf{t} & entraîné \\ \hline
\end{tabularx}
\end{center}

Autres exemples : \all{sagen $\rightarrow$ gesagt} (dit), \all{berichten $\rightarrow$ berichtet} (rapporté), \all{gewinnen $\rightarrow$ gewonnen} (gagné), \all{kommen $\rightarrow$ gekommen} (venu), \all{lesen $\rightarrow$ gelesen} (lu), \all{schreiben $\rightarrow$ geschrieben} (écrit), \all{studieren $\rightarrow$ studiert} (étudié), \all{organisieren $\rightarrow$ organisiert} (organisé).

\begin{attention}[Les verbes en -ieren et les préfixes inséparables n'ont pas de ge-]
Les verbes en \textbf{-ieren} (souvent d'origine française : \all{trainieren}, \all{studieren}, \all{organisieren}, \all{passieren}) forment leur participe passé \textbf{sans ge-} : \all{trainiert} et non *\all{getrainiert}, \all{organisiert} et non *\all{georganisiert}. Erreur très fréquente chez les francophones ! De même, les verbes à préfixe inséparable (\all{berichten}, \all{verstehen}, \all{bekommen}, \all{enthalten}) n'ont pas de ge- : \all{berichtet}, \all{verstanden}, \all{bekommen}, \all{enthalten}.
\end{attention}

\textbf{Le choix de l'auxiliaire : haben ou sein ?}

\begin{propriete}[Haben ou sein au parfait ?]
\begin{itemize}
\item \textbf{haben} : c'est le cas général (la grande majorité des verbes) : \all{ich habe gesagt}, \all{sie hat gewonnen}.
\item \textbf{sein} : pour les verbes de \textbf{mouvement} (\all{kommen}, \all{gehen}, \all{fahren}, \all{fliegen}, \all{laufen}, \all{reisen}) et de \textbf{changement d'état} (\all{sterben}, \all{aufstehen}, \all{einschlafen}, \all{erscheinen}, \all{aufwachen}), ainsi que pour \all{sein} et \all{bleiben} eux-mêmes : \all{ich bin gekommen} (je suis venu), \all{er ist geblieben} (il est resté).
\end{itemize}
\end{propriete}

\begin{exemplebox}[Sein au parfait]
\all{Viele Zuschauer sind ins Stadion gekommen.} « Beaucoup de spectateurs sont venus au stade. »

\all{Die Mannschaft ist nach Douala gefahren.} « L'équipe est allée à Douala. »

\all{Der Artikel ist gestern erschienen.} « L'article a paru hier. »
\end{exemplebox}

\courssub{Le prétérit (das Präteritum) : formation}

Le prétérit est un temps simple : le verbe se conjugue directement, sans auxiliaire.

\textbf{Verbes faibles (réguliers)} : radical + \textbf{-te} + terminaison : \all{machen $\rightarrow$ ich machte}, \all{sagen $\rightarrow$ er sagte}, \all{berichten $\rightarrow$ sie berichtete}.

\textbf{Verbes forts (irréguliers)} : changement de voyelle, à apprendre par cœur : \all{kommen $\rightarrow$ kam}, \all{lesen $\rightarrow$ las}, \all{schreiben $\rightarrow$ schrieb}, \all{gewinnen $\rightarrow$ gewann}, \all{sein $\rightarrow$ war}, \all{haben $\rightarrow$ hatte}.

\begin{center}
\begin{tabular}{|l|l|l|l|}\hline
\rowcolor{popblueL} Personne & sein (war) & haben (hatte) & sagen (sagte) \\ \hline
ich & war & hatte & sagte \\ \hline
du & warst & hattest & sagtest \\ \hline
er/sie/es & war & hatte & sagte \\ \hline
wir & waren & hatten & sagten \\ \hline
ihr & wart & hattet & sagtet \\ \hline
sie/Sie & waren & hatten & sagten \\ \hline
\end{tabular}
\end{center}

Remarque importante : à la 1\iere{} et à la 3\ieme{} personne du singulier, le prétérit n'a \textbf{pas de terminaison} : \all{ich kam}, \all{er kam} (et non *\all{er kamte}).

\begin{exemplebox}[Le prétérit dans la presse]
\all{Viele Zuschauer kamen ins Stadion.} « Beaucoup de spectateurs sont venus (vinrent) au stade. »

\all{Der Kapitän war stolz.} « Le capitaine était fier. »

\all{Der Reporter berichtete über das Spiel.} « Le journaliste a rapporté (rapporta) les faits du match. »

\all{Die Zeitung schrieb einen langen Artikel.} « Le journal a écrit (écrivit) un long article. »
\end{exemplebox}

\courssub{Parfait ou prétérit ? La question du registre}

En français, nous choisissons entre passé composé, imparfait et passé simple selon l'\textbf{aspect} (action ponctuelle ou durative). En allemand, le choix entre parfait et prétérit dépend surtout du \textbf{registre de langue} :

\begin{center}
\begin{tabularx}{\linewidth}{|X|X|X|}\hline
\rowcolor{popblueL} Français & Deutsch & Remarque \\ \hline
Passé composé (oral) & Perfekt & conversation, dialogue, interview \\ \hline
Passé simple / imparfait (récit écrit) & Präteritum & presse, roman, récit, article \\ \hline
J'étais / j'avais & ich war / ich hatte & sein et haben : toujours au prétérit \\ \hline
\end{tabularx}
\end{center}

\begin{propriete}[Emplois des deux temps]
\begin{itemize}
\item Le \textbf{parfait} : langue parlée, conversation, lettre personnelle, interview citée entre guillemets. \all{„Wir haben gut gespielt“, sagte der Trainer.} « Nous avons bien joué, a dit l'entraîneur. »
\item Le \textbf{prétérit} : langue écrite, article de journal, récit, roman, reportage. \all{Das Spiel begann um 16 Uhr.} « Le match a commencé à 16 heures. »
\item \textbf{Obligatoire au prétérit} (même à l'oral) : \all{sein} (\all{war}), \all{haben} (\all{hatte}) et les verbes modaux (\all{musste}, \all{konnte}, \all{wollte}...).
\end{itemize}
\end{propriete}

\begin{attention}[Sein et haben : toujours au prétérit]
Même dans la conversation, on dit \all{ich war} (j'étais) et \all{ich hatte} (j'avais), presque jamais \all{ich bin gewesen} ou \all{ich habe gehabt}. Dans la presse écrite, \all{ich bin gewesen} est tout simplement impossible. Retenez : \all{Gestern war ich in der Schule.} « Hier, j'étais à l'école. »
\end{attention}

\begin{savaistu}[Pourquoi l'article mélange les deux temps]
Dans un article de journal, le \textbf{prétérit domine}, car c'est la langue écrite du récit. Mais quand le journaliste cite une personne entre guillemets (une interview), on trouve le \textbf{parfait}, car c'est de la langue parlée retranscrite : \all{Der Trainer sagte: „Wir haben fantastisch gespielt!“} Observer ce va-et-vient dans un vrai article est un excellent exercice de lecture.
\end{savaistu}

\courssub{Figures récapitulatives}

\begin{popfigure}
\begin{tikzpicture}[font=\small]
\draw[-{Stealth}, thick, popline] (0,0) -- (12,0) node[right, popdark] {Zeit};
\fill[popblue] (3,0) circle (2.5pt) node[above=2pt, popdark] {Vergangenheit};
\fill[popgreen] (9,0) circle (2.5pt) node[above=2pt, popdark] {Gegenwart};
\draw[-{Stealth}, thick, poporange] (3,0) -- (1.2,1.8) node[above, align=center, popdark] {Perfekt\\ oral, interview};
\draw[-{Stealth}, thick, poppurple] (3,0) -- (4.8,1.8) node[above, align=center, popdark] {Präteritum\\ écrit, presse};
\end{tikzpicture}
\legende{Le passé allemand : deux branches selon le registre de langue}
\end{popfigure}

\begin{popfigure}
\begin{tikzpicture}[font=\small]
\node[draw=popblue, fill=popblueL, rounded corners, align=center, minimum width=2.6cm] (v1) at (0,0) {kommen};
\node[draw=poporange, fill=poporangeL, rounded corners, align=center, minimum width=2.6cm] (v2) at (4.5,0) {kam};
\node[draw=popgreen, fill=popgreenL, rounded corners, align=center, minimum width=3.2cm] (v3) at (9.5,0) {ist gekommen};
\node[draw=popblue, fill=popblueL, rounded corners, align=center, minimum width=2.6cm] (w1) at (0,-1.6) {gewinnen};
\node[draw=poporange, fill=poporangeL, rounded corners, align=center, minimum width=2.6cm] (w2) at (4.5,-1.6) {gewann};
\node[draw=popgreen, fill=popgreenL, rounded corners, align=center, minimum width=3.2cm] (w3) at (9.5,-1.6) {hat gewonnen};
\draw[-{Stealth}, thick, popdark] (v1) -- (v2);
\draw[-{Stealth}, thick, popdark] (v2) -- (v3);
\draw[-{Stealth}, thick, popdark] (w1) -- (w2);
\draw[-{Stealth}, thick, popdark] (w2) -- (w3);
\node[above=3pt of v1, popdark] {Infinitiv};
\node[above=3pt of v2, popdark] {Präteritum};
\node[above=3pt of v3, popdark] {Partizip II};
\end{tikzpicture}
\legende{De l'infinitif au prétéritum et au participe II}
\end{popfigure}

\textbf{Dix verbes irréguliers fréquents dans la presse}

\begin{center}
\begin{tabular}{|l|l|l|l|}\hline
\rowcolor{popblueL} Infinitiv & Präteritum & Partizip II & Français \\ \hline
kommen & kam & ist gekommen & venir \\ \hline
gewinnen & gewann & hat gewonnen & gagner \\ \hline
lesen & las & hat gelesen & lire \\ \hline
schreiben & schrieb & hat geschrieben & écrire \\ \hline
finden & fand & hat gefunden & trouver \\ \hline
geben & gab & hat gegeben & donner \\ \hline
sehen & sah & hat gesehen & voir \\ \hline
sprechen & sprach & hat gesprochen & parler \\ \hline
beginnen & begann & hat begonnen & commencer \\ \hline
stehen & stand & hat gestanden & se trouver (debout) \\ \hline
\end{tabular}
\end{center}

\courssub{Texte d'application : un court article au passé}

\begin{textebox}[„Schulfest in Yaoundé“ — Schulzeitung]
Am Samstag fand das große Schulfest unseres Gymnasiums statt. Schon um 8 Uhr kamen viele Schüler und ihre Eltern auf den Schulhof. Der Direktor eröffnete das Fest mit einer kurzen Rede und dankte allen Lehrern für ihre Arbeit.

Die Fußballmannschaft spielte gegen ein Team aus Douala und gewann 3:1. Der Kapitän war nach dem Spiel sehr stolz. „Wir haben zwei Monate hart trainiert“, sagte er unserem Reporter. „Der Sieg ist das Ergebnis unserer Arbeit.“

Am Nachmittag gab es auch Musik und Theater. Eine Gruppe aus der Klasse 11 sang deutsche Lieder, und das Publikum applaudierte lange. Viele Besucher lasen auch die Artikel unserer Schülerzeitung, die an einer Wand hingen.

Am Abend beendete der Direktor das Fest. Alle waren müde, aber glücklich. Nächstes Jahr wollen wir das Fest wieder organisieren!
\end{textebox}

\textbf{Glossaire} : das Schulfest, -e (la fête de l'école) ; der Schulhof, -höfe (la cour de l'école) ; die Rede, -n (le discours) ; danken (remercier) ; der Sieg, -e (la victoire) ; das Ergebnis, -se (le résultat) ; das Publikum (le public) ; applaudieren (applaudir) ; die Wand, -wände (le mur) ; müde (fatigué).

\textbf{Questions de compréhension} : 1. Wann fand das Schulfest statt ? 2. Wer eröffnete das Fest ? 3. Wie viele Tore schoss die Mannschaft ? 4. Was sagte der Kapitän ? 5. Warum waren alle am Abend müde, aber glücklich ?

\courssub{Exercice résolu}

\begin{exoresolu}[Parfait ou prétérit ?]
Complétez avec le verbe entre parenthèses au temps qui convient (contexte : article de journal, sauf les paroles citées).

1. Das Konzert (beginnen) \dots\ um 19 Uhr. \quad 2. Der Sänger (sein) \dots\ sehr gut. \quad 3. Viele Fans (kommen) \dots\ aus ganz Kamerun. \quad 4. „Wir (haben, Perfekt) \dots\ eine tolle Nacht“, sagte ein Fan. \quad 5. Die Zeitung (schreiben) \dots\ einen langen Artikel.
\tcblower
\textbf{Solution détaillée :}

1. \all{begann} — prétérit du verbe fort \all{beginnen} (récit écrit). « Le concert a commencé à 19 heures. »

2. \all{war} — \all{sein} s'emploie toujours au prétérit : « Le chanteur était très bon. »

3. \all{kamen} — prétérit du verbe fort \all{kommen} (langue écrite de la presse). « Beaucoup de fans sont venus de tout le Cameroun. »

4. \all{Wir haben eine tolle Nacht gehabt.} — ici, paroles citées entre guillemets = langue parlée, donc parfait. (Remarque : à l'oral, on préfère souvent \all{„Es war eine tolle Nacht!“})

5. \all{schrieb} — prétérit du verbe fort \all{schreiben}. « Le journal a écrit un long article. »
\end{exoresolu}

\begin{attention}[Pièges fréquents des francophones]
\begin{itemize}
\item Ne placez jamais le participe juste après l'auxiliaire comme en français : *\all{Die Mannschaft hat gewonnen das Finale} est faux. Le participe va \textbf{à la fin} : \all{Die Mannschaft hat das Finale gewonnen.}
\item N'oubliez pas l'alternance de voyelle des verbes forts : \all{kommen $\rightarrow$ kam} (et non *\all{kommte}).
\item Ne traduisez pas « j'étais » par *\all{ich bin gewesen} : dites \all{ich war}.
\item Majuscule aux noms, même au milieu de la phrase : \all{das Finale}, \all{die Zeitung}.
\end{itemize}
\end{attention}

\begin{exercice}[À vous de jouer]
Mettez les phrases suivantes au temps demandé, puis traduisez-les en français.

1. (Perfekt) Die Schüler / lesen / der Artikel. \quad 2. (Präteritum) Der Reporter / berichten / über das Spiel. \quad 3. (Präteritum) Viele Zuschauer / kommen / ins Stadion. \quad 4. (Präteritum) Der Kapitän / sein / stolz. \quad 5. (Perfekt) Das Turnier / stattfinden / am Samstag.
\end{exercice}

\begin{corrige}[Exercice « À vous de jouer »]
1. \all{Die Schüler haben den Artikel gelesen.} « Les élèves ont lu l'article. » (auxiliaire \all{haben} en position 2, participe \all{gelesen} à la fin ; \all{der Artikel} $\rightarrow$ accusatif \all{den Artikel}.)

2. \all{Der Reporter berichtete über das Spiel.} « Le journaliste a rapporté les faits du match. » (verbe faible : radical + -te.)

3. \all{Viele Zuschauer kamen ins Stadion.} « Beaucoup de spectateurs sont venus au stade. » (verbe fort : \all{kommen $\rightarrow$ kamen}.)

4. \all{Der Kapitän war stolz.} « Le capitaine était fier. » (\all{sein} : toujours au prétérit.)

5. \all{Das Turnier hat am Samstag stattgefunden.} « Le tournoi a eu lieu samedi. » (verbe à particule séparable : \all{stattgefunden} en un seul mot à la fin ; auxiliaire \all{sein} car changement de lieu/état.)
\end{corrige}

\begin{aretenir}[L'essentiel de la section]
\begin{itemize}
\item \textbf{Parfait} = haben/sein conjugué (position 2) + participe passé (fin de phrase) ; langue parlée et interviews.
\item \textbf{Prétérit} = verbe conjugué seul ; langue écrite, presse, récit ; obligatoire pour \all{sein} (\all{war}), \all{haben} (\all{hatte}) et les modaux.
\item Participes : \all{gemacht} (faibles), \all{gegangen} (forts), \all{trainiert} (-ieren, sans ge-).
\item Auxiliaire \all{sein} pour les verbes de mouvement et de changement d'état.
\end{itemize}
\end{aretenir}

\coursec{Grammaire 2 : les subordonnées}

Dans la section précédente, nous avons appris à raconter le passé avec le \cle{Perfekt} (\all{die Mannschaft hat gewonnen}) et le \cle{Präteritum} (\all{das Wetter war schön}). Mais un article de journal ne se contente pas d'empiler des phrases courtes : le journaliste relie les idées entre elles. Il explique \emph{pourquoi} les spectateurs sont venus, il rapporte \emph{ce que} l'entraîneur a déclaré. Pour cela, l'allemand utilise des \cle{propositions subordonnées} (\all{die Nebensätze}). C'est l'une des règles les plus importantes — et les plus difficiles pour les francophones — de toute la grammaire allemande : dans une subordonnée, le verbe conjugué change de place.

\courssub{Observation : où est le verbe conjugué ?}

Reprenons deux phrases tirées de notre article de journal et observons-les attentivement.

\begin{exemplebox}[Deux phrases du texte]
\all{Viele Zuschauer kamen, weil das Wetter schön war.}\\
« Beaucoup de spectateurs sont venus parce qu'il faisait beau. »

\medskip
\all{Der Trainer berichtete, dass die Jungen jeden Tag trainiert hatten.}\\
« L'entraîneur a rapporté que les garçons s'étaient entraînés chaque jour. »
\end{exemplebox}

Analysons la place du verbe conjugué :

\begin{itemize}
\item Dans \all{weil das Wetter schön \textbf{war}}, le verbe conjugué \all{war} n'est pas à la deuxième position : il est \textbf{tout à la fin} de la proposition.
\item Dans \all{dass die Jungen jeden Tag trainiert \textbf{hatten}}, l'auxiliaire conjugué \all{hatten} est également \textbf{à la fin}, \emph{après} le participe passé \all{trainiert}.
\item En français, rien ne bouge : « parce que le temps \textbf{était} beau », « que les garçons s'\textbf{étaient} entraînés ». Le verbe reste près du sujet.
\item Remarquez aussi la \textbf{virgule} : en allemand, elle est \textbf{obligatoire} entre la principale et la subordonnée, même quand le français n'en met pas : \all{Er sagt, dass...} (« il dit que... »).
\end{itemize}

\begin{popfigure}
\begin{tikzpicture}[
  box/.style={draw=popblue, thick, rounded corners=2pt, fill=popblueL, align=center, font=\small, inner sep=4pt},
  vbox/.style={draw=poporange, thick, rounded corners=2pt, fill=poporangeL, align=center, font=\small, inner sep=4pt},
  lbl/.style={font=\small\bfseries, text=popdark}
]
\node[lbl] at (-4.5,2.2) {Hauptsatz (principale) : verbe en position 2};
\node[box, align=center] (s1) at (-7.2,1.4){Viele Zuschauer\\ (Sujet)};
\node[vbox, align=center] (v1) at (-4.5,1.4){\textbf{kamen}\\ (Verbe : pos. 2)};
\node[box, align=center] (r1) at (-1.6,1.4){gestern.\\ (Reste)};

\node[lbl] at (-4.5,-0.2) {Nebensatz (subordonnée) : verbe à la fin};
\node[box, align=center] (c2) at (-7.2,-1.0){weil\\ (Conjonction)};
\node[box, align=center] (s2) at (-5.2,-1.0){das Wetter\\ (Sujet)};
\node[box, align=center] (r2) at (-3.0,-1.0){schön\\ (Reste)};
\node[vbox, align=center] (v2) at (-1.0,-1.0){\textbf{war.}\\ (Verbe : fin)};
\draw[-{Stealth}, thick, poporange] (v2.south) to[bend right=25] node[below, font=\scriptsize, text=poporange]{position finale} (c2.south);
\end{tikzpicture}
\legende{Principale : le verbe conjugué occupe la 2\textsuperscript{e} position. Subordonnée : la conjonction « pousse » le verbe conjugué à la toute dernière place.}
\end{popfigure}

\courssub{La règle d'or}

\begin{propriete}[Le verbe conjugué à la fin de la subordonnée]
Après une conjonction de subordination comme \all{dass} (que), \all{weil} (parce que), \all{als} (quand, passé), \all{wenn} (quand, si), \all{ob} (si), le \textbf{verbe conjugué se place en dernière position} de la proposition subordonnée :

\medskip
\all{..., weil das Wetter schön \textbf{war}.}\\
\all{..., dass die Jungen trainiert \textbf{hatten}.}

\medskip
Structure : \textbf{conjonction + sujet + compléments + verbe conjugué (fin)}.

\medskip
Cas particulier : si la subordonnée contient \textbf{deux verbes} (par exemple le parfait : auxiliaire + participe), c'est l'\textbf{auxiliaire conjugué} qui va à la fin, \emph{après} le participe passé :\\
\all{..., dass er das Spiel gewonnen \textbf{hat}.} « ... qu'il a gagné le match. »\\
\all{..., weil sie lange geübt \textbf{hatten}.} « ... parce qu'elles s'étaient longuement entraînées. »
\end{propriete}

\begin{attention}[L'erreur n°1 des francophones]
En français, le verbe ne bouge jamais : « parce que le temps \emph{était} beau ». Le francophone a donc tendance à calquer cette structure et à dire :

\medskip
\all{*weil das Wetter war schön} — \textbf{FAUX !}

\medskip
La seule formulation correcte est : \all{weil das Wetter schön \textbf{war}}. Dès que vous écrivez \all{weil}, \all{dass}, \all{als}, \all{wenn} ou \all{ob}, envoyez mentalement le verbe conjugué à la dernière place. C'est le réflexe le plus important à acquérir en Première.
\end{attention}

\begin{attention}[dass ou das ? Ne confondez pas l'orthographe]
\all{dass} (avec \textbf{ss}) est la conjonction « que » ; \all{das} (avec un seul \textbf{s}) est l'article ou le pronom « le, ce ». Test : si vous pouvez remplacer le mot par \all{dieses} (« ce »), écrivez \all{das} ; sinon, écrivez \all{dass}.

\medskip
\all{Er sagt, \textbf{dass} das Wetter schön ist.} « Il dit que le temps est beau. »\\
\all{\textbf{Das} Wetter ist schön.} « Le temps est beau. » (on peut dire : \all{dieses Wetter})
\end{attention}

\courssub{Les conjonctions de subordination du programme}

\begin{definition}[Les cinq conjonctions à connaître]
\begin{itemize}
\item \all{dass} = « que » : introduit un discours rapporté ou une opinion. \all{Der Trainer sagt, dass die Jungen gut gespielt haben.} « L'entraîneur dit que les garçons ont bien joué. »
\item \all{weil} = « parce que » : exprime la cause. \all{Wir blieben zu Hause, weil es regnete.} « Nous sommes restés à la maison parce qu'il pleuvait. »
\item \all{als} = « quand, lorsque » : un événement \textbf{unique au passé}. \all{Als das Spiel endete, jubelten alle.} « Quand le match s'est terminé, tous ont jubilé. »
\item \all{wenn} = « quand, si » : répétition, condition, présent ou futur. \all{Wenn es regnet, bleibe ich zu Hause.} « Quand il pleut, je reste à la maison. »
\item \all{ob} = « si » (interrogation indirecte). \all{Er fragte, ob die Mannschaft gewonnen hatte.} « Il a demandé si l'équipe avait gagné. »
\end{itemize}
\end{definition}

\begin{popfigure}
\begin{tikzpicture}[
  conj/.style={draw=popblue, thick, rounded corners=2pt, fill=popblueL, font=\small\bfseries, align=center, inner sep=4pt},
  sens/.style={draw=popgreen, thick, rounded corners=2pt, fill=popgreenL, font=\small, align=center, inner sep=4pt},
  ex/.style={draw=poppurple, thick, rounded corners=2pt, fill=poppurpleL, font=\scriptsize, align=center, inner sep=4pt}
]
\node[conj] (c1) at (0,0) {dass};
\node[sens, right=0.4cm of c1] (s1) {que};
\node[ex, right=0.4cm of s1, align=center] (e1){Er sagt, dass sie gewonnen hat.\\ « Il dit qu'elle a gagné. »};

\node[conj] (c2) at (0,-1.3) {weil};
\node[sens, right=0.4cm of c2] (s2) {parce que};
\node[ex, right=0.4cm of s2, align=center] (e2){..., weil das Wetter schön war.\\ « ... parce qu'il faisait beau. »};

\node[conj] (c3) at (0,-2.6) {als};
\node[sens, right=0.4cm of c3, align=center] (s3){quand\\ (passé, 1 fois)};
\node[ex, right=0.4cm of s3, align=center] (e3){Als das Spiel endete, ...\\ « Quand le match s'est terminé, ... »};

\node[conj] (c4) at (0,-3.9) {wenn};
\node[sens, right=0.4cm of c4, align=center] (s4){quand, si\\ (répétition, condition)};
\node[ex, right=0.4cm of s4, align=center] (e4){Wenn es regnet, bleibe ich hier.\\ « Quand il pleut, je reste ici. »};

\node[conj] (c5) at (0,-5.2) {ob};
\node[sens, right=0.4cm of c5, align=center] (s5){si\\ (question indirecte)};
\node[ex, right=0.4cm of s5, align=center] (e5){Er fragte, ob sie gewonnen hatte.\\ « Il a demandé si elle avait gagné. »};
\end{tikzpicture}
\legende{Les cinq conjonctions de subordination : sens et exemples traduits.}
\end{popfigure}

\courssub{als ou wenn ? La distinction capitale}

Le français n'a qu'un seul mot, « quand » ; l'allemand en a deux, et le choix n'est pas libre.

\begin{propriete}[als = passé unique ; wenn = répétition ou condition]
\begin{itemize}
\item \all{als} s'emploie pour un événement qui s'est produit \textbf{une seule fois, au passé} :\\
\all{Als ich ein Kind war, spielte ich jeden Tag Fußball.} « Quand j'étais enfant, je jouais au football chaque jour. » (une période unique du passé)\\
\all{Als der Schiedsrichter pfiff, wurde es still.} « Quand l'arbitre a sifflé, il est devenu silencieux. »
\item \all{wenn} s'emploie pour la \textbf{répétition} (chaque fois que), la \textbf{condition} (si), le \textbf{présent} et le \textbf{futur} :\\
\all{Wenn es regnet, bleibe ich zu Hause.} « Quand il pleut, je reste à la maison. » (à chaque fois)\\
\all{Wenn ich Zeit habe, lese ich die Zeitung.} « Si j'ai le temps, je lis le journal. »
\end{itemize}
\end{propriete}

\begin{attention}[Ne confondez pas als et wenn]
\all{Als ich ein Kind war...} (une seule période passée) mais \all{Wenn es regnet, bleibe ich zu Hause.} (répétition ou condition). Dire \all{*als es regnet} pour une habitude ou une condition est une erreur classique. Question-test : « est-ce arrivé une seule fois dans le passé ? » → oui : \all{als} ; non : \all{wenn}.
\end{attention}

\begin{center}
\begin{tabularx}{\linewidth}{|X|X|X|}
\hline
\rowcolor{popblueL} Français & Deutsch & Remarque \\
\hline
Quand j'étais enfant, ... & Als ich ein Kind war, ... & passé, une seule fois : \all{als} \\
\hline
Quand il pleut, je reste ici. & Wenn es regnet, bleibe ich hier. & répétition : \all{wenn} \\
\hline
Si j'ai le temps, je viens. & Wenn ich Zeit habe, komme ich. & condition : \all{wenn} \\
\hline
Il dit qu'il fait beau. & Er sagt, dass das Wetter schön ist. & verbe à la fin en allemand \\
\hline
Il a demandé si elle venait. & Er fragte, ob sie kam. & « si » interrogatif = \all{ob} \\
\hline
\end{tabularx}
\end{center}

\courssub{La subordonnée en tête de phrase}

Pour varier le style, le journaliste place souvent la subordonnée au début de la phrase. Une règle mécanique s'applique alors.

\begin{propriete}[Subordonnée en tête : le verbe de la principale suit la virgule]
Quand la subordonnée occupe la première position de la phrase, elle compte comme la \textbf{première position} ; le verbe conjugué de la principale vient donc \textbf{immédiatement après la virgule} (position 2 de la phrase, position 1 de la principale), et le sujet de la principale se place après ce verbe :

\medskip
\all{\textbf{Weil das Wetter schön war}, \textbf{kamen} viele Zuschauer.}\\
« Parce qu'il faisait beau, beaucoup de spectateurs sont venus. »

\medskip
\all{\textbf{Als das Spiel endete}, \textbf{jubelten} alle.}\\
« Quand le match s'est terminé, tous ont jubilé. »

\medskip
Schéma : [Subordonnée ... verbe final] , \textbf{verbe} + sujet + reste. On obtient deux verbes conjugués « dos à dos » de part et d'autre de la virgule : \all{..., war, kamen ...}.
\end{propriete}

\begin{exemplebox}[Comparez les deux constructions]
\all{Viele Zuschauer kamen, weil das Wetter schön war.} (principale d'abord : verbe en position 2)\\
\all{Weil das Wetter schön war, kamen viele Zuschauer.} (subordonnée d'abord : le verbe \all{kamen} colle à la virgule)

\medskip
Les deux phrases signifient exactement la même chose ; seul le rythme change. Attention : dans le deuxième cas, on ne dit jamais \all{*Weil das Wetter schön war, viele Zuschauer kamen}.
\end{exemplebox}

\courssub{L'interrogation indirecte avec ob}

Pour rapporter une question fermée (réponse oui/non), le français utilise « si » ; l'allemand utilise \all{ob}, avec le verbe conjugué à la fin.

\begin{exemplebox}[Question directe et question indirecte]
Question directe : \all{Hat die Mannschaft gewonnen?} « L'équipe a-t-elle gagné ? »

\medskip
Question indirecte : \all{Er fragte, ob die Mannschaft gewonnen hatte.} « Il a demandé si l'équipe avait gagné. »

\medskip
Remarquez : le journaliste rapporte souvent au \textbf{plus-que-parfait} (\all{hatte gewonnen}) ce qui s'est passé avant sa déclaration.
\end{exemplebox}

\courssub{Application au texte : un court article}

\begin{textebox}[Ein Bericht aus Yaoundé]
\all{Am Samstag fand in Yaoundé ein großes Fußballspiel statt. Viele Zuschauer kamen früh, weil das Wetter schön war. Ein Journalist berichtete, dass die Jungen aus Douala jeden Tag trainiert hatten. Als das Spiel begann, wurde es sehr laut im Stadion. Der Trainer sagte nach dem Spiel, dass seine Mannschaft sehr gut gespielt hatte. Er erklärte auch, dass er stolz auf die Spieler war. Ein Reporter fragte, ob das Team jetzt der beste Klub der Stadt war. Der Trainer antwortete, dass die Jungen noch viel lernen mussten. Als der Schiedsrichter das Ende pfiff, jubelten alle Fans. Wenn das Wetter nächste Woche wieder schön ist, kommen sicher noch mehr Zuschauer. Viele Leute in Kamerun lieben Fußball, weil der Sport die Menschen zusammenbringt.}

\medskip
\textbf{Glossaire :} das Stadion, die Stadien (le stade) — der Schiedsrichter, - (l'arbitre) — pfeifen, pfiff, hat gepfiffen (siffler) — stolz (fier) — zusammenbringen (rassembler) — der Klub, -s (le club) — stattfinden, fand statt, hat stattgefunden (avoir lieu).

\medskip
\textbf{Questions :} 1. Warum kamen viele Zuschauer früh? 2. Was berichtete der Journalist? 3. Was fragte der Reporter?
\end{textebox}

\courssub{Méthode de vérification et exercice résolu}

\begin{methode}[Vérifier une subordonnée en trois étapes]
\begin{enumerate}
\item \textbf{Soulignez la conjonction} (\all{dass}, \all{weil}, \all{als}, \all{wenn}, \all{ob}) et vérifiez la \textbf{virgule} avant elle.
\item \textbf{Repérez le verbe conjugué} de la proposition qui suit : est-il le \textbf{dernier mot} de cette proposition ? (S'il y a un participe, l'auxiliaire vient après lui.)
\item \textbf{Si la subordonnée est en tête de phrase}, vérifiez que le verbe de la principale suit immédiatement la virgule : \all{..., war, kamen ...}.
\end{enumerate}
\end{methode}

\begin{exoresolu}[Remettez les mots dans l'ordre]
Énoncé : construisez une phrase correcte avec : \all{weil / das Wetter / war / schön / viele Zuschauer / kamen}.

\tcblower
Solution : la conjonction \all{weil} introduit une subordonnée, donc le verbe conjugué \all{war} va à la fin : \all{weil das Wetter schön war}. La principale garde son verbe en position 2 : \all{Viele Zuschauer kamen}. Phrase complète : \all{Viele Zuschauer kamen, weil das Wetter schön war.} « Beaucoup de spectateurs sont venus parce qu'il faisait beau. » Variante avec la subordonnée en tête : \all{Weil das Wetter schön war, kamen viele Zuschauer.}
\end{exoresolu}

\begin{exercice}[À vous de jouer]
Traduisez en allemand : 1. « L'entraîneur dit que les garçons ont bien joué. » 2. « Quand le match s'est terminé, tous ont jubilé. » 3. « Il a demandé si l'équipe avait gagné. » 4. « Parce qu'il pleuvait, nous sommes restés à la maison. » (subordonnée en tête)
\end{exercice}

\begin{corrige}[Exercice : traduction]
1. \all{Der Trainer sagt, dass die Jungen gut gespielt haben.} (parfait dans la subordonnée : \all{haben} après \all{gespielt}). 2. \all{Als das Spiel endete, jubelten alle.} (passé unique : \all{als} ; verbe de la principale après la virgule). 3. \all{Er fragte, ob die Mannschaft gewonnen hatte.} (interrogation indirecte : \all{ob} + plus-que-parfait). 4. \all{Weil es regnete, blieben wir zu Hause.} (subordonnée en tête : \all{blieben} colle à la virgule).
\end{corrige}

\begin{aretenir}[L'essentiel des subordonnées]
\begin{itemize}
\item Après \all{dass}, \all{weil}, \all{als}, \all{wenn}, \all{ob} : le verbe conjugué va \textbf{à la fin}, et la virgule est obligatoire.
\item Avec deux verbes : l'auxiliaire passe \textbf{après} le participe : \all{..., dass er gewonnen hat.}
\item \all{als} = une fois au passé ; \all{wenn} = répétition, condition, présent, futur.
\item Subordonnée en tête : le verbe de la principale suit la virgule : \all{Weil das Wetter schön war, kamen viele Zuschauer.}
\item Orthographe : conjonction \all{dass} (ss) contre article \all{das} (s).
\item Piège majeur : jamais \all{*weil das Wetter war schön} — le français ne déplace pas le verbe, l'allemand si !
\end{itemize}
\end{aretenir}

\coursec{Méthode du commentaire et commentaire modèle}

Dans la section précédente, nous avons appris à construire des subordonnées avec \all{weil}, \all{dass} ou \all{als} : \all{Der Journalist benutzt das Präteritum, weil er eine Geschichte erzählt.} Ces subordonnées vont maintenant nous être très utiles, car un bon commentaire de texte est précisément fait de phrases qui \emph{expliquent} et \emph{justifient} : chaque remarque s'appuie sur une citation et sur une raison introduite par \all{weil} ou \all{dass}. Passons donc de la grammaire à la méthode : comment lire un article de journal comme un germaniste ?

\courssub{Qu'est-ce qu'un commentaire de texte journalistique ?}

\begin{definition}[Le commentaire de texte]
Un \cle{commentaire de texte} (allemand : \all{die Textanalyse}, \all{der Kommentar}) est un exercice écrit qui consiste à \textbf{présenter} un texte (de quoi s'agit-il ?), à \textbf{analyser} sa construction et sa langue (comment est-il écrit ?) et à donner une \textbf{appréciation} personnelle argumentée (qu'en pensez-vous ?). Au lycée, le commentaire se rédige en allemand, en 8 à 12 lignes, avec des phrases simples mais correctes.
\end{definition}

Un article de journal se prête très bien à cet exercice, car sa structure est codifiée : un titre (\all{die Überschrift}), un chapeau (\all{der Vorspann}), un corps d'article (\all{der Artikeltext}) qui répond aux fameuses \cle{5 W} : \all{Wer? Was? Wann? Wo? Warum?} (Qui ? Quoi ? Quand ? Où ? Pourquoi ?).

\begin{attention}[Commentaire n'est pas résumé]
Erreur classique des élèves : recopier le contenu de l'article. Or le \textbf{résumé} dit \emph{ce que} le texte dit ; le \textbf{commentaire} explique \emph{comment} il le dit et \emph{pourquoi} c'est réussi. Dire \all{Die Mannschaft hat gewonnen} est un résumé. Dire \all{Der Journalist benutzt das Präteritum, weil das die Sprache der Presse ist} est un commentaire.
\end{attention}

\courssub{Les quatre étapes du commentaire}

\begin{popfigure}
\begin{tikzpicture}[
  etape/.style={rectangle, rounded corners=4pt, draw=popblue, fill=popblueL, text width=3.6cm, align=center, minimum height=1.1cm, font=\small},
  fleche/.style={-{Stealth[length=3mm]}, thick, poporange}
]
\node[etape, align=center] (e1) at (0,0){\textbf{1. Présentation}\\ type, source, thème, 5 W};
\node[etape, fill=popgreenL, draw=popgreen, align=center] (e2) at (5,0){\textbf{2. Structure}\\ titre, chapeau, corps, ton};
\node[etape, fill=poppurpleL, draw=poppurple, align=center] (e3) at (0,-2.2){\textbf{3. Langue}\\ temps, lexique, subordonnées};
\node[etape, fill=popgoldL, draw=popgold, align=center] (e4) at (5,-2.2){\textbf{4. Appréciation}\\ avis personnel argumenté};
\draw[fleche] (e1) -- (e2);
\draw[fleche] (e2) -- (e3);
\draw[fleche] (e3) -- (e4);
\end{tikzpicture}
\legende{Organigramme : les quatre étapes du commentaire de texte journalistique.}
\end{popfigure}

\textbf{Étape 1 — La présentation.}\par
On commence par identifier le texte : c'est un \all{Zeitungsartikel} (article de journal), paru dans un journal (\all{die Zeitung}), sur un thème précis (\all{das Thema}). On répond ensuite aux 5 W à partir du chapeau de l'article. Phrases utiles :

\all{Der Text ist ein Zeitungsartikel über ein Fußballturnier in Yaoundé.} « Le texte est un article de journal sur un tournoi de football à Yaoundé. »

\all{Das Ereignis findet an einem Gymnasium statt.} « L'événement a lieu dans un lycée. »

\textbf{Étape 2 — L'analyse de la structure.}\par
On observe comment l'article est construit : le titre (\all{die Überschrift}) annonce le sujet et attire le lecteur (\all{der Leser}) ; le chapeau résume l'essentiel ; le corps développe les détails et cite des témoins. On précise le \cle{ton} : un article de presse est en général \all{objektiv} (objectif) et \all{informativ} (informatif), contrairement à un éditorial qui donne une opinion.

\textbf{Étape 3 — L'analyse linguistique.}\par
C'est ici que nos deux leçons de grammaire servent :
\begin{itemize}
\item les \textbf{temps} : le journaliste raconte au \cle{prétérit} (\all{kam}, \all{sagte}, \all{war}), la langue typique de la presse écrite ; le parfait apparaît dans les paroles rapportées des interviews ;
\item le \textbf{vocabulaire} : on relève le champ lexical du thème (ici : \all{der Sport}, \all{die Mannschaft}, \all{das Turnier}, \all{gewinnen}) ;
\item les \textbf{subordonnées} : elles montrent que le journaliste explique et relie les faits (\all{weil}, \all{dass}, \all{als}).
\end{itemize}

\textbf{Étape 4 — L'appréciation personnelle.}\par
On termine par un avis justifié, en lien avec son propre contexte : le sport scolaire au Cameroun, la lecture des journaux à Yaoundé ou à Douala. Formules utiles : \all{Ich finde den Artikel interessant, weil \ldots} / \all{Der Artikel ist wichtig für uns, weil \ldots}

\begin{methode}[Un bon commentaire cite toujours le texte]
\begin{enumerate}
\item Choisis une remarque (temps, vocabulaire, structure).
\item Justifie-la par une \textbf{citation courte} entre guillemets allemands : \all{Der Journalist schreibt: „Die Mannschaft kam um acht Uhr.“}
\item Explique l'effet produit avec une subordonnée : \all{Er benutzt das Präteritum, weil er eine Geschichte erzählt.}
\item Termine par ton avis personnel : \all{Ich finde \ldots}
\end{enumerate}
Règle d'or : pas de remarque sans citation, pas de citation sans explication.
\end{methode}

\courssub{Tableau des formulations utiles}

\begin{center}
\begin{tabularx}{\linewidth}{|X|X|X|}
\hline
\rowcolor{popblueL} Étape & Deutsch & Français \\
\hline
Présentation & Der Artikel berichtet über ein Ereignis in Yaoundé. & L'article rend compte d'un événement à Yaoundé. \\
\hline
Présentation & Die Überschrift informiert sofort den Leser. & Le titre informe immédiatement le lecteur. \\
\hline
Structure & Der Text hat einen Titel und drei Absätze. & Le texte a un titre et trois paragraphes. \\
\hline
Langue & Der Journalist benutzt das Präteritum, weil er eine Geschichte erzählt. & Le journaliste utilise le prétérit parce qu'il raconte. \\
\hline
Langue & Es gibt viele Wörter aus dem Sport. & Il y a beaucoup de mots du domaine du sport. \\
\hline
Appréciation & Ich finde den Artikel interessant, weil Sport an unseren Schulen wichtig ist. & Je trouve l'article intéressant parce que le sport est important dans nos écoles. \\
\hline
\end{tabularx}
\end{center}

\courssub{Commentaire modèle rédigé}

Voici un commentaire complet (environ 10 lignes) sur l'article support de la leçon. Observez les passages clés : la présentation, la citation, l'analyse des temps et l'appréciation.

\begin{textebox}[Commentaire modèle — die Textanalyse]
Der Text ist ein \textbf{Zeitungsartikel} über ein Fußballturnier in Yaoundé. Der Journalist \textbf{berichtet}, dass die Mannschaft des Lycée de Nkol-Eton gewonnen hat. Die \textbf{Überschrift} informiert sofort den Leser über das Ereignis. Der Artikel hat einen Titel, einen Vorspann und drei Absätze: das ist die typische Struktur der Presse. Der Journalist benutzt das \textbf{Präteritum} — „kam“, „sagte“, „war“ —, weil das die Sprache der Presse ist. In den Interviews findet man das Perfekt: „Wir haben gut gespielt!“ Es gibt auch \textbf{Nebensätze} mit „weil“ und „als“; sie erklären die Fakten. Der Wortschatz kommt aus dem Sport: „die Mannschaft“, „das Turnier“, „der Sieg“. Ich finde den Artikel interessant, weil Sport auch an unseren Schulen in Kamerun sehr wichtig ist.
\end{textebox}

Glossaire : \all{der Vorspann, -e} (le chapeau) ; \all{der Absatz, -ä-e} (le paragraphe) ; \all{die Struktur, -en} (la structure) ; \all{der Nebensatz, -"a-tze} (la subordonnée) ; \all{der Wortschatz} (le vocabulaire) ; \all{die Tatsache / das Faktum} — ici \all{die Fakten} (les faits).

Questions de vérification :
\begin{enumerate}
\item Quelle phrase du modèle présente le texte ?
\item Quelle citation prouve l'emploi du prétérit ?
\item Où se trouve l'appréciation personnelle ?
\end{enumerate}

\begin{corrige}[Questions de vérification]
1. La première phrase : \all{Der Text ist ein Zeitungsartikel über ein Fußballturnier in Yaoundé.} 2. La citation : \all{„kam“, „sagte“, „war“}. 3. La dernière phrase : \all{Ich finde den Artikel interessant, weil \ldots}
\end{corrige}

\courssub{Exercice résolu et pièges à éviter}

\begin{exoresolu}[Transformer un résumé en commentaire]
Énoncé — Voici une phrase d'élève : \all{Die Mannschaft hat das Spiel gewonnen.} Est-ce un commentaire ? Sinon, transformez-la en remarque de commentaire.\tcblower
Solution — Non, c'est un résumé : la phrase répète \emph{ce que} dit l'article. Un commentaire explique \emph{comment} le texte est écrit. Correction possible : \all{Der Journalist berichtet, dass die Mannschaft das Spiel gewonnen hat; er benutzt einen Nebensatz mit „dass“, um die Information deutlich zu machen.} « Le journaliste rapporte que l'équipe a gagné le match ; il utilise une subordonnée en \emph{dass} pour rendre l'information claire. » On voit la différence : la deuxième phrase parle du \emph{travail du journaliste}, pas seulement du match.
\end{exoresolu}

\begin{attention}[Pièges fréquents des francophones]
\begin{itemize}
\item \textbf{Place du verbe} : dans la subordonnée, le verbe conjugué va à la fin : \all{\ldots, weil er eine Geschichte \textbf{erzählt}.} et non \emph{weil er erzählt eine Geschichte}.
\item \textbf{Faux ami} : \all{der Artikel} = l'article de journal, mais aussi l'article grammatical (\all{der, die, das}) ; \all{das Ereignis} = l'événement, pas « l'événementiel ».
\item \textbf{Citation} : utilisez les guillemets allemands „\ldots“ et faites suivre le verbe de dire : \all{Der Journalist schreibt: „\ldots“}.
\item \textbf{Temps} : dans votre commentaire, parlez du texte au présent (\all{Der Journalist benutzt\ldots}) et citez le prétérit de l'article sans le « corriger ».
\end{itemize}
\end{attention}

\courssub{Grille d'auto-évaluation}

Avant de rendre votre commentaire, vérifiez chaque point :

\begin{center}
\begin{tabularx}{\linewidth}{|X|l|}
\hline
\rowcolor{popgreenL} Question de contrôle & Oui / Non \\
\hline
Ai-je présenté le texte (type, thème, les 5 W) ? & \\
\hline
Ai-je cité le texte entre guillemets pour justifier chaque remarque ? & \\
\hline
Ai-je utilisé le vocabulaire du thème (\all{die Überschrift}, \all{berichten}, \all{das Ereignis}) ? & \\
\hline
Mes verbes sont-ils bien placés (2\textsuperscript{e} position ; fin dans les subordonnées) ? & \\
\hline
Ai-je analysé les temps (prétérit de narration) ? & \\
\hline
Ai-je terminé par une appréciation personnelle avec \all{weil} ? & \\
\hline
\end{tabularx}
\end{center}

\begin{aretenir}[L'essentiel de la méthode]
Un commentaire de texte journalistique suit quatre étapes : \textbf{présenter} (type, source, 5 W), \textbf{analyser la structure} (titre, chapeau, corps, ton objectif), \textbf{analyser la langue} (prétérit, champ lexical, subordonnées) et \textbf{apprécier} (avis personnel lié à son contexte). Chaque remarque s'appuie sur une citation entre guillemets et sur une justification en \all{weil} ou \all{dass}. Vous êtes maintenant prêts pour la dernière étape de la leçon : rédiger vous-mêmes un court article de journal.
\end{aretenir}

\coursec{Production guidée : écrire un court article}

Dans la section précédente, nous avons appris à \emph{lire} et à \emph{commenter} un article de journal : identifier la \all{Überschrift}, repérer les 5 W, analyser la structure. Nous allons maintenant franchir le pas décisif du journaliste : passer de la lecture à l'\emph{écriture}. À la fin de cette séquence, chacun d'entre vous aura rédigé son propre article de journal en allemand, comme un vrai reporter du lycée de Yaoundé ou de Douala.

\courssub{La tâche : tu es journaliste !}

\begin{definition}[La tâche de production]
Tu rédiges un article de journal de \cle{6 à 8 phrases} sur un événement \textbf{réel} de ton établissement : un match de football, une fête scolaire, la visite d'une personnalité, un concours de lecture, la journée sportive. Ton article doit respecter le plan imposé du journaliste :
\begin{enumerate}
\item une \all{Überschrift} (un titre) ;
\item une phrase-chapeau (\all{Vorspann}) qui répond aux \cle{5 W} ;
\item 3 à 4 phrases de corps (\all{Körper}) au prétérit ou au parfait ;
\item une citation entre guillemets (\all{Zitat}) au parfait.
\end{enumerate}
\end{definition}

\begin{popfigure}
\begin{center}
\begin{tikzpicture}[
box/.style={draw=popblue, thick, rounded corners=3pt, fill=popblueL, align=center, text width=10.5cm, inner sep=6pt},
lab/.style={font=\bfseries, text=popdark}]
\node[box, fill=popgoldL, draw=popgold] (u) at (0,0) {\textbf{1. ÜBERSCHRIFT} --- der Titel : court, accrocheur, majuscules aux noms};
\node[box, fill=poporangeL, draw=poporange, below=0.35cm of u] (v) {\textbf{2. VORSPANN} --- 1 phrase avec les 5 W : Wer? Was? Wann? Wo? Warum?};
\node[box, below=0.35cm of v] (k) {\textbf{3. KÖRPER} --- 3--4 Sätze im Präteritum/Perfekt : Details, Reaktionen, mindestens 2 Nebensätze (weil, dass)};
\node[box, fill=popgreenL, draw=popgreen, below=0.35cm of k] (z) {\textbf{4. ZITAT} --- wörtliche Rede im Perfekt : „Wir haben gut gespielt“, sagte der Kapitän.};
\draw[-{Stealth}, thick, popblue] (u) -- (v);
\draw[-{Stealth}, thick, popblue] (v) -- (k);
\draw[-{Stealth}, thick, popblue] (k) -- (z);
\end{tikzpicture}
\end{center}
\legende{Le gabarit de l'article de journal : quatre blocs obligatoires, du titre à la citation.}
\end{popfigure}

\courssub{Observer d'abord : un article modèle}

Avant d'écrire, observons un article complet rédigé selon ce gabarit. Lis-le deux fois : une fois pour le sens global, une fois pour repérer les quatre blocs et le vocabulaire cible.

\begin{textebox}[Ein Fußballspiel an unserer Schule]
\textbf{Unser Team gewinnt das große Finale}

Dieser \textbf{Artikel} \textbf{berichtet} über ein wichtiges \textbf{Ereignis} an unserer Schule. Der \textbf{Titel} lautet: Unser Team gewinnt das große Finale. Viele \textbf{Leser} werden sich freuen.

Am Freitagnachmittag hat auf dem Sportplatz unserer Schule das Finale des Fußballturniers stattgefunden. Die Klasse 1ère A hat gegen die Klasse Tle C gespielt und mit 2:1 gewonnen.

Viele Schüler und Lehrer kamen auf den Platz, weil das Spiel sehr wichtig war. Um 15 Uhr hat der Schiedsrichter gepfiffen, und das Spiel begann. In der ersten Halbzeit war unser Team müde, aber in der zweiten Halbzeit haben wir besser gespielt. Der Direktor sagte nach dem Spiel, dass er sehr stolz auf die Schüler war. Am Ende gab es Musik, und alle haben getanzt und gelacht.

„Wir haben bis zum letzten Moment gekämpft, und ich bin sehr glücklich“, sagte der Kapitän nach dem Spiel.
\end{textebox}

\textbf{Glossaire :} der Artikel (-) = l'article ; berichten (berichtete, hat berichtet) = rapporter, faire un reportage ; das Ereignis (-se) = l'événement ; der Titel (-) = le titre ; die Überschrift (-en) = le titre (de journal) ; der Leser (-) = le lecteur ; das Finale (-n) = la finale ; das Fußballturnier (-e) = le tournoi de football ; stattfinden (fand statt, hat stattgefunden) = avoir lieu ; der Schiedsrichter (-) = l'arbitre ; pfeifen (pfiff, hat gepfiffen) = siffler ; die Halbzeit (-en) = la mi-temps ; stolz (auf + Akk.) = fier (de) ; kämpfen (kämpfte, hat gekämpft) = se battre, lutter.

\textbf{Observation guidée :} retrouve dans ce texte (1) le titre (\all{Überschrift}), (2) la phrase-chapeau avec les 5 W, (3) les deux subordonnées avec \all{weil} et \all{dass}, (4) la citation au parfait, (5) les cinq mots du lexique de la presse (\all{der Artikel, berichten, das Ereignis, der Titel, der Leser}).

\courssub{La méthode : des 5 W au texte complet}

\begin{methode}[Écrire son article en quatre étapes]
\begin{enumerate}
\item \textbf{Au brouillon, note les 5 W} en mots isolés, sans phrase complète : \all{Wer?} die Klasse 1ère A --- \all{Was?} ein Konzert --- \all{Wann?} am Samstag --- \all{Wo?} in der Aula --- \all{Warum?} das Schulfest.
\item \textbf{Transforme les 5 W en phrase-chapeau} : une seule phrase qui contient tout. Exemple : \all{Am Samstag hat die Klasse 1ère A in der Aula ein Konzert für das Schulfest gegeben.}
\item \textbf{Rédige le corps} : 3 à 4 phrases au prétérit (\all{war, kam, gab}) ou au parfait (\all{hat begonnen, haben gespielt}), avec au moins deux subordonnées introduites par \all{weil} et \all{dass} (verbe conjugué à la fin !).
\item \textbf{Ajoute la citation} au parfait entre guillemets allemands „…“, puis relis : verbe en position 2, majuscules aux noms, participes passés corrects.
\end{enumerate}
\end{methode}

\begin{exemplebox}[Des phrases de démarrage traduites]
\begin{itemize}
\item \all{Am Freitag hat ein Konzert an unserer Schule stattgefunden.} = Vendredi, un concert a eu lieu dans notre école.
\item \all{Die Schüler waren glücklich, weil die Musik schön war.} = Les élèves étaient heureux parce que la musique était belle.
\item \all{Viele Gäste kamen, weil das Fest sehr bekannt war.} = Beaucoup d'invités sont venus parce que la fête était très connue.
\item \all{Der Direktor sagte: „Ich bin sehr stolz auf euch.“} = Le proviseur a dit : « Je suis très fier de vous. »
\item \all{Am Ende gab es ein Foto mit allen Schülern.} = À la fin, il y a eu une photo avec tous les élèves.
\end{itemize}
\end{exemplebox}

\textbf{Banque de phrases de démarrage} (à recopier et compléter) :

\begin{center}
\begin{tabularx}{\linewidth}{|X|X|}
\hline
\rowcolor{popblueL} \textbf{Début de phrase (Deutsch)} & \textbf{Emploi} \\
\hline
\all{Am ... hat ... stattgefunden.} & phrase-chapeau (Wann? Was?) \\
\hline
\all{Viele ... kamen, weil ...} & corps : les participants + raison \\
\hline
\all{Zuerst ..., dann ..., am Ende ...} & corps : la chronologie \\
\hline
\all{Der Direktor sagte, dass ...} & corps : réaction (subordonnée) \\
\hline
\all{Der Direktor sagte: „...“} & citation directe au parfait \\
\hline
\all{Alle haben gesagt, dass das Fest schön war.} & conclusion \\
\hline
\end{tabularx}
\end{center}

\courssub{Les contraintes linguistiques : ta check-list}

Ton article doit obligatoirement contenir :

\begin{center}
\begin{tabularx}{\linewidth}{|X|X|X|}
\hline
\rowcolor{popblueL} \textbf{Contrainte} & \textbf{Exemple allemand} & \textbf{Traduction} \\
\hline
2 subordonnées (\all{weil}, \all{dass}) & \all{Die Schüler kamen, weil das Spiel wichtig war.} & Les élèves sont venus parce que le match était important. \\
\hline
Verbe à la fin de la subordonnée & \all{Der Direktor sagte, dass er stolz war.} & Le proviseur a dit qu'il était fier. \\
\hline
5 mots du lexique de la presse & \all{der Artikel, der Titel, das Ereignis, berichten, der Leser} & l'article, le titre, l'événement, rapporter, le lecteur \\
\hline
Prétérit (\all{sein, haben, es gab}) & \all{Es gab viele Gäste.} & Il y avait beaucoup d'invités. \\
\hline
Parfait (auxiliaire + participe) & \all{Unser Team hat gewonnen.} & Notre équipe a gagné. \\
\hline
Citation au parfait & \all{„Wir haben gut gespielt“, sagte der Kapitän.} & « Nous avons bien joué », a dit le capitaine. \\
\hline
\end{tabularx}
\end{center}

\begin{attention}[La citation directe : le parfait, pas le prétérit !]
Dans la \cle{citation directe} (la langue \emph{parlée} entre guillemets), l'allemand utilise le \textbf{parfait}, comme le français parlé : \all{„Wir haben gut gespielt“, sagte der Kapitän.} = « Nous avons bien joué », a dit le capitaine. N'écris pas \all{„Wir spielten gut“} dans une citation : le prétérit est réservé à la narration écrite du corps de l'article. Attention aussi aux guillemets allemands „…“ (ouvrants en bas, fermants en haut) et à la virgule avant \all{sagte}. Enfin, après \all{sagte}, le sujet suit le verbe : \all{sagte der Kapitän} (inversion obligatoire car la citation occupe la première place).
\end{attention}

\begin{exoresolu}[Transformer un brouillon en article]
Énoncé : voici les notes au brouillon d'un élève. Rédige la phrase-chapeau et une subordonnée. \all{Wer: die Klasse 1ère A / Was: ein Schulfest / Wann: am Samstag / Wo: im Schulhof / Warum: das Ende des Schuljahres. Grund für die Freude: das Wetter war schön.}
\tcblower
\textbf{Solution détaillée.} Étape 1 : la phrase-chapeau rassemble les 5 W dans une seule phrase, verbe en position 2, participe passé à la fin : \all{Am Samstag hat die Klasse 1ère A im Schulhof ein Schulfest für das Ende des Schuljahres organisiert.} (Samedi, la classe de 1ère A a organisé une fête scolaire dans la cour pour la fin de l'année scolaire.) Étape 2 : la subordonnée avec \all{weil}, verbe conjugué à la fin : \all{Alle Schüler waren froh, weil das Wetter schön war.} (Tous les élèves étaient contents parce que le temps était beau.) Vérification : verbe en position 2 dans la principale (\all{hat}), verbe à la fin dans la subordonnée (\all{war}), majuscules aux noms (\all{Samstag, Klasse, Schulhof, Schulfest, Ende, Schuljahres, Wetter}).
\end{exoresolu}

\courssub{Production orale : tu es à la radio !}

\begin{experience}[Der Radioreporter]
Après la rédaction, chaque élève présente son article à la classe \textbf{comme un journaliste à la radio} : debout, voix claire, sans lire mot pour mot si possible. Commence par la formule du journaliste : \all{Guten Tag, liebe Hörer! Ich berichte heute von unserer Schule.} (Bonjour, chers auditeurs ! Je vous parle aujourd'hui depuis notre école.) Termine par : \all{Das war unser Bericht. Auf Wiederhören!} (C'était notre reportage. Au revoir !) La classe écoute et note les 5 W de chaque reportage : qui a donné l'information la plus complète ?
\end{experience}

\courssub{Correction collective au tableau}

\begin{exercice}[Corrige l'article anonyme]
Le professeur écrit au tableau l'article (anonyme) d'un élève, avec ses erreurs. En groupes, trouvez et corrigez les erreurs. Voici un exemple d'article fautif à corriger : \all{Am Montag die Klasse hat ein Spiel gewonnen. Viele Schüler kamen, weil das Spiel war wichtig. Der Kapitän sagte: „Wir spielten gut.“ Die Direktor war stolz.}
\end{exercice}

\begin{corrige}[Exercice : correction de l'article anonyme]
Quatre erreurs : (1) \all{Am Montag die Klasse hat...} : le verbe doit être en position 2 → \all{Am Montag hat die Klasse ein Spiel gewonnen.} (2) \all{weil das Spiel war wichtig} : dans la subordonnée, le verbe conjugué va à la fin → \all{weil das Spiel wichtig war.} (3) \all{„Wir spielten gut“} : dans une citation directe (langue parlée), on utilise le parfait → \all{„Wir haben gut gespielt.“} (4) \all{Die Direktor} : faux genre → \all{der Direktor}. Article corrigé : \all{Am Montag hat die Klasse ein Spiel gewonnen. Viele Schüler kamen, weil das Spiel wichtig war. Der Kapitän sagte: „Wir haben gut gespielt.“ Der Direktor war stolz.}
\end{corrige}

\begin{savaistu}[La relecture : la clé des points au Bac]
Au baccalauréat, la production écrite est notée sur quatre critères : le \textbf{contenu} (respect de la tâche), le \textbf{vocabulaire} (richesse et exactitude), la \textbf{grammaire} (place du verbe, temps, subordonnées) et l'\textbf{orthographe} (majuscules, Umlaute). Une erreur de place du verbe ou une majuscule oubliée coûte des points, même si ton idée est excellente. Les correcteurs allemands ont un dicton : \all{Der Teufel steckt im Detail} — le diable se cache dans les détails. Garde donc toujours cinq minutes pour relire : verbe en position 2 ? verbe à la fin après \all{weil} et \all{dass} ? majuscule à chaque nom ? Relire n'est pas une option : c'est obligatoire !
\end{savaistu}

\begin{aretenir}[L'essentiel de la production]
Un article de journal se construit en quatre blocs : \all{Überschrift} → \all{Vorspann} (les 5 W en une phrase) → \all{Körper} (3--4 phrases au prétérit/parfait, avec \all{weil} et \all{dass}) → \all{Zitat} (citation au parfait entre „…“). La méthode gagnante : d'abord les 5 W au brouillon, ensuite les phrases complètes, enfin la relecture. Le journaliste allemand écrit pour son \all{Leser} : clarté, ordre des mots et exactitude avant tout.
\end{aretenir}

\newpage

\coursec{Activité d'intégration}

\courssub{Objectifs de l'activité}

Cette activité d'intégration clôt la leçon sur l'article de journal. Elle mobilise, dans une tâche authentique et collective, l'ensemble des savoir-faire travaillés :

\begin{itemize}
\item \textbf{Compréhension} : relire et exploiter le texte support de la leçon comme modèle de structure (titre, chapeau, corps d'article).
\item \textbf{Lexique} : employer correctement le vocabulaire de la presse (\all{der Artikel, die Überschrift, der Leser, berichten, das Ereignis, die Redaktion}).
\item \textbf{Grammaire} : conjuguer au \cle{prétérit} pour le récit des faits, au \cle{parfait} pour la citation d'un témoin, et construire au moins deux \cle{subordonnées} (avec \all{weil}, \all{dass}, \all{als} ou un pronom relatif).
\item \textbf{Expression écrite et orale} : rédiger un court article de 6 à 8 phrases, puis le présenter à voix haute comme dans un journal radiophonique.
\item \textbf{Compétence sociale} : travailler en groupe, répartir les rôles, voter et justifier un choix.
\end{itemize}

\courssub{Situation}

Le lycée bilingue de Yaoundé vient de vivre un trimestre riche en événements : un tournoi de football inter-classes, une journée culturelle germano-camerounaise, la visite d'un délégué du Goethe-Institut. Le professeur d'allemand propose à la classe de Première A de créer la une d'un journal scolaire en langue allemande, \all{Die Schülerzeitung}, qui sera affichée dans le hall et envoyée au partenaire allemand du lycée, un Gymnasium de Bonn.

Chaque groupe de quatre élèves forme une petite rédaction (\all{die Redaktion}) et rédige un article sur l'un de ces événements réels. Pour guider les jeunes journalistes, le professeur distribue le modèle suivant, publié dans le journal d'un établissement partenaire :

\begin{textebox}[Modèle : un article du journal partenaire de Bonn]
\textbf{Großer Sieg beim Schulfest}

BONN. Am Samstag hat unsere Schule das jährliche Sportfest gewonnen. Über 200 Schüler kamen um 9 Uhr auf den Sportplatz, weil das Wetter sehr schön war. Die Mannschaft der Klasse 11a spielte gegen die Klasse 11b und gewann das Finale mit 3:1. Der Trainer, der das Team seit September vorbereitet hat, war sehr stolz. „Wir haben hart trainiert und nie aufgegeben“, sagte der Kapitän nach dem Spiel. Die Schüler sangen und tanzten, als der Pokal übergeben wurde. Der Direktor gratulierte allen Spielern und lud die Sieger zu einem Festessen ein.
\end{textebox}

\begin{definition}[Glossaire du modèle]
\begin{itemize}
\item \all{das Sportfest, -e} : la fête du sport
\item \all{die Mannschaft, -en} : l'équipe
\item \all{der Trainer, -} : l'entraîneur
\item \all{der Kapitän, -e} : le capitaine
\item \all{der Pokal, -e} : la coupe, le trophée
\item \all{übergeben (übergab, hat übergeben)} : remettre
\item \all{gratulieren (gratulierte, hat gratuliert)} : féliciter
\item \all{einladen (lud ein, hat eingeladen)} : inviter
\end{itemize}
\end{definition}

\courssub{Consignes}

Travaillez par groupes de quatre. Répartissez d'abord les rôles : un rédacteur en chef (\all{der Chefredakteur}) coordonne, deux journalistes rédigent, un correcteur vérifie la grammaire et l'orthographe.

\begin{enumerate}
\item \textbf{Observer le modèle.} Relisez l'article de Bonn. Soulignez le titre (\all{die Überschrift}), le chapeau (la première phrase, qui résume les faits), et le corps de l'article. Relevez tous les verbes et indiquez leur temps : prétérit ou parfait. Que remarquez-vous sur la citation du capitaine ?

\item \textbf{Choisir l'événement.} En groupe, choisissez un événement réel du trimestre dans votre lycée : le tournoi de football, la journée culturelle, la visite du délégué allemand, une compétition de chant, la fête de la jeunesse. Notez les informations essentielles : qui ? quoi ? quand ? où ? pourquoi ? (\all{wer? was? wann? wo? warum?}).

\item \textbf{Préparer le lexique.} Complétez votre carnet de vocabulaire avec au moins huit mots du thème utiles à votre article : \all{das Ereignis, der Leser, berichten, die Mannschaft, der Besucher, das Konzert, gewinnen, stattfinden, teilnehmen, sich freuen}.

\item \textbf{Rédiger le chapeau.} Écrivez une ou deux phrases qui résument l'événement au \cle{prétérit} : lieu, date, fait principal. Exemple de départ : \all{YAOUNDÉ. Am Freitag fand in unserer Schule ... statt.}

\item \textbf{Rédiger le corps de l'article.} Racontez les faits au \cle{prétérit} en 3 à 4 phrases. Insérez au moins \textbf{deux subordonnées} : une avec \all{weil} ou \all{als}, une avec \all{dass} ou un pronom relatif (\all{der, die, das}). Attention à la place du verbe conjugué : en dernière position dans la subordonnée.

\item \textbf{Ajouter une citation.} Interrogez (réellement ou par imagination réaliste) un témoin : un joueur, un élève, un professeur. Rapportez sa parole entre guillemets allemands „…“ au \cle{parfait}, suivi de \all{sagte ...} : \all{„Wir haben sehr gut gespielt“, sagte der Kapitän.}

\item \textbf{Trouver un titre.} Proposez deux titres courts et frappants (\all{die Überschrift}), puis choisissez le meilleur en groupe. Le titre ne doit pas être une phrase complète : \all{Großer Sieg beim Turnier}, \all{Ein Fest der Kulturen}.

\item \textbf{Corriger et mettre au propre.} Le correcteur vérifie : majuscules des noms, Umlaute, conjugaison du prétérit, place du verbe dans les subordonnées, guillemets allemands. Recopiez l'article sur une grande feuille avec le titre en gros caractères.

\item \textbf{Présenter à l'oral.} Affichez votre article au mur. Chaque groupe lit son texte à voix haute, lentement, comme un journal radiophonique (\all{die Nachrichten}). La classe vote ensuite pour le meilleur titre et justifie son choix en une phrase : \all{Ich wähle diesen Titel, weil er ...}
\end{enumerate}

\courssub{Ressources à mobiliser}

\begin{aretenir}[Boîte à outils de l'article de journal]
\textbf{Structure} : Überschrift (titre) → chapeau (qui, quoi, quand, où) → corps (détails, déroulement) → citation d'un témoin → conclusion brève.

\textbf{Prétérit} (récit des faits) : \all{sein → war, haben → hatte, kommen → kam, spielen → spielte, gewinnen → gewann, singen → sang, finden → fand} (avec \all{statt} : \all{fand ... statt}, particule en fin de phrase).

\textbf{Parfait} (parole rapportée, citation) : \all{haben/sein + Partizip II} : \all{„Wir haben gewonnen.“}, \all{„Wir sind um 8 Uhr gekommen.“}

\textbf{Subordonnées} (verbe conjugué en fin de phrase) :
\begin{itemize}
\item \all{weil} + verbe final : \all{Viele Schüler kamen, weil das Fest interessant war.}
\item \all{als} + verbe final (passé, moment unique) : \all{Alle jubelten, als das Spiel endete.}
\item \all{dass} + verbe final : \all{Der Direktor sagte, dass er stolz war.}
\item Relatif : \all{Der Spieler, der das Tor schoss, heißt Alain.}
\end{itemize}

\textbf{Lexique de la presse} : \all{der Artikel, - ; die Überschrift, -en ; der Titel, - ; der Leser, - ; die Redaktion, -en ; berichten ; das Ereignis, -se ; stattfinden ; teilnehmen (an + Dat.)}.
\end{aretenir}

\begin{attention}[Pièges à éviter]
\begin{itemize}
\item Ne dites pas \all{„Das Fest hat stattgefunden“} dans le corps au prétérit : préférez \all{Das Fest fand statt} (particule \all{statt} en fin de phrase).
\item Dans la subordonnée, le verbe conjugué va \textbf{à la fin} : \all{weil das Wetter schön war} et non \all{weil das Wetter war schön}.
\item La citation est au parfait (oral), le récit au prétérit (écrit) : ne mélangez pas les deux dans la même partie.
\item Guillemets allemands : „ … “ (ouvrant en bas, fermant en haut).
\item Majuscule à tous les noms : \all{das Fest, der Sieg, die Schüler}.
\end{itemize}
\end{attention}

\courssub{Proposition de production et corrigé commenté}

\begin{exoresolu}[Article modèle : la journée culturelle]
Rédigez un article de 6 à 8 phrases sur la journée culturelle germano-camerounaise de votre lycée, en respectant la structure étudiée.
\tcblower
\textbf{Ein Fest der Kulturen in Yaoundé}

YAOUNDÉ. Am Donnerstag fand in unserem Lycée die deutsch-kamerunische Kulturtag statt. Mehr als 300 Schüler und Lehrer kamen in die Aula, weil das Programm sehr reich war. Die Schüler der Klasse Première A, die ein deutsches Theaterstück vorbereitet hatten, spielten vor einem begeisterten Publikum. Danach tanzten die Schüler zu deutscher und kamerunischer Musik, und die Mamas des APEEE servierten Ndolé und deutschen Kuchen. Der Delegierte des Goethe-Instituts, der aus Douala gekommen war, besuchte unsere Ausstellung. „Ich habe viel gelernt und sehr gut gegessen“, sagte ein Schüler nach dem Fest. Der Direktor dankte allen Helfern und sagte, dass die Kulturtag jetzt jedes Jahr stattfinden wird.

\textbf{Commentaire des choix de langue :}
\begin{itemize}
\item \textbf{Titre} : nominal, court, sans verbe conjugué — typique de la presse.
\item \textbf{Chapeau} : \all{fand ... statt} au prétérit, avec lieu (\all{YAOUNDÉ}) et date ; la particule séparable \all{statt} est bien en fin de phrase.
\item \textbf{Corps au prétérit} : \all{kamen, war, spielten, tanzten, servierten, besuchte, dankte, sagte} — temps du récit journalistique.
\item \textbf{Subordonnées} : \all{weil das Programm sehr reich war} (cause, verbe final) ; \all{die ein deutsches Theaterstück vorbereitet hatten} (relative, verbe final) ; \all{der aus Douala gekommen war} (relative) ; \all{dass die Kulturtag jetzt jedes Jahr stattfinden wird} (discours rapporté avec \all{dass}).
\item \textbf{Citation au parfait} : \all{„Ich habe viel gelernt und sehr gut gegessen“} — le parfait est le temps de l'oral ; guillemets allemands „ … “ ; \all{sagte} revient au prétérit car c'est le narrateur qui parle.
\item \textbf{Lexique du thème} : \all{das Fest, das Publikum, die Ausstellung, der Delegierte, das Programm}.
\end{itemize}
\end{exoresolu}

\courssub{Bilan de l'activité}

À l'issue de cette activité, chaque élève a produit, en équipe, un véritable article de presse en allemand : un texte court mais complet, structuré comme un modèle professionnel. Les acquis sont vérifiés sur quatre plans :

\begin{itemize}
\item \textbf{Lexique} : les mots de la presse et de l'événementiel ont été employés dans un contexte réel, non simplement mémorisés.
\item \textbf{Grammaire} : la complémentarité prétérit (récit écrit) / parfait (parole citée) a été mise en pratique, ainsi que la construction de subordonnées avec le verbe en position finale.
\item \textbf{Méthode} : la structure titre–chapeau–corps–citation est désormais un réflexe transferable à toute production écrite journalistique et à l'épreuve de rédaction.
\item \textbf{Oral} : la lecture à voix haute façon journal radiophonique a entraîné la prononciation, l'intonation et l'aisance devant la classe.
\end{itemize}

Le vote pour le meilleur titre a enfin développé l'esprit critique : justifier un choix esthétique en allemand (\all{Ich wähle diesen Titel, weil ...}) est déjà un premier exercice d'argumentation. Les articles affichés au mur du hall prolongent l'activité : les lecteurs du lycée deviennent, à leur tour, un public germanophone réel.

\newpage

\coursec{Méthodes et exercices résolus}

\coursec{Le texte support : un article de journal}

\courssub{Contexte et lecture globale}
\begin{textebox}[Article de journal simulé – Contexte camerounais (Cameroon Tribune / Deutsche Welle)]
„Neuer Markt in Yaoundé eröffnet

YAOUNDE. Der neue Zentralmarkt in der Hauptstadt Kameruns wurde gestern feierlich eröffnet. Der Bürgermeister der Stadt, Herr Mbarga, schnitt das rote Band vor hunderten Zuschauern durch. Der Markt bietet Platz für über 500 Händler und verfügt über moderne sanitäre Anlagen sowie ein Dach, das vor Regen schützt. Bisher verkauften viele Händler ihre Waren am Straßenrand, was den Verkehr behinderte. „Endlich haben wir einen würdigen Arbeitsplatz“, sagte die Händlerin Marie Nkou. Die Bauarbeiten dauerten zwei Jahre und kosteten 3,5 Milliarden CFA-Franc. Die Finanzierung erfolgte durch die Weltbank und den kamerunischen Staat. Der Markt soll die Hygiene verbessern und den Handel in der Region ankurbeln.“
\end{textebox}

\begin{definition}[Glossaire du texte (mots difficiles)]
\begin{itemize}
    \item \cle{der Zentralmarkt, die -märkte} : marché central.
    \item \cle{feierlich} : solennel(le).
    \item \cle{der Bürgermeister, die -} : le maire.
    \item \cle{durchschneiden} (trennbar, fort : \all{schnitt durch – hat durchgeschnitten}) : couper (un ruban).
    \item \cle{der Händler, die -} : le commerçant, le vendeur.
    \item \cle{die sanitären Anlagen} (Pl.) : les installations sanitaires (toilettes, lavabos).
    \item \cle{der Straßenrand, die -änder} : le bord de la route.
    \item \cle{behindern} (faible) : gêner, entraver.
    \item \cle{würdig} : décent, digne.
    \item \cle{die Bauarbeiten} (Pl.) : les travaux de construction.
    \item \cle{der CFA-Franc} (invariable) : le franc CFA.
    \item \cle{die Finanzierung} : le financement.
    \item \cle{erfolgen} (fort : \all{erfolgte – ist erfolgt}) : avoir lieu, être effectuée.
    \item \cle{ankurbeln} (faible) : stimuler, relancer.
\end{itemize}
\end{definition}

\begin{popfigure}
\begin{tikzpicture}[box/.style={draw=popblue, fill=popblueL, rounded corners, text width=9cm, align=center, minimum height=0.9cm}]
\node[box] (a) {\textbf{Titel} (Titre)};
\node[box, below=0.3cm of a] (b) {\textbf{Vorspann / Lead} (Chapeau)};
\node[box, below=0.3cm of b] (c) {\textbf{Hauptteil} (Corps de l'article)};
\node[box, below=0.3cm of c] (d) {\textbf{Zitat / Schluss} (Citation / Conclusion)};
\draw[-{Stealth}, thick, poporange] (a) -- (b);
\draw[-{Stealth}, thick, poporange] (b) -- (c);
\draw[-{Stealth}, thick, poporange] (c) -- (d);
\end{tikzpicture}
\legende{Structure type d'un article de journal (pyramide inversée).}
\end{popfigure}

\coursec{Compréhension du texte}

\courssub{Questions de compréhension globale}
\begin{exercice}[Compréhension globale – \textit{Erstes Verstehen}]
Répondre en français par vrai/faux ou en complétant.
\begin{enumerate}
    \item De quel événement parle l'article ? (Inauguration / Incendie / Manifestation)
    \item Où se déroule la scène ? (Douala / Yaoundé / Bafoussam)
    \item Qui a coupé le ruban ? (Le Président / Le Maire / Un commerçant)
    \item Combien de vendeurs le marché peut-il accueillir ? (100 / 500 / 1000)
    \item Quel était le problème avant ? (Prix trop chers / Vente au bord de la route / Marché trop loin)
\end{enumerate}
\end{exercice}

\courssub{Questions de compréhension détaillée}
\begin{exercice}[Compréhension détaillée – \textit{Genaues Verstehen}]
Répondre en allemand par des phrases complètes ou des éléments de phrase repris du texte.
\begin{enumerate}
    \item Was ist das Thema des Artikels?
    \item Wer hat den Markt eröffnet?
    \item Welche Vorteile hat der neue Markt für die Händler? Nennen Sie zwei Details.
    \item Wie lange dauerten die Bauarbeiten und wer hat gezahlt?
    \item Finden Sie im Text drei Verben im Präteritum und drei formes verbales au Perfekt (ou futur/modal).
    \item Quel connecteur logique introduit la phrase sur le financement ? (\cle{Die Finanzierung erfolgte ...})
\end{enumerate}
\end{exercice}

\begin{corrige}[Exercices de compréhension]
\begin{enumerate}
    \item \textbf{Global :} 1. Inauguration. 2. Yaoundé. 3. Le Maire (Bürgermeister). 4. 500. 5. Vente au bord de la route.
    \item \textbf{Détaillé :}
        \begin{enumerate}
            \item \all{Das Thema ist die Eröffnung eines neuen Zentralmarktes in Yaoundé.}
            \item \all{Der Bürgermeister Herr Mbarga hat den Markt eröffnet (er schnitt das Band durch).}
            \item \all{Der Markt hat moderne sanitäre Anlagen und ein Dach gegen Regen. Die Händler haben endlich einen würdigen Arbeitsplatz (kein Verkauf mehr am Straßenrand).}
            \item \all{Die Bauarbeiten dauerten zwei Jahre und kosteten 3,5 Milliarden CFA-Franc. Die Finanzierung erfolgte durch die Weltbank und den kamerunischen Staat.}
            \item \textbf{Präteritum :} \all{eröffnete, schnitt ... durch, verkauften, behinderte, dauerten, kosteten, erfolgte, sagte}. \textbf{Perfekt (dans citation) :} \all{haben ... gewartet (impliqué dans 'haben wir einen Platz'), soll verbessern / soll ankurbeln (Modal/Futur).} \textit{Note : Le texte narratif utilise presque exclusivement le Präteritum.}
            \item Pas de connecteur conjonctif (subordonnée), mais une phrase principale : \cle{Die Finanzierung erfolgte ...}. (Passif / Verbe d'état).
        \end{enumerate}
\end{enumerate}
\end{corrige}

\coursec{Lexique thématique : la presse et l'article}

\courssub{Vocabulaire de base : noms, verbes, adjectifs}
\begin{definition}[Mots-clés du thème « Presse et Article » (avec article, pluriel, traduction)]
\begin{center}
\begin{tabular}{|l|l|l|}\hline
\rowcolor{popblueL} \textbf{Deutsch (Nom + Article + Pluriel)} & \textbf{Français} & \textbf{Remarque / Famille} \\ \hline
\cle{der Artikel, die Artikel} & l'article (de journal) & \cle{artikeln} (faible) \\ \hline
\cle{die Zeitung, die Zeitungen} & le journal (quotidien) & \cle{Zeitungsartikel} \\ \hline
\cle{der Titel, die Titel} & le titre (général) & Synonyme : \cle{die Überschrift, die -en} \\ \hline
\cle{die Überschrift, die -en} & le titre, la manchette & \cle{über + Akk} \\ \hline
\cle{der Leser, die Leser} & le lecteur & \cle{lesen – las – gelesen} \\ \hline
\cle{der Journalist, die -en} & le journaliste & \cle{der Journalismus} \\ \hline
\cle{der Reporter, die -} & le reporter & \cle{reporten} (anglicisme) \\ \hline
\cle{das Ereignis, die Ereignisse} & l'événement & \cle{ereignen} (réfléchi) \\ \hline
\cle{der Bericht, die Berichte} & le reportage, le compte-rendu & \cle{berichten} \\ \hline
\cle{die Schlagzeile, die -n} & le gros titre (une) & \cle{Schlagzeilen lesen} \\ \hline
\cle{der Vorspann, die Vorspänne} & le chapeau, le lead & \cle{Lead-Absatz} \\ \hline
\cle{die Quelle, die -n} & la source & \cle{aus zuverlässiger Quelle} \\ \hline
\cle{veröffentlichen} (faible) & publier & \cle{die Veröffentlichung} \\ \hline
\cle{erscheinen} (fort: \all{erschien – ist erschienen}) & paraître, sortir & \textbf{Aux. sein !} \\ \hline
\cle{berichten über + Akk} (faible) & reporter sur, faire un article sur & \cle{der Bericht} \\ \hline
\cle{aktualisieren} & actualiser, mettre à jour & \cle{aktuell} (adj.: actuel/recent) \\ \hline
\end{tabular}
\end{center}
\end{definition}

\courssub{Familles de mots, collocations et nominalisation}
\begin{methode}[Mobiliser le lexique de la presse : familles de mots et collocations]
\textbf{Objectif :} Passer du mot isolé à l'expression correcte (genre, pluriel, prépositions, verbes associés) pour l'expression écrite/orale.\par
\begin{enumerate}
    \item \textbf{Apprendre le nom \textbf{avec son article et son pluriel}} : \cle{der Artikel, die Artikel} ; \cle{die Zeitung, die Zeitungen} ; \cle{der Leser, die Leser} ; \cle{das Ereignis, die Ereignisse}.
    \item \textbf{Construire des familles morphologiques} (dérivation) :
        \begin{itemize}
            \item \cle{berichten (Verbum)} $\rightarrow$ \cle{der Bericht, die Berichte} (Substantiv) ; \cle{berichterstatten} (verbe séparable : \all{er berichtet ... / er hat berichtet}).
            \item \cle{veröffentlichen} $\rightarrow$ \cle{die Veröffentlichung} ; \cle{erschienen} (Partizip II adj.) $\rightarrow$ \cle{die Erscheinung}.
            \item \cle{lesen} $\rightarrow$ \cle{der Leser, die Leser} ; \cle{die Lesart, die -en} ; \cle{lesbar}.
        \end{itemize}
    \item \textbf{Mémoriser les collocations (associations figées)} :
        \all{einen Artikel \textbf{veröffentlichen} / \textbf{schreiben} / \textbf{lesen}} ;
        \all{\textbf{über} ein Ereignis \textbf{berichten}} (préposition \cle{über} + Akkusativ) ;
        \all{in der \textbf{Zeitung stehen}} (figurer dans le journal) ;
        \all{die \textbf{Schlagzeile} \textbf{lesen}} ;
        \cle{sich \textbf{für} etwas \textbf{interessieren}} (s'intéresser à qqch – Dativ !).
    \item \textbf{Entraînement – Transformation nominale $\leftrightarrow$ verbale} :
        \all{Die Eröffnung des Marktes} (Nominalisation, Génitif) $\leftrightarrow$ \all{Der Markt wurde eröffnet} (Verbalisation, Passif).
        \all{Die Berichterstattung über das Fest} $\leftrightarrow$ \all{Über das Fest wurde berichtet}.
\end{enumerate}
\cle{Astuce Bac :} Dans le thème, \cle{über} + Akkusatif après \cle{berichten / schreiben / sprechen} (parler de). Ne pas dire \all{*berichten von} (sauf « rendre compte de » au sens large). \cle{sich freuen \textbf{über} + Akk} (présent/passé) vs \cle{sich freuen \textbf{auf} + Akk} (futur).
\end{methode}

\begin{exoresolu}[Thème lexicale et transformation nominale/verbale]
\textbf{Énoncé :} Traduire en allemand les phrases suivantes en utilisant le vocabulaire de la leçon.
\begin{enumerate}
    \item Le journaliste écrit un article sur l'événement.
    \item L'article paraît dans la presse locale.
    \item Les lecteurs lisent le titre avec intérêt.
    \item Le reportage rapporte des faits importants.
    \item Transformation : \all{Die Veröffentlichung des Artikels} $\rightarrow$ phrase verbale (Perfekt, Passiv).
\end{enumerate}
\tcblower
\textbf{Solution détaillée :}
\begin{enumerate}
    \item \all{Der Journalist schreibt einen Artikel \textbf{über} das Ereignis.}
        \textit{Justification :} \cle{Journalist} (masc.) ; \cle{Artikel} (masc., Akk. \cle{einen}) ; \cle{über} + Akk. (\cle{das Ereignis}).
    \item \all{Der Artikel \textbf{erscheint} in der lokalen Presse.} \quad ou \quad \all{Der Artikel \textbf{wird} in der lokalen Presse \textbf{veröffentlicht}.}
        \textit{Justification :} \cle{erscheinen} (fort, intransitif, aux. \cle{sein} : \cle{ist erschienen}) = « paraître ». \cle{veröffentlichen} (faible, transitif) $\rightarrow$ Passif \cle{werden + Partizip II}. \cle{in} + Datif (\cle{der Presse} $\rightarrow$ \cle{der lokalen Presse}).
    \item \all{Die Leser lesen die \textbf{Überschrift} mit Interesse.}
        \textit{Justification :} \cle{Leser} (Pluriel) ; \cle{Überschrift} (fém., synonyme de \cle{Titel} pour journal) ; \cle{mit} + Datif (\cle{Interesse}).
    \item \all{Der Bericht \textbf{behandelt} wichtige Fakten.} \quad (Mieux que \all{berichtet über} pour style écrit). \quad Ou passif : \all{In dem Bericht \textbf{wird über} wichtige Fakten \textbf{berichtet}.}
        \textit{Justification :} \cle{behandeln} (traiter) transitif + Akk. \cle{wichtige Fakten} (Pl. Akk.).
    \item \textbf{Transformation (Nominalisation $\rightarrow$ Verbalisation Passiv Perfekt) :}
        \all{Der Artikel \textbf{ist veröffentlicht worden}.} \quad (Präteritum Passiv : \all{Der Artikel \textbf{wurde veröffentlicht}.})
        \textit{Justification :} \cle{die Veröffentlichung} $\rightarrow$ verbe \cle{veröffentlichen}. Sujet passif = \cle{der Artikel}. Perfekt Passif = \cle{ist + Partizip II (veröffentlicht) + worden}. \cle{worden} sans \cle{ge-}.
\end{enumerate}
\textbf{Glossaire difficile :} \cle{die Presse} (pas de pluriel usuel) ; \cle{die Überschrift, die -en} ; \cle{die Tatsache, die -n} / \cle{der Fakt, die -en} ; \cle{lokal} (adj. local) ; \cle{behandeln} (traiter un sujet).
\end{exoresolu}

\coursec{Grammaire 1 : Le parfait (Perfekt) et le prétérit (Präteritum)}

\courssub{Règle fondamentale : Contexte écrit vs oral}
\begin{propriete}[Choix entre Präteritum et Perfekt – Niveau A2/B1]
\begin{enumerate}
    \item \textbf{Präteritum (Imperfekt) = Temps du récit écrit, de la narration objective, du journalisme, de la littérature.}
        \begin{itemize}
            \item Verbes d'état/auxiliaires : \cle{sein, haben, werden} $\rightarrow$ \textbf{toujours} Präteritum (\all{war, hatte, wurde}).
            \item Verbes modaux : \cle{können, müssen, wollen, sollen, dürfen, mögen} $\rightarrow$ \textbf{préférentiellement} Präteritum (\all{konnte, musste, wollte, sollte, durfte, mochte}).
            \item Verbes d'action dans un article/roman/récit : \all{Der Bürgermeister \textbf{eröffnete} den Markt.}
        \end{itemize}
    \item \textbf{Perfekt = Temps de l'oral, du dialogue, du résumé subjectif, du résultat présent, des citations directes (style direct).}
        \begin{itemize}
            \item Conversation : \all{„Was \textbf{haben Sie gemacht}?“} – \all{Ich \textbf{bin nach Hause gegangen}.}
            \item Citations directes dans l'article : \all{„Endlich \textbf{haben} wir einen Platz“, \textbf{sagte} sie.} (\cle{sagte} = Präteritum verbe introducteur ; \cle{haben} = Perfekt dans la bouche du locuteur).
            \item Verbes de mouvement/changement d'état (auxiliaire \cle{sein}) : \cle{gehen, kommen, werden, sterben, aufstehen, einschlafen}.
        \end{itemize}
    \item \textbf{Test de substitution mental :}
        \begin{itemize}
            \item Si je peux dire \all{„Damals ...“} (À l'époque, récit) $\rightarrow$ \textbf{Präteritum}.
            \item Si je peux dire \all{„Schon ...“} (Déjà) / résultat maintenant / expérience vécue $\rightarrow$ \textbf{Perfekt}.
        \end{itemize}
\end{enumerate}
\cle{Règle d'or :} \textbf{Narration (texte journalistique) = Präteritum.} \textbf{Parole citée / Oral / Résultat présent = Perfekt.}
\end{propriete}

\courssub{Conjugaison des verbes du thème (Tableau de référence)}
\begin{definition}[Verbes forts, faibles et mixtes essentiels pour le thème « Article de journal »]
\begin{center}
\begin{tabular}{|l|c|c|c|c|}\hline
\rowcolor{popblueL} \textbf{Infinitif} & \textbf{Präteritum (ich/er/sie/es)} & \textbf{Partizip II} & \textbf{Auxiliaire} & \textbf{Sens} \\ \hline
\cle{schreiben} & \all{schrieb} & \all{geschrieben} & haben & écrire \\ \hline
\cle{lesen} & \all{las} & \all{gelesen} & haben & lire \\ \hline
\cle{erscheinen} & \all{erschien} & \all{erschienen} & \textbf{sein} & paraître, sortir \\ \hline
\cle{berichten} & \all{berichtete} (faible) & \all{berichtet} & haben & reporter, faire un article \\ \hline
\cle{veröffentlichen} & \all{veröffentlichte} (faible) & \all{veröffentlicht} & haben & publier \\ \hline
\cle{dauern} & \all{dauerte} & \all{gedauert} & haben & durer \\ \hline
\cle{kosten} & \all{kostete} & \all{gekostet} & haben & coûter \\ \hline
\cle{eröffnen} & \all{eröffnete} (faible) & \all{eröffnet} & haben & inaugurer, ouvrir \\ \hline
\cle{durchschneiden} & \all{schnitt durch} & \all{durchgeschnitten} & haben & couper (ruban) \\ \hline
\cle{sagen} & \all{sagte} (faible/irrégulier) & \all{gesagt} & haben & dire \\ \hline
\cle{verkaufen} & \cle{verkaufte} (faible) & \cle{verkauft} & haben & vendre \\ \hline
\cle{behindern} & \cle{behinderte} (faible) & \cle{behindert} & haben & gêner \\ \hline
\cle{erfolgen} & \cle{erfolgte} (fort: a$\rightarrow$o) & \cle{erfolgt} & \textbf{sein} & avoir lieu, s'ensuivre \\ \hline
\cle{verbessern} & \cle{verbesserte} (faible) & \cle{verbessert} & haben & améliorer \\ \hline
\cle{ankurbeln} & \cle{ankurbelte} (faible) & \cle{ankurbelt} & haben & stimuler \\ \hline
\end{tabular}
\end{center}
\cle{Attention :} \cle{erscheinen} et \cle{erfolgen} prennent \textbf{sein} au Perfekt ! \all{Der Artikel ist erschienen. / Die Finanzierung ist erfolgt.}
\end{definition}

\courssub{Exercices d'application : Conjugaison au temps juste}
\begin{exoresolu}[Application : Contexte journalistique vs Témoignage oral]
\textbf{Énoncé :} Compléter le texte suivant en conjuguant les verbes entre parenthèses au \textbf{Präteritum} (récit journalistique) ou au \textbf{Perfekt} (témoignage oral / citation). Justifier chaque choix.
\textbf{Texte :}
„Gestern \_\_\_\_\_\_\_\_\_ (eröffnen) der neue Markt. Die Bauarbeiten \_\_\_\_\_\_\_\_\_ (dauern) zwei Jahre. Eine Händlerin \_\_\_\_\_\_\_\_\_ (sagen): ‚Ich \_\_\_\_\_\_\_\_\_ (warten) lange auf diesen Tag.‘ Der Bürgermeister \_\_\_\_\_\_\_\_\_ (schneiden) das Band. Die Leute \_\_\_\_\_\_\_\_\_ (freuen) sich sehr.“
\tcblower
\textbf{Solution détaillée :}
\begin{enumerate}
    \item \all{eröffnete} (Präteritum). \textit{Justification :} Verbe d'action principal du récit journalistique (\cle{eröffnen} faible). Temps de la narration écrite.
    \item \all{dauerten} (Präteritum). \textit{Justification :} Verbe d'état/durée dans le récit (\cle{dauern} fort : a $\rightarrow$ au). Contexte : fait passé clos, narré.
    \item \all{sagte} (Präteritum). \textit{Justification :} Verbe introducteur de citation (\cle{sagen} faible/irrégulier). Standard journalistique pour introduire le style direct.
    \item \all{habe ... gewartet} (Perfekt : \all{ich habe lange gewartet}). \textit{Justification :} \textbf{Style direct (Zitat)}. La personne parle de son expérience vécue (oralité). \cle{warten} (faible) + \cle{haben}.
    \item \all{schnitt} (Präteritum). \textit{Justification :} Verbe fort \cle{schneiden} (a $\rightarrow$ a). Action ponctuelle, narrative, passée. Sujet : \cle{der Bürgermeister} (3ᵉ pers. sg.). Particule \cle{durch} en fin de proposition principale : \all{schnitt das Band durch}.
    \item \all{freuten} (Präteritum). \textit{Justification :} \cle{sich freuen} (réfléchi, faible). Description de l'état de la foule dans le récit journalistique. (À l'oral : \all{haben sich gefreut}).
\end{enumerate}
\textbf{Règle d'or résumée :} \textbf{Narration (texte) = Präteritum.} \textbf{Parole citée / Oral / Résultat = Perfekt.}
\end{exoresolu}

\coursec{Grammaire 2 : Les subordonnées (Nebensätze)}

\courssub{Structure immuable et connecteurs}
\begin{propriete}[Ordre des mots dans la subordonnée allemande]
\textbf{Structure :} \cle{Conjonction} + \cle{Sujet} + \cle{Compléments (Temps, Manière, Lieu...)} + \cle{VERBE CONJUGUÉ À LA FIN}.\par
\all{Weil der Markt neu \textbf{ist},} (pas *\all{ist der Markt neu}).\par
Si la subordonnée est en \textbf{position 1} (début de phrase), le verbe de la principale vient en \textbf{position 2 (V2)} : \all{Weil der Markt neu \textbf{ist}, \textbf{hat} er viele Kunden.}
\end{propriete}

\begin{definition}[Connecteurs principaux pour subordonnées (Niveau A2/B1)]
\begin{center}
\begin{tabularx}{\linewidth}{|X|X|X|}\hline
\rowcolor{popblueL} \textbf{Fonction} & \textbf{Conjonction (Subordonnant $\rightarrow$ Verbe FIN)} & \textbf{Exemple (verbe en fin)} \\ \hline
Cause & \cle{weil, da} & \all{Da es regnete, blieb ich zu Hause.} \\ \hline
But & \cle{damit} (sujet différent) / \cle{um ... zu} (infinitif, même sujet) & \all{Er kam, \textbf{damit} er hilft.} / \all{Er kam, \textbf{um} zu \textbf{helfen}.} \\ \hline
Temps (antérieur) & \cle{bevor, nachdem} & \all{\textbf{Nachdem} er \textbf{gekommen} \textbf{war}, aß er.} (Plusquamperfekt dans sub.) \\ \hline
Temps (simultané) & \cle{während, als} (ponctuel passé), \cle{wenn} (répété/présent/futur) & \all{\textbf{Als} er \textbf{kam}, regnete es.} \\ \hline
Concession & \cle{obwohl, obgleich} & \all{\textbf{Obwohl} es \textbf{regnete}, ging er raus.} \\ \hline
Condition & \cle{wenn, falls} & \all{\textbf{Wenn} es \textbf{regnet}, bleibe ich.} \\ \hline
Relatif & \cle{der, die, das} (accordé genre/nombre, cas selon fonction) & \all{Der Markt, \textbf{der} gestern \textbf{eröffnet wurde}, ist groß.} \\ \hline
\end{tabularx}
\end{center}
\end{definition}

\begin{definition}[Adverbes conjonctifs (Konjunktionaladverbien $\rightarrow$ V2 !)]
\begin{center}
\begin{tabular}{|l|l|l|}\hline
\rowcolor{popblueL} \textbf{Fonction} & \textbf{Adverbe (V2)} & \textbf{Exemple} \\ \hline
Conséquence & \cle{deshalb, darum, daher, deswegen} & \all{Es regnete. \textbf{Deshalb} \textbf{blieb} ich zu Hause.} \\ \hline
Concession (adversatif) & \cle{trotzdem, jedoch} & \all{Es regnete. \textbf{Trotzdem} \textbf{ging} er raus.} \\ \hline
Ajout & \cle{außerdem, außerdem} & \all{Er ist müde. \textbf{Außerdem} \textbf{hat} er Hunger.} \\ \hline
\end{tabular}
\end{center}
\cle{PIÈGE MAJEUR :} \cle{weil/obwohl/dass} = Subordonnant (Verbe FIN). \cle{deshalb/trotzdem/aber} = Adverbe/Coordination (Verbe V2). \textbf{Jamais les deux ensemble !}
\end{definition}

\courssub{Exercices de transformation et thème}
\begin{exoresolu}[Transformation et thème : Subordonnées (Cause, Concession, But, Temps)]
\textbf{Énoncé :}
\begin{enumerate}
    \item Relier par \cle{weil} (cause) : \all{Der Markt ist modern. Er hat ein Dach.}
    \item Relier par \cle{obwohl} (concession) : \all{Es regnete. Die Eröffnungsfeier fand statt.}
    \item Relier par \cle{damit} (but) : \all{Die Händler arbeiten hart. Sie wollen Geld verdienen.}
    \item Thème (Français $\rightarrow$ Allemand) : « Bien que le marché soit nouveau, il y a encore peu de clients, parce que la route est mauvaise. »
    \item Repérer l'erreur et corriger : \all{*Weil der Markt neu ist, deshalb hat er viele Kunden.}
\end{enumerate}
\tcblower
\textbf{Solution détaillée :}
\begin{enumerate}
    \item \all{Der Markt ist modern, \textbf{weil er ein Dach hat}.}
        \textit{Justification :} \cle{weil} $\rightarrow$ verbe \cle{hat} à la fin. Sujet \cle{er}.
    \item \all{\textbf{Obwohl} es regnete, fand die Eröffnungsfeier statt.} \quad (ou \all{Die Eröffnungsfeier fand statt, \textbf{obwohl} es regnete.})
        \textit{Justification :} \cle{obwohl} + verbe \cle{regnete} à la fin. \cle{stattfinden} (séparable) $\rightarrow$ \cle{fand ... statt} (Präteritum). Subordonnée en position 1 $\rightarrow$ Verbe principal V2 (\cle{fand}).
    \item \all{Die Händler arbeiten hart, \textbf{damit sie Geld verdienen}.}
        \textit{Justification :} \cle{damit} + sujet \cle{sie} + verbe \cle{verdienen} à la fin. (Alternative infinitif si même sujet : \all{um Geld zu verdienen}).
    \item \textbf{Thème :} \all{\textbf{Obwohl} der Markt neu \textbf{ist}, gibt es noch wenige Kunden, \textbf{weil} die Straße schlecht \textbf{ist}.}
        \textit{Analyse :}
        \begin{itemize}
            \item « Bien que » = \cle{obwohl} (concession) $\rightarrow$ verbe \cle{ist} à la fin.
            \item « il y a » = \cle{es gibt} (impersonnel, invariables) + Accusatif (\cle{wenige Kunden}).
            \item « parce que » = \cle{weil} (cause) $\rightarrow$ verbe \cle{ist} à la fin.
            \item \cle{die Straße} (fém.) $\rightarrow$ \cle{schlecht} (adj. épithète/attribut non décliné après \cle{ist}).
        \end{itemize}
    \item \textbf{Erreur :} Double connecteur de cause (\cle{weil} + \cle{deshalb}) + ordre des mots incorrect.
        \textbf{Correction A (subordonnée seule) :} \all{\textbf{Weil} der Markt neu \textbf{ist}, hat er viele Kunden.}
        \textbf{Correction B (adverbe seul) :} \all{Der Markt ist neu. \textbf{Deshalb} \textbf{hat} er viele Kunden.}
        \textit{Règle :} On ne mélange pas \cle{weil} (subordonnant, verbe fin) et \cle{deshalb} (adverbe, V2). On choisit l'un ou l'autre.
\end{enumerate}


\end{exoresolu}

\coursec{Méthode du commentaire et commentaire modèle}

\courssub{Démarche en trois lectures}
\begin{methode}[Comprendre un article de journal : démarche en trois lectures]
\textbf{Objectif :} Extraire l'information essentielle (qui, quoi, quand, où, pourquoi) et comprendre la structure journalistique (titre, chapeau, corps, conclusion).\par
\begin{enumerate}
    \item \textbf{1ʳᵉ lecture (globale) :} Lire le \cle{Titel} (titre) et le \cle{Vorspann} (chapeau/sous-titre). Identifier le \cle{Thema} (sujet) et le \cle{Textsort} (type de texte : reportage, commentaire, brève). Se demander : \all{„Worum geht es?“} (« De quoi s'agit-il ? »).
    \item \textbf{2ᵉ lecture (détaillée) :} Lire par paragraphes (\cle{Absätze}). Souligner les \cle{Schlüsselwörter} (mots-clés) : noms propres, chiffres, verbes d'action. Repérer les connecteurs logiques (\cle{deshalb, jedoch, außerdem}). Noter les \cle{W-Fragen} (questions QQOQCP) : \all{Wer? Was? Wann? Wo? Warum? Wie?}
    \item \textbf{3ᵉ lecture (vérification) :} Relire pour valider les réponses aux questions de compréhension. Vérifier les \cle{Falsche Freunde} (faux amis) et les temps de narration (\cle{Präteritum} pour le récit, \cle{Perfekt} pour le témoignage oral/résumé).
\end{enumerate}
\cle{Réflexe Bac :} Ne jamais traduire mot à mot. Répondre en allemand par des phrases complètes ou des éléments de phrase repris du texte.
\end{methode}

\courssub{Commentaire modèle de l'article « Neuer Markt in Yaoundé »}
\begin{exoresolu}[Commentaire structuré – \textit{Textanalyse}]
\textbf{Consigne :} Analyser la structure, la langue et l'intention de l'article.
\tcblower
\textbf{1. Structure (Aufbau) :}
L'article respecte la \textbf{pyramide inversée} :
\begin{itemize}
    \item \textbf{Titel :} \all{Neuer Markt in Yaoundé eröffnet} (Infinitif passé / nominalisé, informatif).
    \item \textbf{Ortzeile / Dateline :} \all{YAOUNDE.} (Lieu en majuscules, standard agency).
    \item \textbf{Lead (Vorspann / 1. Satz) :} \all{Der neue Zentralmarkt ... wurde gestern feierlich eröffnet.} Répond aux 5 W : \cle{Was?} (Eröffnung), \cle{Wo?} (Yaoundé), \cle{Wann?} (gestern), \cle{Wer?} (Zentralmarkt), \cle{Wie?} (feierlich).
    \item \textbf{Hauptteil (Corps) :} Détails par ordre d'importance : Description (Platz, Anlagen), Contexte/Problème passé (Straßenrand, Verkehr), Citation (Marie Nkou), Chiffres/Financement (2 Jahre, 3,5 Mrd, Weltbank/Staat).
    \item \textbf{Schluss (Ausblick) :} \all{Der Markt soll die Hygiene verbessern ...} (Modal \cle{sollen} + infinitif = perspective officielle/future).
\end{itemize}

\textbf{2. Langue et style (Sprache/Stil) :}
\begin{itemize}
    \item \textbf{Temps :} Dominance du \textbf{Präteritum} (\all{eröffnete, schnitt, dauerten, kosteten, erfolgte, sagte, verkauften, behinderte}) $\rightarrow$ Narration objective, distance journalistique.
    \item \textbf{Passif :} \all{wurde ... eröffnet, erfolgte} (focus sur l'événement, pas sur l'agent).
    \item \textbf{Style direct :} \all{„Endlich haben wir einen würdigen Arbeitsplatz“} (Perfekt \cle{haben} = oralité, subjectivité, émotion). Verbe introducteur \cle{sagte} (Präteritum).
    \item \textbf{Vocabulaire précis :} \cle{sanitäre Anlagen, würdig, behindern, ankurbeln, Finanzierung}.
    \item \textbf{Connecteurs :} \cle{Bisher} (contraste passé/présent), \cle{sowie} (ajout), \cle{und} (coordination).
\end{itemize}

\textbf{3. Intention (Absicht) :}
Informer le public sur une réalisation concrète (infrastructure), valoriser l'action publique (Mairie, État, Banque Mondiale), donner la parole aux bénéficiaires (commerçante), annoncer les retombées positives (hygiène, commerce).
\end{exoresolu}

\coursec{Production guidée : écrire un court article}

\courssub{Structure type et phrases-modèles (Bausteine)}
\begin{methode}[Rédiger un court article de journal : structure et style]
\textbf{Objectif :} Produire un texte cohérent de 80-100 mots respectant les codes journalistiques.\par
\begin{enumerate}
    \item \textbf{TITRE (Titel) :} Court, nominal, informatif, présent ou infinitif passé. \all{Neuer Brunnen für das Dorf} / \all{Einweihung der neuen Schule}.
    \item \textbf{LEAD (Vorspann / Erster Absatz) :} 1 phrase. Répond aux 5 W (\cle{W-Fragen}) : \cle{Wo?} \cle{Was?} \cle{Wann?} \cle{Wer?} \cle{Warum?}.
        \all{In Yaoundé wurde gestern ein neuer Brunnen eingeweiht.}
    \item \textbf{CORPS (Hauptteil) :} Développement par paragraphes (\cle{Absätze}). Ordre : important $\rightarrow$ détails. \textbf{Präteritum}. Connecteurs : \cle{zuerst, dann, schließlich, außerdem, jedoch}.
    \item \textbf{CITATION (Zitat) :} Apporte vie et subjectivité. Style direct (guillemets „ “) + verbe introducteur au Präteritum (\cle{sagen, erklären, betonen, freuen, meinen}).
        \all{„Das Wasser ist sauber“, \textbf{sagte} ein Anwohner.}
    \item \textbf{CONCLUSION (Ausblick / Schluss) :} Perspective future (Futur I \cle{werden} ou modal \cle{sollen/wollen/müssen}) ou bilan.
        \all{Der Brunnen \textbf{wird} die Gesundheit \textbf{verbessern}.}
    \item \textbf{Check-list finale :} Majuscules aux noms ? Verbe V2 (principale) / Verbe fin (subordonnée) ? \cle{ß} vs \cle{ss} ? Umlaute ? Articles/pluriels/cas corrects ?
\end{enumerate}
\cle{Modèles de phrases-types (Bausteine) pour l'écrit :}
\begin{itemize}
    \item \all{Wie die Behörde / der Sprecher mitteilte, ...} (Selon les autorités / le porte-parole... + sub. \cle{dass} ou infinitif).
    \item \all{Es wird erwartet, dass ...} (On s'attend à ce que... + \cle{dass}).
    \item \all{Laut einem Bericht der Zeitung ...} (Selon un rapport du journal...).
    \item \all{Wie Augenzeugen berichteten, ...} (Selon des témoins oculaires...).
\end{itemize}
\end{methode}

\courssub{Production d'un article pour le journal scolaire}
\begin{exoresolu}[Production guidée : « Ein Sportfest an unserer Schule »]
\textbf{Consigne :} Écrire un court article (env. 90 mots) pour le journal scolaire. Utiliser les éléments donnés et respecter la structure Titre / Lead / Corps / Citation / Conclusion. Conjuguer au Präteritum (sauf citation).
\textbf{Mots-clés :} \cle{gestern, Sportfest, Schulhof, viele Schüler, laufen, springen, Preise, Direktor, eröffnen, Spaß, nächstes Jahr, wieder}.
\tcblower
\textbf{Proposition de rédaction modèle (Corrigé type Bac) :}
\begin{textebox}[Article élève – Journal scolaire]
„Sportfest auf dem Schulhof

YAOUNDE. Gestern fand auf dem Schulhof unseres Lyzeums das jährliche Sportfest statt. Hunderte Schüler nahmen an den Wettkämpfen teil. Zuerst liefen die Jungen und Mädchen die 100 Meter. Dann sprangen sie weit und hoch. Der Direktor, Herr Fouda, eröffnete die Feier um 9 Uhr. Das Wetter war gut, obwohl es am Morgen geregnet hatte. „Es hat großen Spaß gemacht“, sagte die Schülerin Amina. Nächstes Jahr soll das Sportfest wieder stattfinden.“
\end{textebox}

\textbf{Analyse grammaticale et lexicale du modèle :}
\begin{itemize}
    \item \textbf{Titre :} Nominal (\cle{Sportfest auf dem Schulhof}).
    \item \textbf{Lead :} \all{fand ... statt} (Präteritum, verbe séparable \cle{stattfinden}). \cle{unseres Lyzeums} (Génitif chaîne nominale). \cle{Hunderte Schüler} (Accusatif pluriel).
    \item \textbf{Corps :} \cle{nahmen ... teil} (Präteritum \cle{teilnehmen}, séparable, + Dativ). \cle{liefen} (Präteritum \cle{laufen} fort). \cle{sprangen} (Präteritum \cle{springen} fort). \cle{eröffnete} (Präteritum \cle{eröffnen} faible).
    \item \textbf{Subordonnée temporelle (antériorité) :} \all{obwohl es am Morgen geregnet \textbf{hatte}} (Plus-que-parfait / \cle{Plusquamperfekt} : auxiliaire \cle{hatte} + Partizip II \cle{geregnet} – niveau A2+/B1, bienvenu pour l'antériorité).
    \item \textbf{Citation :} Style direct \all{„Es hat großen Spaß gemacht“} (Perfekt oral : \cle{machen} + \cle{Spaß}). Verbe introducteur \cle{sagte} (Präteritum).
    \item \textbf{Conclusion :} \all{soll ... stattfinden} (Modal \cle{sollen} + infinitif – projet futur officiel).
    \item \textbf{Vocabulaire clé :} \cle{der Schulhof, die Höfe} ; \cle{der Wettkampf, die -e} ; \cle{teilnehmen an + Dat} ; \cle{der Direktor, die -en} ; \cle{die Schülerin, die -nen} ; \cle{der Preis, die -e}.
\end{itemize}
\textbf{Note de l'examinateur :} Le texte respecte la structure, les temps (Präteritum narration, Perfekt citation, Plusquamperfekt antériorité, Modal futur), l'ordre des mots (V2, verbe fin), la majuscule aux noms. Niveau A2+ validé.


\end{exoresolu}

\courssub{Check-list finale et erreurs à éviter (Préparation Bac)}
\begin{attention}[Les 5 erreurs qui coûtent des points au Bac (Allemand LV2 – 1ère A)]
\begin{enumerate}
    \item \textbf{Majuscules oubliées sur les noms :} \all{*der markt, *die schule, *ein artikel.} $\rightarrow$ \textbf{Systématique :} \all{der Markt, die Schule, ein Artikel.} \textit{(Règle absolue : Tout nom = Majuscule).}
    \item \textbf{Verbe pas à la fin dans la subordonnée (\cle{weil, dass, obwohl, wenn, als, nachdem, bevor}) :} \all{*Weil ich bin müde. *Obwohl er kommt nicht.} $\rightarrow$ \textbf{Correct :} \all{Weil ich müde \textbf{bin}. Obwohl er nicht \textbf{kommt}.} \textit{(Réflexe : Conjonction $\rightarrow$ « parking » pour le verbe à la fin).}
    \item \textbf{Confusion \cle{werden} (futur/passif présent) / \cle{worden} (passif parfait) / \cle{wurden} (passif prétérit) :}
        \begin{itemize}
            \item \all{*Der Artikel ist veröffentlicht.} (Manque \cle{worden} pour passif Perfekt).
            \item \textbf{Correct :} \all{Der Artikel \textbf{ist veröffentlicht worden}.} (Passif Perfekt) vs \all{Der Artikel \textbf{wurde veröffentlicht}.} (Passif Präteritum) vs \all{Der Artikel \textbf{wird veröffentlicht}.} (Passif Présent/Futur).
        \end{itemize}
    \item \textbf{Faux amis et calques du français :}
        \begin{itemize}
            \item \all{*aktuell} = « actuel / en ce moment / récent » (PAS « actuel » au sens de « réel / effectif » $\rightarrow$ \cle{tatsächlich, wirklich}).
            \item \all{*kontrollieren} = « vérifier / contrôler / inspecter » (PAS « maîtriser / commander » $\rightarrow$ \cle{beherrschen, leiten}).
            \item \all{*das Event} (anglicisme) $\rightarrow$ Préférer \cle{das Ereignis, die Ereignisse}.
            \item \cle{berichten \textbf{über} + Akk} (faire un reportage sur) vs \cle{berichten \textbf{von}} (raconter/parler de qqch de manière narrative).
            \item \cle{sich freuen \textbf{über} + Akk} (se réjouir DE qqch présent/passé) vs \cle{sich freuen \textbf{auf} + Akk} (avoir hâte DE qqch futur).
        \end{itemize}
    \item \textbf{Cas (Kasus) après prépositions et verbes (Pièges classiques) :}
        \begin{itemize}
            \item \cle{wegen / während / trotz / statt / außerhalb / innerhalb} + \textbf{Génitif} (écrit formel) : \all{wegen des Regens} (pas *wegen dem Regen – oral seulement).
            \item \cle{helfen / danken / gefallen / folgen / glauben / begegnen / gratulieren} + \textbf{Datif} : \all{Ich helfe \textbf{dem} Mann.} (Pas *den Mann).
            \item \cle{warten auf / sich freuen über / berichten über / sprechen über / schreiben über} + \textbf{Akkusatif} : \all{Ich warte auf \textbf{den} Bus.}
        \end{itemize}
\end{enumerate}
\cle{Conseil final :} Relisez votre copie \textbf{à l'envers} (dernière phrase $\rightarrow$ première) pour ne voir que la forme (mots, majuscules, verbe fin, cas) et non le sens.
\end{attention}

\newpage

\coursec{Exercices}

\courssub{Je vérifie mes connaissances}

\begin{exercice}[QCM : la presse et l'article]
Entoure la bonne réponse. Une seule réponse est correcte.
\begin{enumerate}
\item \all{Die Überschrift} signifie :
\begin{itemize}
\item[a)] le lecteur
\item[b)] le titre (d'un article)
\item[c)] le journaliste
\end{itemize}
\item \all{Der Artikel} est un nom :
\begin{itemize}
\item[a)] masculin, pluriel \all{die Artikel}
\item[b)] féminin, pluriel \all{die Artikeln}
\item[c)] neutre, pluriel \all{die Artikel}
\end{itemize}
\item Dans un article de journal allemand, le temps le plus fréquent est :
\begin{itemize}
\item[a)] le Futur
\item[b)] le Präteritum
\item[c)] le Konjunktiv II
\end{itemize}
\item \all{berichten} signifie :
\begin{itemize}
\item[a)] lire
\item[b)] raconter / faire un reportage
\item[c)] annuler
\end{itemize}
\item Dans la subordonnée \all{Weil es regnete, wurde das Spiel abgesagt}, le verbe conjugué se trouve :
\begin{itemize}
\item[a)] en deuxième position
\item[b)] au milieu de la phrase
\item[c)] à la fin de la subordonnée
\end{itemize}
\end{enumerate}
\end{exercice}

\begin{exercice}[Vrai ou faux ? Justifie ta réponse]
Dis si les affirmations suivantes sont vraies ou fausses et justifie chaque réponse par une phrase en français.
\begin{enumerate}
\item \all{Der Leser} signifie « le journaliste ».
\item Au Perfekt, on utilise l'auxiliaire \all{haben} ou \all{sein} + le participe passé.
\item Dans une subordonnée introduite par \all{dass}, le verbe conjugué reste en deuxième position.
\item \all{Das Ereignis} est un nom féminin.
\item \all{als} s'emploie pour parler d'un événement unique au passé.
\end{enumerate}
\end{exercice}

\begin{exercice}[Familles de mots]
Complète le tableau suivant avec le mot manquant de chaque famille, en donnant l'article et le pluriel des noms.

\begin{center}
\begin{tabularx}{\linewidth}{|X|X|X|}
\hline
\rowcolor{popblueL} Verbe & Nom (article + pluriel) & Traduction du nom \\
\hline
lesen & der Leser, die Leser & le lecteur \\
\hline
berichten & \ldots & \ldots \\
\hline
informieren & \ldots & \ldots \\
\hline
schreiben & \ldots & \ldots \\
\hline
\ldots & die Überschrift, die Überschriften & le titre \\
\hline
\end{tabularx}
\end{center}
\end{exercice}

\begin{exercice}[Conjugaison : Perfekt ou Präteritum ?]
Mets le verbe entre parenthèses à la forme correcte. Choisis le Präteritum pour le récit journalistique (phrases 1 à 5) et le Perfekt pour le dialogue oral (phrases 6 à 10).
\begin{enumerate}
\item Gestern \ldots\ (finden) ein großes Konzert in Yaoundé statt.
\item Viele Jugendliche \ldots\ (kommen) aus allen Stadtteilen.
\item Der Bürgermeister \ldots\ (eröffnen) die Veranstaltung um 18 Uhr.
\item Es \ldots\ (geben) Musik, Tanz und ein Fußballspiel.
\item Die Zeitung \ldots\ (berichten) am nächsten Tag über das Fest.
\item — \ldots\ du gestern das Konzert \ldots\ (sehen) ?
\item — Ja, ich \ldots\ mit meinen Freunden \ldots\ (gehen).
\item — Was \ldots\ ihr \ldots\ (machen) ?
\item — Wir \ldots\ lange \ldots\ (tanzen) und viel \ldots\ (lachen).
\item — Mein Bruder \ldots\ leider zu Hause \ldots\ (bleiben).
\end{enumerate}
\end{exercice}

\courssub{Je m'entraîne}

\begin{exercice}[Thème : phrases à traduire]
Traduis les phrases suivantes en allemand. Attention à la place du verbe et au choix du temps.
\begin{enumerate}
\item Samedi, un tournoi a eu lieu dans notre lycée.
\item Beaucoup d'élèves sont venus parce qu'il faisait beau.
\item Le directeur a ouvert la fête à neuf heures.
\item Le journal de l'école a écrit un article sur cet événement.
\item Les lecteurs ont trouvé l'article très intéressant.
\end{enumerate}
\end{exercice}

\begin{exercice}[Des phrases simples aux subordonnées]
Relie les deux phrases avec la conjonction indiquée entre parenthèses. Fais attention à la place du verbe conjugué.
\begin{enumerate}
\item Es regnete. Das Spiel wurde abgesagt. (\all{weil})
\item Der Reporter schreibt einen Artikel. Das Fest war schön. (\all{dass})
\item Das Konzert begann. Es war acht Uhr. (\all{als})
\item Viele Menschen kamen. Der Eintritt war frei. (\all{weil})
\item Die Schüler freuen sich. Die Ferien beginnen morgen. (\all{dass})
\end{enumerate}
\end{exercice}

\begin{exercice}[Texte à trous : un début d'article]
Complète l'article suivant avec les mots de la liste : \all{berichtete — Ereignis — Leser — Überschrift — Artikel — statt}.

\begin{textebox}[Ein Fest in der Schule]
Unter der \ldots\ „Ein Tag voller Musik“ \ldots\ die Schülerzeitung gestern über ein großes \ldots\ : Das Schulfest fand am Samstag \ldots\ . Mehr als zweihundert Gäste kamen. Der \ldots\ war sehr lang, aber die \ldots\ fanden ihn spannend und informativ.
\end{textebox}
\end{exercice}

\begin{exercice}[Appariements : mots et définitions]
Relie chaque mot allemand (1 à 6) à sa définition française (a à f).
\begin{enumerate}
\item \all{der Titel}
\item \all{der Leser}
\item \all{das Ereignis}
\item \all{berichten}
\item \all{die Zeitung}
\item \all{die Überschrift}
\end{enumerate}
\begin{itemize}
\item[a)] la personne qui lit un article
\item[b)] raconter un fait dans un média
\item[c)] le nom d'un article ou d'un livre
\item[d)] un fait important qui se passe
\item[e)] le texte en gros caractères au-dessus de l'article
\item[f)] un média imprimé qui paraît chaque jour
\end{itemize}
\end{exercice}

\courssub{Je me prépare au Bac}

\begin{exercice}[Compréhension écrite : un événement culturel]
Lis l'article suivant, puis réponds aux questions.

\begin{textebox}[Kulturfestival in Douala begeistert die Jugend]
Am Samstag fand in Douala ein großes Kulturfestival statt. Mehr als tausend Besucher kamen in den Stadtpark, obwohl es am Morgen regnete. Das Festival begann um 15 Uhr mit einem Konzert junger Musiker aus Kamerun und Deutschland. Danach tanzten Gruppen aus fünf Regionen des Landes traditionelle Tänze. Die Organisatorin, Frau Mbarga, erklärte: „Wir wollen zeigen, dass Kultur die Jugend verbindet.“ Viele Schüler aus den Gymnasien der Stadt waren auch da. Sie schrieben kurze Berichte für ihre Schülerzeitungen. Ein deutscher Journalist aus Berlin berichtete ebenfalls über das Festival. Sein Artikel erschien am Montag in einer großen Zeitung. Am Abend gab es noch ein Fußballspiel zwischen zwei Schulmannschaften. Das Spiel endete 2:2, und alle feierten zusammen. Die Besucher fanden das Festival sehr schön, weil es Musik, Tanz und Sport verband. Nächstes Jahr wollen die Organisatoren das Festival wiederholen.
\end{textebox}

\textbf{A. Questions de compréhension} — Réponds en allemand par des phrases complètes.
\begin{enumerate}
\item Wann fand das Festival statt ?
\item Wer tanzte traditionelle Tänze ?
\item Was erklärte Frau Mbarga ?
\item Warum fanden die Besucher das Festival schön ?
\item Was wollen die Organisatoren nächstes Jahr machen ?
\end{enumerate}

\textbf{B. Richtig oder falsch ?} — Justifie chaque réponse par une citation du texte.
\begin{enumerate}
\item Das Festival begann am Morgen.
\item Ein deutscher Journalist schrieb einen Artikel über das Festival.
\item Das Fußballspiel endete 3:1.
\end{enumerate}

\textbf{C. Questions de langue}
\begin{enumerate}
\item Releve dans le texte trois verbes au Präteritum et donne leur infinitif.
\item Releve une subordonnée avec \all{weil} et recopie-la.
\item Trouve dans le texte le contraire de : \all{klein}, \all{hässlich}.
\end{enumerate}
\end{exercice}

\begin{exercice}[Thème et version]
\textbf{A. Thème} — Traduis en allemand les phrases suivantes (fais attention au temps et à la place du verbe).
\begin{enumerate}
\item Samedi, un tournoi de football a eu lieu à Yaoundé.
\item Beaucoup d'élèves sont venus parce qu'il faisait beau.
\item Le match a commencé quand le directeur est arrivé.
\item Un journaliste a écrit un article sur cet événement.
\item Les lecteurs ont dit que l'article était très bon.
\end{enumerate}

\textbf{B. Version} — Traduis en français l'article suivant.

\begin{textebox}[Neue Bibliothek eröffnet]
Am Freitag eröffnete die Schule eine neue Bibliothek. Der Direktor hielt eine kurze Rede und dankte allen Lehrern. Viele Schüler kamen, weil sie die neuen Bücher sehen wollten. Die Bibliothek hat jetzt zweitausend Bücher auf Deutsch, Französisch und Englisch. Eine Schülerin sagte: „Ich komme jetzt jeden Tag, weil ich gern lese.“ Die Schülerzeitung berichtete am Montag über die Eröffnung.
\end{textebox}
\end{exercice}

\begin{exercice}[Production écrite guidée : un article de journal]
Tu es journaliste pour le journal de ton lycée. Rédige en allemand un court article (8 à 10 phrases, environ 80 à 100 mots) sur la fête de la jeunesse du 11 février dans ta ville. Respecte le plan imposé suivant :
\begin{enumerate}
\item \textbf{Überschrift} : donne un titre à ton article.
\item \textbf{Einleitung} : quand et où la fête a eu lieu (Präteritum).
\item \textbf{Hauptteil} : décris deux ou trois activités (défilé, discours, match de football, danses) ; utilise au moins une subordonnée avec \all{weil} et une avec \all{als}.
\item \textbf{Zitat} : cite une phrase d'un élève ou d'un organisateur (guillemets allemands „ … “).
\item \textbf{Schluss} : donne ton opinion ou une perspective pour l'année prochaine.
\end{enumerate}
Mots utiles : \all{stattfinden, der Umzug, die Rede, die Jugend, feiern, berichten, das Ereignis}.
\end{exercice}

\newpage

\coursec{Corrigés des exercices}

\begin{corrige}[Exercice 1 — QCM : la presse et l'article]
\begin{enumerate}
\item \textbf{b)} \all{Die Überschrift} signifie « le titre (d'un article) ». C'est le texte en gros caractères au-dessus de l'article. « Le lecteur » se dit \all{der Leser} et « le journaliste » \all{der Journalist}.
\item \textbf{a)} \all{Der Artikel} est un nom \textbf{masculin}, pluriel \all{die Artikel} (sans terminaison ajoutée). Attention au faux ami : en français « un article » peut être un objet ; en allemand, \all{der Artikel} désigne surtout le texte de journal (et aussi le mot grammatical : \all{der, die, das}).
\item \textbf{b)} Le temps le plus fréquent dans un article de journal allemand est le \textbf{Präteritum}, temps du récit écrit. Le Perfekt domine à l'oral, le Konjunktiv II sert au conditionnel ou au discours rapporté.
\item \textbf{b)} \all{berichten} signifie « raconter / faire un reportage ». « Lire » se dit \all{lesen} et « annuler » \all{absagen}.
\item \textbf{c)} Dans la subordonnée introduite par \all{weil}, le verbe conjugué (\all{regnete}) se trouve \textbf{à la fin} : \all{Weil es regnete, wurde das Spiel abgesagt.} C'est la règle fondamentale de la subordonnée allemande.
\end{enumerate}
\textbf{Points de vigilance} : retenir le trio \all{der Titel / die Überschrift / der Artikel} ; ne pas confondre \all{der Leser} (le lecteur) et \all{der Lehrer} (le professeur).
\end{corrige}

\begin{corrige}[Exercice 2 — Vrai ou faux ? Justifie ta réponse]
\begin{enumerate}
\item \textbf{Faux.} \all{Der Leser} signifie « le lecteur », c'est-à-dire la personne qui lit ; « le journaliste » se dit \all{der Journalist}.
\item \textbf{Vrai.} Le Perfekt se forme avec l'auxiliaire \all{haben} ou \all{sein} conjugué au présent, suivi du participe passé placé à la fin : \all{Ich habe einen Artikel gelesen.} / \all{Er ist gekommen.}
\item \textbf{Faux.} Dans une subordonnée introduite par \all{dass}, le verbe conjugué va \textbf{à la fin} : \all{Die Leser sagten, dass der Artikel gut war.} La deuxième position est réservée à la phrase principale.
\item \textbf{Faux.} \all{Das Ereignis} est un nom \textbf{neutre} (pluriel : \all{die Ereignisse}). Il faut apprendre chaque nom avec son article.
\item \textbf{Vrai.} \all{als} introduit un événement unique et ponctuel au passé : \all{Als das Konzert begann, war es acht Uhr.} Pour une habitude au passé, on utilise \all{wenn}.
\end{enumerate}
\end{corrige}

\begin{corrige}[Exercice 3 — Familles de mots]
\begin{center}
\begin{tabularx}{\linewidth}{|X|X|X|}
\hline
\rowcolor{popblueL} Verbe & Nom (article + pluriel) & Traduction du nom \\
\hline
lesen & der Leser, die Leser & le lecteur \\
\hline
berichten & \textbf{der Bericht, die Berichte} & \textbf{le reportage / le compte rendu} \\
\hline
informieren & \textbf{die Information, die Informationen} & \textbf{l'information} \\
\hline
schreiben & \textbf{der Schriftsteller, die Schriftsteller} (ou \all{die Schrift, die Schriften}) & \textbf{l'écrivain} (ou « l'écriture ») \\
\hline
\textbf{überschreiben} & die Überschrift, die Überschriften & le titre \\
\hline
\end{tabularx}
\end{center}
\textbf{Remarques} : \all{überschreiben} est un verbe \textbf{inséparable} (participe : \all{überschrieben}) car le préfixe \all{über-} y est inaccentué. Pour \all{schreiben}, on accepte aussi \all{der Schreiber, die Schreiber} (le scribe), mais \all{der Schriftsteller} est la réponse attendue au lycée.
\end{corrige}

\begin{corrige}[Exercice 4 — Conjugaison : Perfekt ou Präteritum ?]
\textbf{Récit journalistique (Präteritum) :}
\begin{enumerate}
\item Gestern \textbf{fand} ein großes Konzert in Yaoundé \textbf{statt}. (\all{stattfinden} : verbe fort à particule séparable ; Präteritum \all{fand ... statt})
\item Viele Jugendliche \textbf{kamen} aus allen Stadtteilen. (\all{kommen — kam — ist gekommen})
\item Der Bürgermeister \textbf{eröffnete} die Veranstaltung um 18 Uhr. (verbe faible régulier)
\item Es \textbf{gab} Musik, Tanz und ein Fußballspiel. (\all{geben — gab — hat gegeben})
\item Die Zeitung \textbf{berichtete} am nächsten Tag über das Fest. (verbe faible ; \all{berichten über + Akkusativ})
\end{enumerate}
\textbf{Dialogue oral (Perfekt) :}
\begin{enumerate}
\setcounter{enumi}{5}
\item — \textbf{Hast} du gestern das Konzert \textbf{gesehen} ? (\all{sehen — sah — hat gesehen})
\item — Ja, ich \textbf{bin} mit meinen Freunden \textbf{gegangen}. (\all{gehen} : verbe de déplacement, auxiliaire \all{sein})
\item — Was \textbf{habt} ihr \textbf{gemacht} ?
\item — Wir \textbf{haben} lange \textbf{getanzt} und viel \textbf{gelacht}. (verbes faibles, auxiliaire \all{haben})
\item — Mein Bruder \textbf{ist} leider zu Hause \textbf{geblieben}. (\all{bleiben} : toujours avec \all{sein} !)
\end{enumerate}
\textbf{Points de vigilance} : le participe passé se place \textbf{à la fin} de la phrase ; les verbes de déplacement (\all{gehen, kommen}) et \all{bleiben} prennent l'auxiliaire \all{sein}.
\end{corrige}

\begin{corrige}[Exercice 5 — Thème : phrases à traduire]
\begin{enumerate}
\item \all{Am Samstag fand ein Turnier in unserem Gymnasium statt.} (Präteritum de \all{stattfinden} ; \all{am Samstag} en tête, le verbe reste en deuxième position.)
\item \all{Viele Schüler kamen, weil das Wetter schön war.} (ou \all{... weil es schönes Wetter gab.} — dans la subordonnée avec \all{weil}, le verbe \all{war} va à la fin.)
\item \all{Der Direktor eröffnete das Fest um neun Uhr.}
\item \all{Die Schülerzeitung schrieb einen Artikel über dieses Ereignis.} (\all{schreiben — schrieb} ; \all{über + Akkusativ} : \all{dieses Ereignis}, neutre.)
\item \all{Die Leser fanden den Artikel sehr interessant.} (\all{finden — fanden} ; \all{den Artikel} : accusatif masculin.)
\end{enumerate}
\textbf{Points de vigilance} : ne jamais écrire « \all{Viele Schüler kamen, weil das Wetter war schön} » : après \all{weil}, le verbe conjugué va obligatoirement à la fin.
\end{corrige}

\begin{corrige}[Exercice 6 — Des phrases simples aux subordonnées]
\begin{enumerate}
\item \all{Das Spiel wurde abgesagt, weil es regnete.} (ou \all{Weil es regnete, wurde das Spiel abgesagt.})
\item \all{Der Reporter schreibt in seinem Artikel, dass das Fest schön war.} (subordonnée avec \all{dass} : verbe \all{war} à la fin ; on ne peut pas écrire « \all{einen Artikel, dass ...} », il faut \all{in seinem Artikel}.)
\item \all{Als das Konzert begann, war es acht Uhr.} (\all{als} : événement unique au passé ; si la subordonnée ouvre la phrase, le verbe de la principale \all{war} vient juste après la virgule.)
\item \all{Viele Menschen kamen, weil der Eintritt frei war.}
\item \all{Die Schüler freuen sich, dass die Ferien morgen beginnen.}
\end{enumerate}
\textbf{Règle rappelée} : dans la subordonnée, le verbe conjugué occupe la \textbf{dernière place} ; quand la subordonnée est en tête, elle compte comme la première position et la principale commence par son verbe : \all{Als ..., war ...}.
\end{corrige}

\begin{corrige}[Exercice 7 — Texte à trous : un début d'article]
\begin{textebox}[Ein Fest in der Schule — corrigé]
Unter der \textbf{Überschrift} „Ein Tag voller Musik“ \textbf{berichtete} die Schülerzeitung gestern über ein großes \textbf{Ereignis} : Das Schulfest fand am Samstag \textbf{statt}. Mehr als zweihundert Gäste kamen. Der \textbf{Artikel} war sehr lang, aber die \textbf{Leser} fanden ihn spannend und informativ.
\end{textebox}
\textbf{Justifications} : \all{unter der Überschrift} (datif féminin après \all{unter}) ; \all{berichtete} : Präteritum de \all{berichten}, temps du récit de presse ; \all{stattfinden} est séparable : \all{fand ... statt} ; \all{die Leser} : pluriel identique au singulier, seul l'article change.
\end{corrige}

\begin{corrige}[Exercice 8 — Appariements : mots et définitions]
\begin{center}
\begin{tabularx}{\linewidth}{|X|X|}
\hline
\rowcolor{popblueL} Mot allemand & Définition \\
\hline
1. der Titel & c) le nom d'un article ou d'un livre \\
\hline
2. der Leser & a) la personne qui lit un article \\
\hline
3. das Ereignis & d) un fait important qui se passe \\
\hline
4. berichten & b) raconter un fait dans un média \\
\hline
5. die Zeitung & f) un média imprimé qui paraît chaque jour \\
\hline
6. die Überschrift & e) le texte en gros caractères au-dessus de l'article \\
\hline
\end{tabularx}
\end{center}
\textbf{Astuce mémo} : \all{die Überschrift} est littéralement « l'écriture au-dessus » (\all{über} + \all{die Schrift}) ; \all{das Ereignis} vient de \all{sich ereignen} (se produire).
\end{corrige}

\begin{corrige}[Exercice 9 — Compréhension écrite : un événement culturel]
\textbf{Barème indicatif (20 points)} : A. Questions de compréhension : 5 × 2 = 10 pts (1 pt pour l'information exacte, 1 pt pour la correction de la phrase) ; B. Richtig oder falsch : 3 × 2 = 6 pts (1 pt pour la réponse, 1 pt pour la citation) ; C. Questions de langue : 4 pts.

\textbf{A. Questions de compréhension — réponses modèles :}
\begin{enumerate}
\item \all{Das Festival fand am Samstag in Douala statt.}
\item \all{Gruppen aus fünf Regionen des Landes tanzten traditionelle Tänze.}
\item \all{Frau Mbarga erklärte, dass Kultur die Jugend verbindet.} (ou citation directe : \all{Sie erklärte: „Wir wollen zeigen, dass Kultur die Jugend verbindet.“})
\item \all{Die Besucher fanden das Festival schön, weil es Musik, Tanz und Sport verband.}
\item \all{Nächstes Jahr wollen die Organisatoren das Festival wiederholen.}
\end{enumerate}

\textbf{B. Richtig oder falsch ?}
\begin{enumerate}
\item \textbf{Falsch.} \all{Das Festival begann um 15 Uhr mit einem Konzert.} (Il a commencé l'après-midi, pas le matin ; le matin, il pleuvait.)
\item \textbf{Richtig.} \all{Ein deutscher Journalist aus Berlin berichtete ebenfalls über das Festival. Sein Artikel erschien am Montag in einer großen Zeitung.}
\item \textbf{Falsch.} \all{Das Spiel endete 2:2.} (et non 3:1.)
\end{enumerate}

\textbf{C. Questions de langue :}
\begin{enumerate}
\item Trois verbes au Präteritum (au choix) : \all{fand} (\all{finden}), \all{kamen} (\all{kommen}), \all{begann} (\all{beginnen}), \all{tanzten} (\all{tanzen}), \all{erklärte} (\all{erklären}), \all{schrieben} (\all{schreiben}), \all{berichtete} (\all{berichten}), \all{erschien} (\all{erscheinen}), \all{gab} (\all{geben}), \all{endete} (\all{enden}), \all{verband} (\all{verbinden}).
\item Subordonnée avec \all{weil} : \all{Die Besucher fanden das Festival sehr schön, weil es Musik, Tanz und Sport verband.} (verbe \all{verband} à la fin.)
\item Contraires : \all{klein} → \all{groß} (\all{ein großes Kulturfestival}) ; \all{hässlich} → \all{schön} (\all{sehr schön}).
\end{enumerate}
\textbf{Points de vigilance} : répondre par des phrases complètes avec le verbe en deuxième position ; pour justifier, recopier la citation \textbf{mot pour mot}, sans la modifier.
\end{corrige}

\begin{corrige}[Exercice 10 — Thème et version]
\textbf{Barème indicatif (20 points)} : A. Thème : 5 × 2 = 10 pts (1 pt pour le temps et la conjugaison, 1 pt pour la place du verbe et le vocabulaire) ; B. Version : 10 pts (fidélité du sens : 6 pts ; qualité du français : 4 pts).

\textbf{A. Thème — corrigé :}
\begin{enumerate}
\item \all{Am Samstag fand ein Fußballturnier in Yaoundé statt.} (verbe séparable : \all{fand ... statt}.)
\item \all{Viele Schüler kamen, weil das Wetter schön war.} (subordonnée : \all{war} à la fin.)
\item \all{Das Spiel begann, als der Direktor ankam.} (\all{als} + événement unique au passé ; \all{ankam} : Präteritum du verbe séparable \all{ankommen}, à la fin de la subordonnée.)
\item \all{Ein Journalist schrieb einen Artikel über dieses Ereignis.}
\item \all{Die Leser sagten, dass der Artikel sehr gut war.} (subordonnée avec \all{dass} : \all{war} à la fin.)
\end{enumerate}

\textbf{B. Version — traduction modèle :}

« Vendredi, l'école a inauguré une nouvelle bibliothèque. Le directeur a prononcé un court discours et a remercié tous les professeurs. Beaucoup d'élèves sont venus parce qu'ils voulaient voir les nouveaux livres. La bibliothèque possède maintenant deux mille livres en allemand, en français et en anglais. Une élève a déclaré : „Je viens maintenant tous les jours, parce que j'aime lire.“ Le journal des élèves a rendu compte de l'inauguration lundi. »

\textbf{Points de vigilance} : \all{eröffnen} = « inaugurer / ouvrir officiellement » ; \all{eine Rede halten} = « prononcer un discours » (et non « tenir une redite » !) ; \all{die Eröffnung} = « l'inauguration » ; au thème, bien choisir le Präteritum pour ce récit écrit.
\end{corrige}

\begin{corrige}[Exercice 11 — Production écrite guidée : un article de journal]
\textbf{Barème indicatif (20 points)} : respect du plan en 5 parties (Überschrift, Einleitung, Hauptteil, Zitat, Schluss) : 5 pts ; emploi correct du Präteritum : 4 pts ; subordonnées avec \all{weil} et \all{als} (verbe à la fin) : 4 pts ; vocabulaire du thème : 4 pts ; orthographe (majuscules des noms, Umlaute) : 3 pts.

\textbf{Proposition de rédaction modèle :}

\begin{textebox}[Jugendfest in Yaoundé: Ein Tag voller Freude]
Am 11. Februar fand in Yaoundé das große Jugendfest statt. Tausende Jugendliche kamen am Morgen zum Stadion, weil das Wetter schön war. Der Bürgermeister hielt eine Rede und dankte allen Schülern. Als der Umzug begann, spielte eine Musikgruppe traditionelle Lieder. Viele Schüler tanzten und sangen zusammen. Am Nachmittag gab es ein Fußballspiel zwischen zwei Gymnasien der Stadt. Eine Schülerin sagte: „Dieses Fest ist das schönste Ereignis des Jahres, weil wir hier alle zusammen sind.“ Am Abend endete das Fest mit einem großen Konzert. Ich fand diesen Tag wunderbar, und ich hoffe, dass wir nächstes Jahr wieder so feiern.
\end{textebox}

\textbf{Traduction du modèle} : « Fête de la jeunesse à Yaoundé : une journée pleine de joie. Le 11 février, la grande fête de la jeunesse a eu lieu à Yaoundé. Des milliers de jeunes sont venus au stade le matin parce qu'il faisait beau. Le maire a prononcé un discours et a remercié tous les élèves. Quand le défilé a commencé, un groupe de musique a joué des chants traditionnels. Beaucoup d'élèves ont dansé et chanté ensemble. L'après-midi, il y a eu un match de football entre deux lycées de la ville. Une élève a dit : „Cette fête est le plus bel événement de l'année, parce que nous sommes tous ensemble ici.“ Le soir, la fête s'est terminée par un grand concert. J'ai trouvé cette journée merveilleuse, et j'espère que nous fêterons encore ainsi l'année prochaine. »

\textbf{Points de vigilance :}
\begin{itemize}
\item Date : \all{am 11. Februar} (avec un point après le chiffre : c'est un ordinal).
\item Subordonnées vérifiées : \all{..., weil das Wetter schön war.} / \all{Als der Umzug begann, spielte ...} (verbe conjugué à la fin de la subordonnée ; si elle est en tête, la principale commence par le verbe).
\item Citation avec les guillemets allemands „ … “ et deux-points avant la parole.
\item Majuscule à tous les noms : \all{das Fest, der Umzug, die Rede, das Ereignis}.
\item Verbes forts au Präteritum : \all{fand ... statt, kam, hielt, begann, gab, sagte, endete}.
\end{itemize}
\end{corrige}

\newpage

\coursec{Fiche bilan}

\begin{popfigure}
\centering
\begin{center}\tcbox[colback=popblue,colframe=popblue,coltext=white,arc=9pt,boxsep=4pt,left=6pt,right=6pt]{\bfseries\small L'article de journal}\end{center}
\vspace{-2pt}
\begin{tcbraster}[raster columns=3,raster equal height=rows,raster column skip=6pt,raster row skip=6pt,enhanced,blankest]
\begin{tcolorbox}[enhanced,colback=poporange,colframe=poporange,coltext=white,boxrule=0.9pt,arc=3pt,left=4pt,right=4pt,top=3pt,bottom=3pt,boxsep=1pt,halign=left]\bfseries\scriptsize Texte support \& Compréhension\end{tcolorbox}
\begin{tcolorbox}[enhanced,colback=popgreen,colframe=popgreen,coltext=white,boxrule=0.9pt,arc=3pt,left=4pt,right=4pt,top=3pt,bottom=3pt,boxsep=1pt,halign=left]\bfseries\scriptsize Lexique thématique (die Presse)\end{tcolorbox}
\begin{tcolorbox}[enhanced,colback=poppurple,colframe=poppurple,coltext=white,boxrule=0.9pt,arc=3pt,left=4pt,right=4pt,top=3pt,bottom=3pt,boxsep=1pt,halign=left]\bfseries\scriptsize Grammaire 1 Perfekt / Präteritum\end{tcolorbox}
\begin{tcolorbox}[enhanced,colback=poppink,colframe=poppink,coltext=white,boxrule=0.9pt,arc=3pt,left=4pt,right=4pt,top=3pt,bottom=3pt,boxsep=1pt,halign=left]\bfseries\scriptsize Grammaire 2 Subordonnées (dass, weil, wenn)\end{tcolorbox}
\begin{tcolorbox}[enhanced,colback=popgold,colframe=popgold,coltext=white,boxrule=0.9pt,arc=3pt,left=4pt,right=4pt,top=3pt,bottom=3pt,boxsep=1pt,halign=left]\bfseries\scriptsize Méthode Commentaire structuré\end{tcolorbox}
\begin{tcolorbox}[enhanced,colback=popdark,colframe=popdark,coltext=white,boxrule=0.9pt,arc=3pt,left=4pt,right=4pt,top=3pt,bottom=3pt,boxsep=1pt,halign=left]\bfseries\scriptsize Production Écrire un court article\end{tcolorbox}
\end{tcbraster}

\legende{Carte mentale de la leçon ALA18 : Média et communication}
\end{popfigure}

\begin{bilan}

\end{bilan}

\courssub{Lexique clé : die Presse und der Zeitungsartikel}

\begin{center}
\begin{tabularx}{\linewidth}{|>{\bfseries}l|X|X|}
\hline
\rowcolor{popblueL} \textbf{Deutsch (Artikel, Plural)} & \textbf{Français} & \textbf{Contexte / Exemple} \\
\hline
der Artikel, die Artikel & l'article & \all{Ich lese einen interessanten Artikel.} « Je lis un article intéressant. » \\
\hline
die Zeitung, die Zeitungen & le journal (quotidien) & \all{Die Zeitung erscheint täglich.} « Le journal paraît quotidiennement. » \\
\hline
die Überschrift, die Überschriften & le titre, la manchette & \all{Die Überschrift ist groß gedruckt.} « Le titre est en gros caractères. » \\
\hline
der Titel, die Titel & le titre (général) & \all{Was ist der Titel des Artikels?} « Quel est le titre de l'article ? » \\
\hline
der Leser, die Leser / die Leserin, die Leserinnen & le lecteur / la lectrice & \all{Der Leser will informiert werden.} « Le lecteur veut être informé. » \\
\hline
der Autor, die Autoren / die Autorin, die Autorinnen & l'auteur / l'autrice & \all{Der Autor berichtet objektiv.} « L'auteur rapporte objectivement. » \\
\hline
berichten (berichtete, hat berichtet) & rapporter, informer & \all{Er berichtet über das Ereignis.} « Il rapporte sur l'événement. » \\
\hline
das Ereignis, die Ereignisse & l'événement & \all{Ein wichtiges Ereignis gestern.} « Un événement important hier. » \\
\hline
die Meldung, die Meldungen & la dépêche, la brève & \all{Eine kurze Meldung im Radio.} « Une brève à la radio. » \\
\hline
der Bericht, die Berichte & le reportage, le compte-rendu & \all{Ein Bericht aus Yaoundé.} « Un reportage depuis Yaoundé. » \\
\hline
die Schlagzeile, die Schlagzeilen & la une, le gros titre & \all{Die Schlagzeile schockiert.} « La une choque. » \\
\hline
objektiv / subjektiv & objectif / subjectif & \all{Ein guter Artikel ist objektiv.} « Un bon article est objectif. » \\
\hline
die Quelle, die Quellen & la source & \all{Nennen Sie Ihre Quellen!} « Citez vos sources ! » \\
\hline
veröffentlichen & publier & \all{Der Artikel wurde gestern veröffentlicht.} « L'article a été publié hier. » \\
\hline
\end{tabularx}
\end{center}

\courssub{Grammaire 1 : Perfekt vs. Präteritum (Récit au passé)}

\begin{center}
\begin{tabularx}{\linewidth}{|>{\bfseries}p{2.5cm}|X|X|}
\hline
\rowcolor{poppurpleL} \textbf{Aspekt} & \textbf{Perfekt (Gesprächspassiv)} & \textbf{Präteritum (Schriftsprache)} \\
\hline
\cle{Emploi} & Langage oral, lettres, emails, dialogues. Action terminée au passé. & \cle{Écrit} : romans, journaux, récits formels, articles de presse. \\
\hline
\cle{Formation} & \all{haben/sein + Partizip II} (à la fin) & \cle{Verbes faibles} : radical + \cle{-te} + terminaison.\\
& \all{Ich habe gelesen.} & \cle{Verbes forts} : voyelle modifiée + terminaison. \\
& \all{Ich bin gegangen.} & \all{Ich las.} / \all{Ich ging.} \\
\hline
\cle{Auxiliaire} & \all{haben} (transitifs, la plupart) & Pas d'auxiliaire. Forme simple. \\
& \all{sein} (mouvement, changement d'état, \all{sein, bleiben, werden}) & \\
\hline
\cle{Marqueurs temporels} & \all{gestern, vor zwei Tagen, letzte Woche} (aussi en Präteritum) & \all{damals, damals, als, nachdem} (contexte narratif) \\
\hline
\cle{Exemple (Presse)} & \all{„Der Autor hat einen Artikel geschrieben.“} (Interview, oral) & \all{„Der Autor schrieb einen Artikel.“} (Texte journalistique standard) \\
\hline
\end{tabularx}
\end{center}

\courssub{Grammaire 2 : Les subordonnées (Nebensätze) — Structure verbe-final}

\begin{center}
\begin{tabularx}{\linewidth}{|>{\bfseries}p{2.2cm}|X|X|}
\hline
\rowcolor{popgreenL} \textbf{Conjonction} & \textbf{Sens / Fonction} & \textbf{Structure \& Exemple} \\
\hline
\cle{dass} & Complément d'objet (dire, penser, savoir, lire) & Verbe à la \cle{fin}. \\
& & \all{Der Leser weiß, \textbf{dass} der Artikel \textbf{erscheint}.} \\
\hline
\cle{weil} & Cause (réponse à \all{warum?}) & Verbe à la \cle{fin}. \\
& & \all{Der Artikel ist wichtig, \textbf{weil} er \textbf{informiert}.} \\
\hline
\cle{wen / wenn} & Temporel (présent/futur : \cle{quand} ; passé : \cle{lorsque}) & Verbe à la \cle{fin}. \\
& Condition (\cle{si} — avec Konjunktiv II plus tard) & \all{\textbf{Wenn} ich die Zeitung \textbf{lese}, lerne ich viel.} \\
\hline
\cle{ob} & Question indirecte fermée (si / whether) & Verbe à la \cle{fin}. \\
& & \all{Ich frage mich, \textbf{ob} der Titel \textbf{stimmt}.} \\
\hline
\cle{Word Order} & \textbf{Conjonction + Sujet + Compléments + VERBE (fin)} & \all{..., \textbf{weil} \textbf{ich} \textbf{Deutsch} \textbf{lerne}.} \\
\hline
\end{tabularx}
\end{center}

\courssub{Méthodes Express}

\begin{methode}[Commenter un article de journal (3 étapes)]
\begin{enumerate}
\item \textbf{Analyser la forme :} Identifier \cle{Überschrift}, \cle{Schlagzeile}, \cle{Autor}, \cle{Quelle}, structure (lead, corps, conclusion), ton (objektiv/subjektiv), temps (souvent \cle{Präteritum}).
\item \textbf{Comprendre le fond :} Quel \cle{Ereignis} ? Qui ? Quoi ? Où ? Quand ? Pourquoi ? (\cle{W-Fragen}). Distinguer faits (\cle{Tatsachen}) et opinions (\cle{Meinungen}).
\item \textbf{Structurer la réponse :} Introduction (source, sujet) $\rightarrow$ Résumé structuré (connecteurs : \all{zuerst, dann, schließlich}) $\rightarrow$ Appréciation personnelle (vocabulaire : \all{meiner Meinung nach, ich finde, dass...}).
\end{enumerate}
\end{methode}

\begin{methode}[Écrire un court article (Schreibplan)]
\begin{enumerate}
\item \textbf{Trouver un sujet :} Événement scolaire (tournoi, fête), actualité locale (Yaoundé, marché), fait divers.
\item \textbf{Rédiger la \cle{Überschrift} :} Courte, accrocheuse, nominale (\all{Fußballturnier in Yaoundé}).
\item \textbf{Le \cle{Lead} (premier paragraphe) :} Répondre aux 5 W (\all{Wer? Was? Wo? Wann? Warum?}) au \cle{Präteritum}.
\item \textbf{Le corps :} Développement chronologique, citations (\cle{dass}-Sätze : \all{Der Trainer sagte, dass...}), détails.
\item \textbf{Conclusion :} Conséquences, suite, avis. Relecture : majuscules noms, verbe position 2 (principal) / fin (subordonnée), \cle{Perfekt/Präteritum} cohérent.
\end{enumerate}
\end{methode}



\begin{aretenir}[Checklist avant le Bac — Leçon ALA18]
\begin{itemize}
\item $\square$ Je connais le vocabulaire de la presse (\cle{der Artikel, die Überschrift, der Leser, berichten, das Ereignis, veröffentlichen}) avec genre et pluriel.
\item $\square$ Je sais former le \cle{Partizip II} des verbes faibles (\all{gelesen}), forts (\all{geschrieben, gegangen}) et mixtes (\all{gebracht}).
\item $\square$ Je choisis l'auxiliaire correct (\cle{haben} vs \cle{sein}) pour le Perfekt.
\item $\square$ Je conjugue les verbes forts/irréguliers au \cle{Präteritum} (\all{ich las, ich ging, ich war, ich hatte}).
\item $\square$ Je sais quand utiliser \cle{Perfekt} (oral) et \cle{Präteritum} (écrit, presse, récit).
\item $\square$ Je place le verbe à la \cle{fin} dans toute subordonnée introduite par \cle{dass, weil, wenn, ob}.
\item $\square$ Je construis correctement l'ordre des mots : \all{Conjonction + Sujet + Compléments + Verbe}.
\item $\square$ Je sais répondre aux \cle{W-Fragen} (Wer, Was, Wo, Wann, Warum, Wie) sur un texte.
\item $\square$ Je sais rédiger un \cle{Lead} journalistique (5 W) au Präteritum.
\item $\square$ Je distingue \cle{objektiv} (faits, Präteritum) et \cle{subjektiv} (opinions, Konjunktiv II / modalverben).
\item $\square$ Je relis mes phrases : majuscules aux noms, Umlauts (ä/ö/ü), ß après voyelle longue, guillemets „ “.
\end{itemize}
\end{aretenir}

\findecours

\end{document}
